% Dissertation : produire une dissertation — Français Tle Tech — Cours signé OnBuch+
% Programme officiel MINESEC. Compiler avec : tectonic N02-dissertation-produire-dissertation.tex
\newcommand{\DOCMATIERE}{Français}
\newcommand{\DOCNIVEAU}{Tle Tech}
\newcommand{\DOCMODULE}{Français – classe de Terminale technique}
\newcommand{\DOCLECON}{N02}
\newcommand{\DOCTITRE}{Dissertation : produire une dissertation}
% OnBuch+ — préambule « cours signés » ENRICHI (compilé avec tectonic / XeLaTeX)
% Système visuel « Pop » d'OnBuch+ : fond crème (jamais blanc), bordures encre
% épaisses, ombres « dures » décalées, titres Archivo Black en capitales.
% Version MATHÉMATIQUES : graphes (pgfplots/tikz), boîtes pédagogiques
% (définition, propriété, méthode, attention, le savais-tu, exercice résolu,
% exercice, TP, bilan), page de garde et signature OnBuch+ ancrée en bas.
%
% Macros attendues dans main.tex AVANT \input{preamble} :
%   \DOCMATIERE  \DOCNIVEAU  \DOCMODULE  \DOCLECON (numéro)  \DOCTITRE
\documentclass[11pt]{article}

\usepackage{fontspec}
\usepackage[french]{babel}
\usepackage[a4paper,margin=2cm,top=3.4cm,bottom=2.8cm,headheight=1.4cm,footskip=1.3cm]{geometry}
\usepackage{amsmath,amssymb,amsfonts}
\usepackage{mathtools} % \xmapsto, \coloneqq... souvent utilisés par les modèles
\usepackage{cancel} % simplifications barrées (bilans de forces)
% \abs{z} (module, valeur absolue) et \norm{u} (norme), utilisés par les modèles
\providecommand{\abs}{}\let\abs\relax\DeclarePairedDelimiter{\abs}{\lvert}{\rvert}
\providecommand{\norm}{}\let\norm\relax\DeclarePairedDelimiter{\norm}{\lVert}{\rVert}
\usepackage[table]{xcolor} % [table] : fournit aussi \rowcolors (alternance de lignes)
\usepackage{pagecolor}
\usepackage{tikz}
\usetikzlibrary{arrows.meta,positioning,calc,shapes.geometric,shapes.misc,shapes.arrows,decorations.pathmorphing,decorations.pathreplacing,decorations.markings,patterns,shadows,mindmap,trees,fit,backgrounds,matrix,intersections,angles,quotes}
\usepackage{pgfplots}
\usepackage[version=4]{mhchem} % équations nucléaires \ce{^{235}_{92}U}
\usepackage[european,siunitx]{circuitikz} % schémas électriques
\pgfplotsset{compat=1.18}
\usepgfplotslibrary{fillbetween,groupplots} % aires entre courbes (calcul intégral)
\usepackage{enumitem}
\usepackage{fancyhdr}
% [most] charge toutes les bibliothèques tcolorbox (listings, minted, raster,
% documentation...) et gonfle inutilement la mémoire TeX sur un document long
% (~300 boîtes) : on ne charge que ce qui est réellement utilisé.
\usepackage[skins,breakable]{tcolorbox}
\usepackage{graphicx}
\usepackage{colortbl}
\usepackage{booktabs}
\usepackage{tabularx}
\usepackage{multirow}
\usepackage{diagbox} % case d'en-tête diagonale des tableaux à double entrée
% Colonnes extensibles centrée / gauche / droite (C, L, R), souvent utilisées par les modèles
\newcolumntype{C}{>{\centering\arraybackslash}X}
\newcolumntype{L}{>{\raggedright\arraybackslash}X}
\newcolumntype{R}{>{\raggedleft\arraybackslash}X}
\usepackage{array}
\usepackage{multicol}
\usepackage{siunitx}
% parse-numbers=false : accepte aussi \SI{2,5 \times 10^{-3}}{...}
\sisetup{locale=FR,output-decimal-marker={,},group-minimum-digits=4,parse-numbers=false}
\DeclareSIUnit\ohm{\ensuremath{\symup{\Omega}}} % U+2126 absent de la police
\DeclareSIUnit\division{div} % réglages d'oscilloscope (ms/div, V/div)
\DeclareSIUnit\tr{tr} % tours (vitesse de rotation, ex. \SI{1500}{\tr\per\minute})
\DeclareSIUnit\molar{M} % concentration molaire (ex. \SI{3}{\molar}, \SI{1}{\micro\molar})
\DeclareSIUnit\an{an} % année (le modèle confond parfois avec \year, un compteur TeX) — ex. \SI{5}{\kilo\meter\per\an}
\DeclareSIUnit\FCFA{FCFA} % franc CFA (ex. \SI{500}{\FCFA\per\kilo\gram})
\DeclareSIUnit\degreeBrix{\textdegree Brix} % teneur en sucre d'un sirop/jus (réfractométrie)
\DeclareSIUnit\month{mois} % le modèle utilise parfois \month, un compteur TeX, comme unité de durée
\DeclareSIUnit\microliter{\micro\liter} % le modèle écrit parfois \microliter en un seul token
\DeclareSIUnit\million{millions} % dénombrement (ex. \SI{15}{\million\per\mL} spermatozoïdes)
\usepackage{qrcode}
% \degree : commande fréquemment utilisée par les modèles pour un angle ou une
% température en dehors de \SI{}{} (non fournie par les paquets chargés).
\providecommand{\degree}{^{\circ}}
% \qed (fin de démonstration), utilisé par les modèles sans amsthm
\providecommand{\qed}{\hfill$\square$}
% \barre : double barre des valeurs interdites dans un tableau de variations (tkz-tab)
\providecommand{\barre}{\ensuremath{\|}}
% environnement proof (amsthm non chargé) utilisé par les modèles
\ifdefined\proof\else\newenvironment{proof}[1][Démonstration]{\par\noindent\textit{#1.}\ }{\qed\par}\fi
\usepackage{lastpage}
\usepackage{hyperref}
\hypersetup{hidelinks}

% ============================================================
% Palette « Pop » OnBuch+ — mêmes hex que la classe P (plus_ui.dart)
% ============================================================
\definecolor{poporange}{HTML}{F59321}
\definecolor{popdark}{HTML}{E07A0C}
\definecolor{popcream}{HTML}{FFF6E8}
\definecolor{popink}{HTML}{1C1714}
\definecolor{poppurple}{HTML}{7A5AE0}
\definecolor{popgreen}{HTML}{2E9E6A}
\definecolor{poppink}{HTML}{E0457B}
\definecolor{popblue}{HTML}{2F7DE1}
\definecolor{popgold}{HTML}{C98A12}
\definecolor{popmuted}{HTML}{8A7F76}
\definecolor{popline}{HTML}{EADFD2}
\definecolor{popgray}{HTML}{8A7F76}
\colorlet{popgrey}{popgray}
% Alias tolérants (le modèle nomme parfois les couleurs différemment)
\colorlet{popwhite}{white}
\colorlet{popblack}{popink}
\colorlet{popred}{poppink}
\colorlet{popyellow}{popgold}
% Teintes claires (fonds de tableaux/encadrés) — le modèle nomme parfois
% les couleurs avec un suffixe L (ex. popblueL) sans les avoir définies.
\colorlet{poporangeL}{poporange!20}
\colorlet{popdarkL}{popdark!20}
\colorlet{poppurpleL}{poppurple!20}
\colorlet{popgreenL}{popgreen!20}
\colorlet{poppinkL}{poppink!20}
\colorlet{popblueL}{popblue!20}
\colorlet{popgoldL}{popgold!20}
\colorlet{popredL}{popred!20}
\colorlet{popgrayL}{popgray!20}
% Teintes claires pour fonds de tableaux / zones
\colorlet{popblueL}{popblue!10!white}
\colorlet{poporangeL}{poporange!14!white}
\colorlet{popgreenL}{popgreen!12!white}
\colorlet{poppurpleL}{poppurple!12!white}
\colorlet{poppinkL}{poppink!12!white}
\colorlet{popgoldL}{popgold!14!white}

\pagecolor{popcream}

% ============================================================
% Typographie
% ============================================================
\defaultfontfeatures{Ligatures=TeX}
\setmainfont{PlusJakartaSans-Medium.ttf}[Path=../../../fonts/,BoldFont=PlusJakartaSans-Bold.ttf]
% Maths en sans-serif (Fira Math) avec chiffres et lettres droites de Plus
% Jakarta, pour que les formules se fondent dans le texte.
\usepackage[math-style=ISO,bold-style=ISO]{unicode-math}
\setmathfont{FiraMath-Regular.otf}[Path=../../../fonts/]
\setmathfont{PlusJakartaSans-Medium.ttf}[Path=../../../fonts/,range={up/num,up/{Latin,latin},\mathup}]
\usepackage{icomma} % virgule décimale sans espace en mode math
% Caractères mathématiques Unicode absents des polices de texte (Plus Jakarta,
% Archivo Black) : rendus par leur équivalent mathématique, en texte comme en
% formule — sinon ils s'affichent en carré vide (ex. « Arithmétique dans ℤ »).
\usepackage{newunicodechar}
\newunicodechar{ℝ}{\ensuremath{\mathbb{R}}}
\newunicodechar{ℤ}{\ensuremath{\mathbb{Z}}}
\newunicodechar{ℂ}{\ensuremath{\mathbb{C}}}
\newunicodechar{ℕ}{\ensuremath{\mathbb{N}}}
\newunicodechar{ℚ}{\ensuremath{\mathbb{Q}}}
\newunicodechar{𝔻}{\ensuremath{\mathbb{D}}}
\newunicodechar{α}{\ensuremath{\alpha}}
\newunicodechar{β}{\ensuremath{\beta}}
\newunicodechar{γ}{\ensuremath{\gamma}}
\newunicodechar{δ}{\ensuremath{\delta}}
\newunicodechar{ε}{\ensuremath{\varepsilon}}
\newunicodechar{θ}{\ensuremath{\theta}}
\newunicodechar{λ}{\ensuremath{\lambda}}
\newunicodechar{μ}{\ensuremath{\mu}}
\newunicodechar{σ}{\ensuremath{\sigma}}
\newunicodechar{τ}{\ensuremath{\tau}}
\newunicodechar{φ}{\ensuremath{\varphi}}
\newunicodechar{ψ}{\ensuremath{\psi}}
\newunicodechar{ω}{\ensuremath{\omega}}
\newunicodechar{Σ}{\ensuremath{\Sigma}}
% majuscules produites par \MakeUppercase dans les titres (α-aminés -> Α)
\newunicodechar{Α}{\ensuremath{\alpha}}
\newunicodechar{Β}{\ensuremath{\beta}}
\newunicodechar{Γ}{\ensuremath{\Gamma}}
\newunicodechar{Δ}{\ensuremath{\Delta}}
\newunicodechar{Ε}{\ensuremath{\varepsilon}}
\newunicodechar{Θ}{\ensuremath{\Theta}}
\newunicodechar{Λ}{\ensuremath{\Lambda}}
\newunicodechar{Μ}{\ensuremath{\mu}}
\newunicodechar{Τ}{\ensuremath{\tau}}
\newunicodechar{Φ}{\ensuremath{\Phi}}
\newunicodechar{Ψ}{\ensuremath{\Psi}}
\newunicodechar{Ω}{\ensuremath{\Omega}}
\newunicodechar{↦}{\ensuremath{\mapsto}}
\newunicodechar{∈}{\ensuremath{\in}}
\newunicodechar{∉}{\ensuremath{\notin}}
\newunicodechar{∧}{\ensuremath{\wedge}}
\newunicodechar{⇔}{\ensuremath{\Leftrightarrow}}
\newunicodechar{⟺}{\ensuremath{\Longleftrightarrow}}
\newunicodechar{⇒}{\ensuremath{\Rightarrow}}
\newunicodechar{⟹}{\ensuremath{\Longrightarrow}}
\newunicodechar{∘}{\ensuremath{\circ}}
\newunicodechar{∪}{\ensuremath{\cup}}
\newunicodechar{∩}{\ensuremath{\cap}}
\newunicodechar{≡}{\ensuremath{\equiv}}
\newunicodechar{⊂}{\ensuremath{\subset}}
\newunicodechar{⊥}{\ensuremath{\perp}}
\newunicodechar{∃}{\ensuremath{\exists}}
\newunicodechar{∀}{\ensuremath{\forall}}
\newunicodechar{∛}{\ensuremath{\sqrt[3]{\,}}}
\newunicodechar{∜}{\ensuremath{\sqrt[4]{\,}}}
\newunicodechar{⃗}{\ensuremath{{}^{\to}}}
\newunicodechar{ⁿ}{\ifmmode^{n}\else\textsuperscript{\ensuremath{n}}\fi}
\newunicodechar{ˣ}{\ifmmode^{x}\else\textsuperscript{\ensuremath{x}}\fi}
\newunicodechar{ᵇ}{\ifmmode^{b}\else\textsuperscript{\ensuremath{b}}\fi}
\newunicodechar{ʳ}{\ifmmode^{r}\else\textsuperscript{\ensuremath{r}}\fi}
\newunicodechar{ᵘ}{\ifmmode^{u}\else\textsuperscript{\ensuremath{u}}\fi}
\newunicodechar{ʸ}{\ifmmode^{y}\else\textsuperscript{\ensuremath{y}}\fi}
\newunicodechar{ᵃ}{\ifmmode^{a}\else\textsuperscript{\ensuremath{a}}\fi}
\newunicodechar{ᵉ}{\ifmmode^{e}\else\textsuperscript{\ensuremath{e}}\fi}
\newunicodechar{ᵏ}{\ifmmode^{k}\else\textsuperscript{\ensuremath{k}}\fi}
\newunicodechar{ᵖ}{\ifmmode^{p}\else\textsuperscript{\ensuremath{p}}\fi}
\newunicodechar{ᴺ}{\ifmmode^{N}\else\textsuperscript{\ensuremath{N}}\fi}
\newunicodechar{ᶜ}{\ifmmode^{c}\else\textsuperscript{\ensuremath{c}}\fi}
\newunicodechar{ʰ}{\ifmmode^{h}\else\textsuperscript{\ensuremath{h}}\fi}
\newunicodechar{ᵠ}{\ifmmode^{q}\else\textsuperscript{\ensuremath{q}}\fi}
\newunicodechar{ₙ}{\ifmmode_{n}\else\textsubscript{\ensuremath{n}}\fi}
\newunicodechar{ₐ}{\ifmmode_{a}\else\textsubscript{\ensuremath{a}}\fi}
\newunicodechar{ᵢ}{\ifmmode_{i}\else\textsubscript{\ensuremath{i}}\fi}
\newunicodechar{ᵣ}{\ifmmode_{r}\else\textsubscript{\ensuremath{r}}\fi}
\newunicodechar{ₖ}{\ifmmode_{k}\else\textsubscript{\ensuremath{k}}\fi}
\newunicodechar{ᵦ}{\ifmmode_{\beta}\else\textsubscript{\ensuremath{\beta}}\fi}
\newunicodechar{ₓ}{\ifmmode_{x}\else\textsubscript{\ensuremath{x}}\fi}
\newfontfamily\popdisplay{ArchivoBlack-Regular.ttf}[Path=../../../fonts/]
\newfontfamily\poplogo{SpaceGrotesk-Bold.ttf}[Path=../../../fonts/]
\newfontfamily\popsemi{PlusJakartaSans-SemiBold.ttf}[Path=../../../fonts/]
\newfontfamily\popxbold{PlusJakartaSans-ExtraBold.ttf}[Path=../../../fonts/]
\newfontfamily\popxboldls{PlusJakartaSans-ExtraBold.ttf}[Path=../../../fonts/,LetterSpace=3.0]

\setlist[enumerate]{leftmargin=1.7em,itemsep=5pt,topsep=4pt}
\setlist[itemize]{leftmargin=1.7em,itemsep=4pt,topsep=4pt,label={\color{poporange}\textbullet}}
\setlength{\parindent}{0pt}
\setlength{\parskip}{5pt}
\renewcommand{\arraystretch}{1.25}

% pgfplots : style Pop par défaut
\pgfplotsset{
  every axis/.append style={axis line style={popink,line width=1pt},tick style={popink},
    grid style={popline,line width=0.5pt},label style={font=\small},tick label style={font=\footnotesize}},
  popcurve/.style={poporange,line width=1.8pt,no markers,smooth},
}
% Même style disponible dans un \draw TikZ simple (hors pgfplots)
\tikzset{popcurve/.style={poporange,line width=1.8pt}}
\tikzset{popgrid/.style={popline,very thin}}
\colorlet{popcurve}{poporange} % le nom du style est parfois utilisé comme couleur

% ============================================================
% Pill / squiggle
% ============================================================
\newcommand{\poppill}[3]{%
  \tikz[baseline=-0.55ex]\node[draw=popink,line width=1.2pt,rounded corners=2.6pt,%
    fill=#2,inner xsep=6.5pt,inner ysep=3pt]{%
    \popxboldls\fontsize{7}{7}\selectfont\color{#3}\MakeUppercase{#1}};%
}
\newcommand{\popsquiggle}[1]{%
  \tikz[baseline=-0.4ex]{
    \draw[#1,line width=2.6pt,line cap=round]
      plot [smooth] coordinates {(0,0.09) (0.34,0.01) (0.68,0.21) (1.02,0.01) (1.36,0.10)};
  }%
}
% Mot-clé surligné (marqueur orange sous le texte)
\newcommand{\cle}[1]{{\popxbold\color{popdark}#1}}

% ============================================================
% En-tête / pied
% ============================================================
\pagestyle{fancy}
\fancyhf{}
\renewcommand{\headrulewidth}{0pt}
\renewcommand{\footrulewidth}{0pt}
\lhead{\raisebox{-0.15ex}{\poplogo\fontsize{13}{13}\selectfont\color{popink}OnBuch\color{poporange}+}}
\rhead{\poppill{\DOCMATIERE\ \textbullet\ \DOCNIVEAU\ \textbullet\ Leçon \DOCLECON}{white}{popink}}
\lfoot{\popxbold\fontsize{7.6}{7.6}\selectfont\color{popmuted}\MakeUppercase{Cours signé OnBuch+}\ \textbullet\ {\popsemi\DOCTITRE}}
\rfoot{\tikz[baseline=-0.6ex]\node[rounded corners=8.5pt,fill=popink,minimum height=17pt,minimum width=17pt,inner xsep=5pt,inner sep=0]{\popxbold\fontsize{8}{8}\selectfont\color{popcream}\thepage\,/\,\pageref*{LastPage}};}

% ============================================================
% Titres
% ============================================================
\newcommand{\chaptitre}[2]{%
  \begin{tcolorbox}[enhanced,colback=poporange,colframe=popink,boxrule=2.6pt,arc=11pt,
      drop shadow=popink,left=15pt,right=15pt,top=12pt,bottom=11pt,before skip=3pt,after skip=18pt]
    \poppill{#2}{popink}{popcream}\par\vspace{9pt}
    {\raggedright\hyphenpenalty=10000\exhyphenpenalty=10000\popdisplay\fontsize{22}{24}\selectfont\color{popcream}\MakeUppercase{#1}\par}
  \end{tcolorbox}
}
% Section numérotée : pastille orange avec le numéro + titre Archivo
\newcounter{popsec}
\newcommand{\coursec}[1]{%
  \stepcounter{popsec}\setcounter{popsub}{0}%
  \par\vspace{16pt}\noindent
  \begin{minipage}{\linewidth}
  \tikz[baseline=-0.7ex]\node[draw=popink,line width=1.6pt,fill=poporange,rounded corners=4pt,minimum size=22pt,inner sep=0]{\popdisplay\fontsize{12}{12}\selectfont\color{popcream}\thepopsec};%
  \hspace{8pt}{\popdisplay\fontsize{15}{17}\selectfont\color{popink}\MakeUppercase{#1}}\par
  \vspace{4pt}\noindent{\color{poporange}\rule{36pt}{3.6pt}}
  \end{minipage}\par\nopagebreak\vspace{8pt}
}
\newcounter{popsub}
\newcommand{\courssub}[1]{%
  \stepcounter{popsub}%
  \par\vspace{9pt}\noindent{\popxbold\fontsize{11.5}{13}\selectfont\color{popink}{\color{poporange}\thepopsec.\thepopsub}\hspace{6pt}#1}\par\nopagebreak\vspace{4pt}
}

% ============================================================
% Boîtes pédagogiques — même recette : fond blanc, bordure épaisse colorée,
% ombre dure encre, badge plein. Titre optionnel [#1].
% ============================================================
\tcbset{popbase/.style={breakable,enhanced,colback=white,boxrule=2.3pt,arc=9pt,
  shadow={3pt}{-3pt}{0pt}{popink},left=14pt,right=14pt,top=11pt,bottom=11pt,before skip=11pt,after skip=15pt,
  fontupper=\popsemi\fontsize{10.6}{14.5}\selectfont\color{popink},
  fontlower=\popsemi\fontsize{10.6}{14.5}\selectfont\color{popink}}}
% Coche dessinée (le glyphe ✓ n'existe pas dans les polices du document)
\newcommand{\popcheck}[1]{\tikz[baseline=-0.5ex]\draw[#1,line width=1.6pt,line cap=round,line join=round](-0.13,0.0)--(-0.04,-0.09)--(0.13,0.1);}
\providecommand{\coche}{\popcheck{popgreen}}
\providecommand{\resultat}[1]{\par\smallskip\noindent\fcolorbox{popgreen}{popgreenL}{\parbox{\dimexpr\linewidth-2\fboxsep-2\fboxrule\relax}{\textbf{Résultat.} #1}}\par\smallskip}
\newcommand{\popboxhead}[3]{\poppill{#1}{#2}{white}\ifx\relax#3\relax\else\hspace{6pt}{\popxbold\fontsize{10.6}{12}\selectfont\color{popink}#3}\fi\par\vspace{7pt}}

\newtcolorbox{aretenir}[1][]{popbase,colframe=popgreen,before upper={\popboxhead{À retenir}{popgreen}{#1}}}
\newtcolorbox{exemplebox}[1][]{popbase,colframe=poporange,before upper={\popboxhead{Exemple}{poporange}{#1}}}
\newtcolorbox{definition}[1][]{popbase,colframe=popblue,before upper={\popboxhead{Définition}{popblue}{#1}}}
\newtcolorbox{propriete}[1][]{popbase,colframe=poppurple,before upper={\popboxhead{Propriété}{poppurple}{#1}}}
\newtcolorbox{methode}[1][]{popbase,colframe=popgold,before upper={\popboxhead{Méthode}{popgold}{#1}}}
\newtcolorbox{attention}[1][]{popbase,colframe=poppink,before upper={\popboxhead{Attention}{poppink}{#1}}}
\newtcolorbox{experience}[1][]{popbase,colframe=popblue,colback=popblueL,before upper={\popboxhead{Expérience}{popblue}{#1}}}
% « Le savais-tu ? » : carte violette pleine (contexte, culture, Cameroun)
\newtcolorbox{savaistu}[1][]{popbase,colback=poppurple,colframe=popink,
  fontupper=\popsemi\fontsize{10.4}{14.2}\selectfont\color{white},
  before upper={\poppill{Le savais-tu\,?}{white}{poppurple}\ifx\relax#1\relax\else\hspace{6pt}{\popxbold\color{white}#1}\fi\par\vspace{7pt}}}
% Exercice résolu : énoncé en haut, \tcblower puis la solution sur fond crème
\newcounter{popexo}\newcounter{popexr}
\newtcolorbox{exoresolu}[1][]{popbase,colframe=poporange,colbacklower=popcream,
  segmentation style={popink,line width=1.2pt,dashed},
  before upper={\stepcounter{popexr}\popboxhead{Exercice résolu \thepopexr}{poporange}{#1}},
  before lower={\poppill{Solution}{popink}{popcream}\par\vspace{6pt}}}
% Exercice (non résolu en place) — la correction est dans la partie Corrigés
\newtcolorbox{exercice}[1][]{popbase,colframe=popink,
  before upper={\stepcounter{popexo}\popboxhead{Exercice \thepopexo}{popink}{#1}}}
% Corrigé
\newtcolorbox{corrige}[1][]{popbase,colframe=popgreen,colback=popgreenL,
  before upper={\popboxhead{Corrigé}{popgreen}{#1}}}
% Objectifs / prérequis / bilan
\newtcolorbox{objectifs}{popbase,colframe=popink,colback=poporangeL,
  before upper={\popboxhead{Objectifs de la leçon}{popink}{}}}
\newtcolorbox{prerequis}{popbase,colframe=popink,colback=popblueL,
  before upper={\popboxhead{Prérequis}{popblue}{}}}
\newtcolorbox{bilan}{popbase,colframe=popink,colback=white,boxrule=2.8pt,
  before upper={\popboxhead{Fiche bilan}{poporange}{L'essentiel à connaître pour l'examen}}}
% Figure encadrée avec légende
\newtcolorbox{popfigure}[1][]{enhanced,breakable=false,colback=white,colframe=popline,boxrule=1.4pt,arc=8pt,
  left=10pt,right=10pt,top=10pt,bottom=8pt,before skip=10pt,after skip=12pt,halign=center,
  lower separated=false,halign lower=center,
  fontlower=\popsemi\fontsize{9}{11.5}\selectfont\color{popmuted},
  #1}
\newcommand{\legende}[1]{\par\vspace{6pt}{\popsemi\fontsize{9}{11.5}\selectfont\color{popmuted}#1}}

% ============================================================
% Signature OnBuch+
% ============================================================
\newcommand{\poplogoinline}[1]{{\poplogo\fontsize{#1}{#1}\selectfont\color{popink}OnBuch\color{poporange}+}}
\newcommand{\signatureonbuch}{%
  \begin{tcolorbox}[enhanced,colback=popink,colframe=popink,boxrule=2.2pt,arc=10pt,
      drop shadow=poporange,left=14pt,right=14pt,top=10pt,bottom=10pt,before skip=0pt,after skip=0pt]
    \tikz[baseline=-0.7ex]\node[circle,fill=popgreen,draw=popcream,line width=1.3pt,minimum size=22pt,inner sep=0]{\popcheck{white}};%
    \hspace{9pt}\begin{minipage}[c]{0.62\linewidth}
      {\popxbold\fontsize{12}{13}\selectfont\color{popcream}Signé\ \textbullet\ {\poplogo OnBuch\color{poporange}+}}\par\vspace{2pt}
      {\raggedright\popsemi\fontsize{8}{10.5}\selectfont\color{popline}\MakeUppercase{Équipe pédagogique OnBuch+}\\\MakeUppercase{Conforme au programme officiel MINESEC}\par}
    \end{minipage}\hfill
    {\poplogo\fontsize{22}{22}\selectfont\color{popcream}OnBuch\color{poporange}+}
  \end{tcolorbox}%
}

% ============================================================
% Page de garde
% #1 titre · #2 matière · #3 niveau · #4 module · #5 n° leçon · #6 durée ·
% #7 description / objectifs (texte ou itemize)
% ============================================================
\newcommand{\pagegarde}[7]{%
  \thispagestyle{empty}
  \newgeometry{a4paper,margin=1.7cm,top=1.6cm,bottom=1.4cm}%
  \begin{tikzpicture}[remember picture,overlay]
    \fill[poporange!18] ([xshift=-3.2cm,yshift=-3.6cm]current page.north east) circle (4.6cm);
    \fill[poppurple!14] ([xshift=2.4cm,yshift=7.6cm]current page.south west) circle (3.8cm);
  \end{tikzpicture}%
  {\poplogo\fontsize{30}{30}\selectfont\color{popink}OnBuch\color{poporange}+}\hfill
  \raisebox{0.6ex}{\poppill{Cours signé \textbullet\ Premium}{popink}{popcream}}\par
  \vspace{1pt}\popsquiggle{poppurple}\par
  \vspace{8pt}
  \begin{tcolorbox}[enhanced,colback=poporange,colframe=popink,boxrule=2.8pt,arc=13pt,
      drop shadow=popink,left=18pt,right=18pt,top=12pt,bottom=13pt,after skip=10pt]
    \poppill{#2 \textbullet\ #3}{popink}{popcream}\hspace{5pt}\poppill{#4}{white}{popink}\par\vspace{8pt}
    {\popdisplay\fontsize{14}{14}\selectfont\color{popink}LEÇON #5}\par\vspace{4pt}
    {\raggedright\hyphenpenalty=10000\exhyphenpenalty=10000\popdisplay\fontsize{25}{27}\selectfont\color{popcream}\MakeUppercase{#1}\par}
    \vspace{8pt}
    {\popxbold\fontsize{9.5}{9.5}\selectfont\color{popink}\MakeUppercase{Durée conseillée : #6}}
  \end{tcolorbox}
  \begin{tcolorbox}[enhanced,colback=white,colframe=popink,boxrule=2.2pt,arc=10pt,
      drop shadow=popink,left=16pt,right=16pt,top=10pt,bottom=10pt,after skip=8pt]
    \poppill{Dans cette leçon}{poppurple}{white}\par\vspace{5pt}
    \popsemi\fontsize{9.3}{12.6}\selectfont\color{popink}#7
  \end{tcolorbox}
  \noindent\begin{minipage}[t]{0.487\linewidth}
    \begin{tcolorbox}[enhanced,colback=popgreen,colframe=popink,boxrule=2.2pt,arc=10pt,
        drop shadow=popink,left=11pt,right=11pt,top=10pt,bottom=10pt,height=3.1cm,valign=center]
      \tcbox[colback=white,colframe=popink,boxrule=1.3pt,arc=5pt,boxsep=0pt,nobeforeafter,
        left=3pt,right=3pt,top=3pt,bottom=3pt,valign=center]{\qrcode[height=1.9cm]{https://play.google.com/store/apps/details?id=com.luvvix.onbuch&hl=fr}}%
      \hspace{8pt}\begin{minipage}[c]{0.5\linewidth}
        {\popdisplay\fontsize{11}{12.5}\selectfont\color{white}\MakeUppercase{Télécharge OnBuch}}\par\vspace{4pt}
        {\raggedright\popsemi\fontsize{8.4}{11}\selectfont\color{white}Gratuit \textbullet\ Google Play\\Tuteur IA, annales, concours, résultats\par}
      \end{minipage}
    \end{tcolorbox}
  \end{minipage}\hfill
  \begin{minipage}[t]{0.487\linewidth}
    \begin{tcolorbox}[enhanced,colback=poppurple,colframe=popink,boxrule=2.2pt,arc=10pt,
        drop shadow=popink,left=11pt,right=11pt,top=10pt,bottom=10pt,height=3.1cm,valign=center]
      \tikz[baseline=-0.7ex]\node[circle,fill=white,draw=popink,line width=1.3pt,minimum size=1.6cm,inner sep=0]{\popdisplay\fontsize{24}{24}\selectfont\color{poppurple}+};%
      \hspace{8pt}\begin{minipage}[c]{0.6\linewidth}
        {\popdisplay\fontsize{11}{12.5}\selectfont\color{white}\MakeUppercase{Passe à OnBuch+}}\par\vspace{4pt}
        {\raggedright\popsemi\fontsize{8.4}{11}\selectfont\color{white}Tous les cours signés, Tuteur IA illimité, fiches PDF — onglet OnBuch+ de l'app\par}
      \end{minipage}
    \end{tcolorbox}
  \end{minipage}
  \vspace{8pt}
  \popcontenu
  \vfill
  \signatureonbuch
  \restoregeometry
  \newpage
}

% Rangée « Ce cours contient » (page de garde)
\newcommand{\poptile}[3]{% #1 couleur #2 gros texte #3 légende
  \begin{tcolorbox}[enhanced,colback=#1,colframe=popink,boxrule=2pt,arc=9pt,shadow={2.5pt}{-2.5pt}{0pt}{popink},
      left=6pt,right=6pt,top=9pt,bottom=9pt,halign=center,valign=center,height=1.75cm,width=\linewidth,nobeforeafter]
    {\popdisplay\fontsize{12.5}{13}\selectfont\color{white}\MakeUppercase{#2}}\par\vspace{3pt}
    {\centering\popsemi\fontsize{7.8}{9.5}\selectfont\color{white}#3\par}
  \end{tcolorbox}}
\newcommand{\popcontenu}{%
  \par\noindent{\popxboldls\fontsize{7.5}{7.5}\selectfont\color{popmuted}CE COURS SIGNÉ CONTIENT}\par\vspace{7pt}
  \noindent\begin{minipage}[t]{0.235\linewidth}\poptile{popblue}{Cours}{complet, illustré, pas à pas}\end{minipage}\hfill
  \begin{minipage}[t]{0.235\linewidth}\poptile{poporange}{Méthodes}{et exercices résolus}\end{minipage}\hfill
  \begin{minipage}[t]{0.235\linewidth}\poptile{poppink}{Exercices}{type examen + corrigés}\end{minipage}\hfill
  \begin{minipage}[t]{0.235\linewidth}\poptile{popgreen}{Fiche bilan}{l'essentiel à retenir}\end{minipage}\par}

% Sommaire
\newtcolorbox{sommaire}{enhanced,colback=white,colframe=popink,boxrule=2.2pt,arc=10pt,
  drop shadow=popink,left=18pt,right=18pt,top=13pt,bottom=13pt,before skip=6pt,after skip=20pt,
  before upper={\poppill{Sommaire}{poppurple}{white}\par\vspace{9pt}\popsemi\fontsize{10.6}{14.5}\selectfont\color{popink}}}

% Signature de fin de cours — poussée tout en bas de la dernière page
\newcommand{\findecours}{%
  \par\vfill\nopagebreak
  \begin{center}\popsquiggle{poporange}\end{center}\vspace{4pt}
  \signatureonbuch
}
\AtBeginDocument{\def\llbracket{\mathopen{[\mkern-3.5mu[}}\def\rrbracket{\mathclose{]\mkern-3.5mu]}}} % intervalles d'entiers (glyphe ⟦ absent de la police)
% Glyphes absents de Fira Math : redéfinis à partir de glyphes disponibles
\AtBeginDocument{%
  \def\setminus{\mathbin{\backslash}}\let\smallsetminus\setminus
  \def\vdots{\mathord{\vbox{\baselineskip3.2pt\lineskiplimit0pt\kern4pt\hbox{.}\hbox{.}\hbox{.}}}}%
  \def\ddots{\mathinner{\mkern1mu\raise6.5pt\hbox{.}\mkern2mu\raise3.6pt\hbox{.}\mkern2mu\raise0.7pt\hbox{.}\mkern1mu}}%
  \def\triangle{\mathord{\scalebox{0.9}{$\symup{\Delta}$}}}%
  \def\nabla{\mathord{\raisebox{\depth}{\scalebox{1}[-1]{$\symup{\Delta}$}}}}%
  \def\checkmark{\popcheck{popdark}}%
  \def\star{\mathbin{\ast}}\def\bigstar{\mathord{\ast}}%
  \def\diamond{\mathbin{\scalebox{0.6}{$\Diamond$}}}%
  \def\bigcup{\mathop{\mathchoice{\raisebox{-0.3em}{\scalebox{1.7}{$\cup$}}}{\scalebox{1.25}{$\cup$}}{\cup}{\cup}}\displaylimits}%
  \def\bigcap{\mathop{\mathchoice{\raisebox{-0.3em}{\scalebox{1.7}{$\cap$}}}{\scalebox{1.25}{$\cap$}}{\cap}{\cap}}\displaylimits}%
  \def\lesseqqgtr{\mathrel{\lessgtr}}\def\bigoplus{\mathop{\oplus}\displaylimits}%
}
\providecommand{\ssi}{si et seulement si} % abréviation fréquente
\AtBeginDocument{\providecommand{\wideparen}{\overparen}} % arc de cercle

%% FIGFIX-BEGIN v5 — correctifs globaux des figures (docs/onbuch-generation/tools/figfix)
\usepackage{adjustbox}\usepackage{etoolbox}\tcbuselibrary{raster}
% toute figure TikZ/pgfplots plus large que la ligne est réduite à la largeur disponible
\BeforeBeginEnvironment{tikzpicture}{\begin{adjustbox}{max width=\linewidth}}
\AfterEndEnvironment{tikzpicture}{\end{adjustbox}}
% légendes pgfplots lisibles par-dessus les courbes
\pgfplotsset{every axis/.append style={legend style={fill=white,fill opacity=0.92,text opacity=1,draw=black!15,inner sep=3pt}}}
% étiquettes d'arêtes (syntaxe edge["..."]) avec fond blanc
\tikzset{every edge quotes/.append style={fill=white,inner sep=1.2pt}}
%% FIGFIX-END
\begin{document}

\pagegarde{Dissertation : produire une dissertation}{Français}{Tle Tech}{Français – classe de Terminale technique}{N02}{15 séances (fiche de progression harmonisée)}{Cette leçon de 15 séances guide l'élève de Terminale Technique dans l'acquisition des compétences fondamentales pour réussir l'épreuve de dissertation au Baccalauréat : analyse approfondie de l'œuvre au programme, maîtrise des outils de liaison syntaxique, compréhension des stratégies argumentatives, et surtout, appropriation pas à pas de la méthode dissertationnaire (problématique, plan, introduction, développement, conclusion, réécriture).
\vspace{4pt}
\begin{multicols}{2}\small
\begin{itemize}
\item À la fin de cette leçon, je sais lire méthodiquement une œuvre littéraire (Le Vieux nègre et la médaille) en identifiant les thèmes, les personnages et la portée sociale.
\item Je sais identifier et classifier les outils de liaison (conjonctions, adverbes, pronoms relatifs) pour assurer la cohérence de mes écrits.
\item Je sais analyser les stratégies argumentatives (thèses, arguments, procédés rhétoriques) dans un texte de presse ou littéraire.
\item Je sais construire une introduction de dissertation type (accroche, sujet posé, sujet divisé, problématique, annonce du plan).
\item Je sais formuler une problématique pertinente et élaborer un plan détaillé dialectique ou thématique équilibré.
\item Je sais rédiger un paragraphe argumenté respectant la structure : idée directrice, argument(s), exemple(s) précis, lien/transition.
\end{itemize}\end{multicols}}

\begin{sommaire}
\begin{multicols}{2}
\begin{enumerate}
\item 1. « Le Vieux nègre et la médaille » : Lecture méthodique et appropriation de l'œuvre (4 séances)
\item 2. Maîtrise de la langue : Les outils de liaison pour une cohérence textuelle (2 séances)
\item 3. Le texte argumentatif : Stratégies et analyse (1 séance)
\item 4. La dissertation : Architecture de l'introduction à la conclusion (3 séances)
\item 5. De la problématique à la rédaction finale : Processus complet d'élaboration (4 séances)
\item 6. Bilan, Compte-rendu et Remédiation individualisée (1 séance)
\item Activité d'intégration
\item Méthodes et exercices résolus
\item Exercices
\item Corrigés des exercices
\item Fiche bilan
\end{enumerate}
\end{multicols}
\end{sommaire}

\begin{objectifs}
\begin{itemize}
\item À la fin de cette leçon, je sais lire méthodiquement une œuvre littéraire (Le Vieux nègre et la médaille) en identifiant les thèmes, les personnages et la portée sociale.
\item Je sais identifier et classifier les outils de liaison (conjonctions, adverbes, pronoms relatifs) pour assurer la cohérence de mes écrits.
\item Je sais analyser les stratégies argumentatives (thèses, arguments, procédés rhétoriques) dans un texte de presse ou littéraire.
\item Je sais construire une introduction de dissertation type (accroche, sujet posé, sujet divisé, problématique, annonce du plan).
\item Je sais formuler une problématique pertinente et élaborer un plan détaillé dialectique ou thématique équilibré.
\item Je sais rédiger un paragraphe argumenté respectant la structure : idée directrice, argument(s), exemple(s) précis, lien/transition.
\item Je sais rédiger une conclusion synthétique (réponse à la problématique, ouverture) et relire mon devoir pour corriger la syntaxe et l'orthographe.
\item Je sais mobiliser mes connaissances sur l'œuvre et la langue pour traiter un sujet de dissertation complet en temps limité (épreuve type Bac).
\end{itemize}
\end{objectifs}

\begin{prerequis}
\begin{itemize}
\item Structure de base du paragraphe (idée force / explication / illustration) acquise en Première.
\item Notion de phrase simple et complexe (propositions principales / subordonnées) vue en classe de Seconde.
\item Définition de la thèse, de l'argument et de l'exemple dans le cadre de l'argumentation directe.
\item Connaissance du contexte historique et social du Cameroun post-indépendance (années 60-70) pour situer l'œuvre.
\item Maîtrise minimale des connecteurs logiques de base (donc, mais, parce que, si).
\end{itemize}
\end{prerequis}

\newpage

\coursec{« Le Vieux nègre et la médaille » : Lecture méthodique et appropriation de l'œuvre (4 séances)}

Au carrefour de la rue de l'Indépendance à Douala, un vieil homme médaillé tend la main : sa médaille de la valeur ne paie pas le pain. Comment, dans une dissertation, articuler l'indignation face à l'ingratitude de l'État et la dignité humaine bafouée, sans tomber dans le simple pathos ? C'est tout l'enjeu de la maîtrise de l'argumentation et de la structure dissertationnaire que nous abordons ici, à travers le prisme du \emph{Vieux nègre et la médaille}.

\begin{popfigure}
\begin{tikzpicture}[node distance=2.5cm, auto, every node/.style={font=\small}]
\node[draw, rectangle, rounded corners, fill=popblue!15, thick, minimum width=5cm, minimum height=1.2cm] (titre) {\textbf{Le Vieux nègre et la médaille}};
\node[draw, rectangle, rounded corners, fill=popgreen!15, thick, below of=titre, minimum width=6cm, align=center] (axes){\textbf{Axes d'étude}\\1. Dénonciation du colonialisme\\2. Éveil de la conscience\\3. Ironie tragique};
\node[draw, ellipse, fill=poporange!15, thick, left of=axes, xshift=-3cm, minimum width=3.5cm, align=center] (perso){\textbf{Personnages}\\Meka, Kelara, Commandant, Prêtre};
\node[draw, ellipse, fill=poppurple!15, thick, right of=axes, xshift=3cm, minimum width=3.5cm, align=center] (themes){\textbf{Thèmes majeurs}\\Ingratitude, Dignité, Choc des cultures};
\node[draw, rectangle, rounded corners, fill=poppink!15, thick, below of=axes, yshift=-1cm, minimum width=5cm, align=center] (style){\textbf{Procédés stylistiques}\\Ironie, Oralité, Proverbes, Structure circulaire};
\draw[->, thick, popblue] (titre) -- (axes);
\draw[->, thick, popgreen] (axes) -- (perso);
\draw[->, thick, popgreen] (axes) -- (themes);
\draw[->, thick, poporange] (axes) -- (style);
\end{tikzpicture}
\legende{Carte mentale de l'œuvre : axes d'étude, personnages, thèmes et procédés stylistiques.}
\end{popfigure}

\courssub{Présentation de l'auteur et du contexte historique (Séance 1)}

\begin{savaistu}[Ferdinand Oyono : un intellectuel au service de la nation]
Ferdinand \textbf{Oyono} (1929--2010) est une figure majeure des lettres et de la politique camerounaises. Né à Ngoulemakong, il poursuit des études en France où il rédige ses trois romans : \emph{Une vie de boy} (1956), \emph{Le Vieux nègre et la médaille} (1956) et \emph{Chemin d'Europe} (1960). Diplomate de carrière, il fut ambassadeur puis ministre sous Ahmadou Ahidjo et Paul Biya. \emph{Le Vieux nègre et la médaille}, publié aux Éditions Julliard dans la collection « Lettres nouvelles », est souvent considéré comme le \textbf{premier grand roman camerounais de langue française}. Il obtient le prix Sainte-Beuve en 1957.
\end{savaistu}

\begin{definition}[Contexte de publication et portée historique]
L'année 1956 marque l'apogée de la guerre d'Algérie et la veille des indépendances africaines. Au Cameroun, l'Union des populations du Cameroun (UPC) mène une lutte armée réprimée dans le sang par l'administration coloniale française. Le roman d'Oyono paraît dans ce climat : c'est une \textbf{satire féroce de l'administration coloniale}, qui dévoile les mécanismes de la domination par le rire et l'ironie. L'œuvre ne se contente pas de dénoncer ; elle montre comment le colonisé intériorise sa soumission avant de s'en affranchir par une prise de conscience brutale.
\end{definition}

\begin{exemplebox}[Repères chronologiques pour le dossier Bac]
\begin{itemize}
\item \textbf{1929} : Naissance à Ngoulemakong (région du Sud, Cameroun).
\item \textbf{1956} : Publication simultanée de \emph{Une vie de boy} et \emph{Le Vieux nègre et la médaille}.
\item \textbf{1960} : Indépendance du Cameroun ; Oyono rentre au pays.
\item \textbf{1961--1990} : Carrière diplomatique (ambassadeur en Grande-Bretagne, aux Nations unies, en Algérie, etc.) et ministérielle (Fonction publique, Éducation, Affaires étrangères).
\item \textbf{2010} : Décès à Yaoundé.
\end{itemize}
\end{exemplebox}

\courssub{Lecture méthodique n°1 : Titre, incipit et cadre spatio-temporel (Séance 1--2)}

\begin{definition}[Lecture méthodique]
La \cle{lecture méthodique} est une lecture guidée par des questionnaires précis qui font émerger le sens littéral puis implicite du texte. Elle procède par étapes : \textbf{observation} (que vois-je ?), \textbf{analyse} (comment est-ce construit ?), \textbf{interprétation} (que signifie cela ?). Elle évite l'impressionnisme et ancre l'analyse dans le texte.
\end{definition}

\textbf{Analyse du titre.} « Le Vieux nègre et la médaille » : trois éléments s'opposent et se répondent. \emph{Le Vieux nègre} : désignation déshumanisante, raciste, qui réduit l'identité à l'âge et à la race (le terme « nègre » est assumé par Oyono comme marqueur colonial). \emph{La médaille} : symbole de la reconnaissance officielle, métal froid, valeur symbolique vide. \emph{Et} : conjonction de coordination qui lie deux réalités inconciliables — l'homme et l'objet, la chair et le métal, la vie et le symbole.

\textbf{Étude de l'incipit (premier chapitre).} Le roman s'ouvre sur la préparation de la cérémonie : « Meka se leva avant l'aube [...] Il voulait être prêt quand le soleil se lèverait. » Le cadre spatio-temporel se déploie :
\begin{itemize}
\item \textbf{Espace} : le village (Doum), la case de Meka, la place publique où se tiendra la cérémonie. L'espace est clos, villageois, régi par les coutumes.
\item \textbf{Temps} : l'aube d'un jour de fête, temps suspendu de l'attente. Le temps colonial (la cérémonie officielle) va s'imposer au temps villageois.
\item \textbf{Atmosphère} : excitation fébrile, fierté naïve. Meka « polissait sa médaille » comme un objet sacré.
\end{itemize}

\begin{methode}[Comment analyser un incipit ?]
\begin{enumerate}
\item Identifier les \textbf{indices spatio-temporels} (adverbes, compléments de lieu/temps, verbes de mouvement).
\item Repérer le \textbf{point de vue narratif} (qui voit ? qui parle ?).
\item Relever les \textbf{indices tonaux} (lexique, figures de style) qui annoncent la couleur du roman.
\item Formuler une \textbf{hypothèse de lecture} : quel contrat de lecture l'incipit propose-t-il au lecteur ?
\end{enumerate}
\end{methode}

\begin{exoresolu}[Analyse de l'incipit]
\textbf{Énoncé :} À partir des premières lignes du roman (jusqu'à « Il était fier »), relevez les éléments qui montrent la naïveté de Meka et l'ironie narrative.

\textbf{Solution détaillée :}
\begin{enumerate}
\item \textbf{Le lexique de la propreté et de la préparation} : « se lava soigneusement », « polissait sa médaille », « mit son plus beau pagne ». Meka traite la médaille comme une relique, un objet sacré.
\item \textbf{Le point de vue interne} : le narrateur adopte la focalisation de Meka (« Il était fier », « Il pensait »). Nous voyons le monde \emph{à travers} ses yeux naïfs.
\item \textbf{L'ironie dramatique} : le lecteur, lui, sait (par le titre, par le genre romanesque) que cette fierté est vouée à l'échec. L'écart entre la conscience du personnage et la connaissance du lecteur crée l'ironie.
\item \textbf{La symbolique de l'aube} : « avant l'aube », « quand le soleil se lèverait » — l'aube annonce normalement un renouveau ; ici, elle précède la chute. L'ironie est déjà à l'œuvre dans la structure temporelle.
\end{enumerate}
\end{exoresolu}

\begin{attention}[Piège à éviter : Confondre focalisation et narration]
Ne dites pas « le narrateur est ironique » sans nuance. Précisez : « Le narrateur \emph{adopte une focalisation interne sur Meka} qui rend visible sa naïveté, tandis que \emph{la distance narrative} (le ton, le choix des détails) installe une ironie que Meka ne perçoit pas. » L'ironie naît de l'écart entre les deux instances.
\end{attention}

\courssub{Lecture méthodique n°2 : Personnages et relations de pouvoir (Séance 2)}

\begin{center}
\begin{tabularx}{\linewidth}{|l|X|X|}
\hline
\rowcolor{popblueL}
\textbf{Personnage} & \textbf{Portrait physique / social} & \textbf{Fonction symbolique / relation de pouvoir} \\
\hline
\textbf{Meka} & Vieux chef de village, ancien tirailleur, médaillé, polygame, respecté localement. & Figure du \textbf{colonisé loyal} qui a cru au pacte colonial. Sa trajectoire : naïveté $\to$ désillusion $\to$ révolte muette. \\
\hline
\textbf{Kelara} & Première épouse de Meka, femme lucide, pragmatique, attachée aux biens matériels. & \textbf{Voix de la raison}, du réalisme. Elle pressent l'arnaque coloniale (« À quoi ça sert, ta médaille ? »). Incarnation de la sagesse populaire. \\
\hline
\textbf{Le Commandant} & Administrateur colonial, cérémonial, distant, paternaliste. & \textbf{Visage du pouvoir colonial} : la médaille est un outil de domination, non de reconnaissance. Il incarne l'institution qui broie l'individu. \\
\hline
\textbf{Le Prêtre (Père Vandermayer)} & Missionnaire, médiateur culturel, ambigu. & \textbf{Relais idéologique} : la religion cautionne l'ordre colonial. Sa bénédiction sacralise la médaille. \\
\hline
\end{tabularx}
\end{center}

\textbf{Relations de pouvoir analysées :}
\begin{itemize}
\item \textbf{Meka / Commandant} : relation asymétrique. Meka attend reconnaissance et justice ; le Commandant octroie une distinction honorifique qui ne change rien à la condition matérielle. Le pouvoir colonial \emph{donne pour mieux asservir}.
\item \textbf{Meka / Kelara} : relation de complémentarité conflictuelle. Kelara voit clair ; Meka veut croire. Leur dialogue (chapitre 2) est le lieu où s'exprime le conflit entre \emph{illusion coloniale} et \emph{réalité vécue}.
\item \textbf{Meka / ses fils} : l'échec de la transmission. Les fils sont partis, l'un à l'école (aliénation culturelle), l'autre en ville (prolétarisation). La médaille ne sauve pas la lignée.
\end{itemize}

\begin{exemplebox}[Citation clé pour le dossier Bac -- Relation Meka/Kelara]
\begin{quote}
« -- À quoi ça sert, ta médaille ? demanda Kelara. \\
-- C'est une marque d'honneur, répondit Meka. \\
-- L'honneur ne mange pas, dit Kelara. »
\end{quote}
Cette réplique en deux versets (structure biblique, oralité) condense le conflit central : \textbf{valeur symbolique vs valeur d'usage}, \textbf{honneur colonial vs dignité vitale}.
\end{exemplebox}

\courssub{Lecture méthodique n°3 : Thèmes majeurs (Séance 3)}

\begin{aretenir}[Quatre thèmes structurants pour la dissertation]
\begin{enumerate}
\item \textbf{L'ingratitude coloniale (ou la « médaille vide »)} : La médaille récompense un service (guerre 14--18, 39--45) mais ne garantit ni pension, ni soins, ni respect. Le colonisé a donné son sang ; la métropole lui rend du métal. C'est le thème central, politique et éthique.
\item \textbf{La naïveté et le réveil de la conscience} : Meka n'est pas dupe au début ? Il \emph{veut} croire. Le roman suit une \textbf{courbe d'apprentissage inversée} : le vieux « apprend » qu'il a été dupé. L'éveil est douloureux, solitaire, sans issue révolutionnaire — mais il existe.
\item \textbf{Le choc des cultures} : Cérémonie coloniale (discours, défilé, médaille) vs coutumes villageoises (libations, danses, paroles des anciens). Le colonisateur impose son rituel ; le colonisé le subit en le détournant (ironie des chants, regards complices).
\item \textbf{L'ironie tragique} : Le comble : Meka meurt (ou sombre dans le silence) le jour de sa « gloire ». La médaille, censée l'élever, devient son linceul symbolique. L'ironie n'est pas seulement stylistique : elle est \textbf{structurelle}, elle organise le destin du personnage.
\end{enumerate}
\end{aretenir}

\begin{exoresolu}[Repérage thématique]
\textbf{Énoncé :} Relevez dans le chapitre de la cérémonie (chapitre 4) trois passages illustrant le thème de l'ingratitude coloniale.

\textbf{Solution détaillée :}
\begin{enumerate}
\item \textbf{Le discours du Commandant} : « La France n'oublie pas ses enfants. » \emph{Analyse} : formule creuse, rituelle. « Enfants » = paternalisme infantilisant. La France \emph{a} oublié : pas de pension, pas de médecin.
\item \textbf{La remise de la médaille} : « Il épingla la médaille sur la poitrine de Meka. » \emph{Analyse} : geste mécanique, clinique. La médaille pèse sur la chair comme une charge, pas comme un honneur.
\item \textbf{Le banquet officiel} : « Le champagne coulait à flots pour les officiels ; le vin de palme pour les indigènes. » \emph{Analyse} : ségrégation matérialisée par la boisson. L'apartheid festif révèle la vérité du pouvoir.
\end{enumerate}
\end{exoresolu}

\courssub{Lecture méthodique n°4 : Style, oralité et structure narrative (Séance 4)}

\begin{propriete}[Procédés stylistiques majeurs -- Exigibles au Bac]
\begin{itemize}
\item \textbf{L'ironie} (figure de style phare) : dire le contraire de ce qu'on pense, ou laisser un écart entre l'énoncé et la situation. Ex. : « La France est généreuse » (pensée de Meka) vs réalité du banquet ségrégé.
\item \textbf{Le sous-entendu (l'ellipse significative)} : ce qui n'est pas dit pèse plus que ce qui est dit. Les silences de Meka à la fin, les regards échangés entre villageois.
\item \textbf{L'oralité et les proverbes} : « L'honneur ne mange pas » (Kelara), « Le singe ne voit pas sa propre fesse » (proverbe inséré). Ils ancrent le roman dans la culture orale béti/bassa et donnent une portée universelle aux situations.
\item \textbf{La structure narrative circulaire} : Le roman s'ouvre sur l'attente de la cérémonie et se referme sur le retour à la case, la médaille rangée, le silence. Le cercle est bouclé, mais \emph{la conscience a changé}. Structure : \textbf{Naïveté $\to$ Cérémonie (point culminant) $\to$ Désillusion $\to$ Silence lucide}.
\end{itemize}
\end{propriete}

\begin{methode}[Analyser un passage pour le commentaire ou la dissertation]
\begin{enumerate}
\item \textbf{Contextualiser} : où sommes-nous dans l'intrigue ? quel personnage ? quel enjeu ?
\item \textbf{Segmenter} : repérer les unités de sens (paragraphes, séquences dialoguées, descriptions).
\item \textbf{Repérer les procédés} : lexique, syntaxe, figures, rythmes, registres (tragique, ironique, pathétique).
\item \textbf{Interpréter} : en quoi ces procédés servent-ils le thème ? quelle vision du monde transmettent-ils ?
\item \textbf{Conclure} : ouvrir sur la portée générale de l'œuvre.
\end{enumerate}
\end{methode}

\begin{exemplebox}[Exemple d'analyse stylistique -- Dernier paragraphe du roman]
\textit{« Meka rentra chez lui. Il posa la médaille sur la table. Il s'assit dans son fauteuil. Il ne dit rien. La nuit tomba sur le village. »}

\begin{itemize}
\item \textbf{Syntaxique} : phrases courtes, nominales ou verbales simples, parataxe (juxtaposition). Rythme haché = fatigue, vide, fin.
\item \textbf{Lexical} : verbes d'action banale (« posa », « s'assit ») contrastant avec l'événement exceptionnel. « La médaille sur la table » : l'objet sacré redevient objet banal.
\item \textbf{Symbolique} : « La nuit tomba » = fin du jour de gloire, mais aussi métaphore de la conscience qui s'obscurcit / s'éclaircit (ambiguïté féconde).
\item \textbf{Portée} : le silence de Meka est sa première parole libre. Il ne remercie pas, ne maudit pas : il \emph{constate}.
\end{itemize}
\end{exemplebox}

\courssub{Constitution de la fiche de lecture synthétique pour le dossier Bac (Travail transversal sur les 4 séances)}

\begin{methode}[Fiche Bac : 4 étapes obligatoires]
\begin{enumerate}
\item \textbf{Fiche d'identité} : Auteur (Ferdinand Oyono), Date (1956), Genre (Roman, satire postcoloniale), Collection (Julliard, « Lettres nouvelles »), Contexte (fin de l'époque coloniale, guerre d'Algérie, nationalisme camerounais).
\item \textbf{Résumé en 10 lignes maximum} : Meka, vieux tirailleur décoré, attend la cérémonie où le Commandant doit lui remettre la médaille de la Légion d'honneur. Fierté naïve, préparation minutieuse. Kelara, sa femme, doute de l'utilité de la médaille. Cérémonie fastueuse mais ségrégée. Discours creux du Commandant. Retour au village : Meka, silencieux, range sa médaille. La nuit tombe. Il a compris que l'honneur colonial ne nourrit pas son peuple.
\item \textbf{3 axes d'étude avec 2 citations chacun} :
\begin{itemize}
\item Axe 1 -- \emph{La médaille comme symbole de l'imposture coloniale} : « La France n'oublie pas ses enfants » (Commandant) / « L'honneur ne mange pas » (Kelara).
\item Axe 2 -- \emph{L'éveil douloureux d'une conscience} : « Il pensait à ses fils partis » (focalisation sur Meka) / « Il ne dit rien » (dernière phrase).
\item Axe 3 -- \emph{L'ironie comme arme de résistance} : « Le champagne pour les officiels, le vin de palme pour les indigènes » / Proverbe « Le singe ne voit pas sa propre fesse ».
\end{itemize}
\item \textbf{Portée personnelle / universelle} : L'œuvre dépasse le Cameroun des années 50. Elle interroge toute société sur la \textbf{reconnaissance due aux aînés, aux combattants, aux travailleurs de l'ombre}. Au Cameroun aujourd'hui, la question des anciens combattants (vétérans des guerres d'Indochine, d'Algérie, ou des crises récentes) et de leur prise en charge sociale fait écho direct au roman.
\end{enumerate}
\end{methode}

\begin{exemplebox}[Plan de travail pour l'appropriation complète]
L'œuvre compte environ \textbf{120 pages} (édition Julliard / Livre de Poche). Prévoyez :
\begin{itemize}
\item \textbf{3h de lecture personnelle} (à raison de 40 pages/heure en lecture cursive, plus lente pour l'annotation).
\item \textbf{4h de cours} pour les 4 lectures méthodiques guidées (1h par lecture).
\item \textbf{2h de révision} pour constituer la fiche Bac et mémoriser les 6 citations clés.
\end{itemize}
Total : \textbf{9h} environ. Commencez la lecture personnelle \textbf{avant} la première séance.
\end{exemplebox}

\begin{attention}[Erreur fréquente au Bac : Confondre résumé et analyse]
\begin{itemize}
\item \textbf{Résumé} : « Meka reçoit sa médaille, il est content, puis il rentre chez lui déçu. » $\to$ Récit de l'intrigue, pas d'analyse.
\item \textbf{Analyse} : « La remise de la médaille (chap. 4) met en scène l'ironie tragique : le discours du Commandant ("La France n'oublie pas ses enfants") entre en contradiction flagrante avec la ségrégation du banquet (champagne vs vin de palme), révélant que la distinction honorifique masque l'exploitation réelle. » $\to$ Citation + procédure + interprétation + lien au thème.
\end{itemize}
\textbf{Règle d'or} : En dissertation comme en commentaire, \textbf{chaque affirmation doit être étayée par une référence textuelle précise} (citation ou renvoi de page).
\end{attention}

\courssub{Mise en relation avec l'actualité camerounaise : mémoire et reconnaissance}

\begin{savaistu}[Actualité : Les anciens combattants camerounais aujourd'hui]
Le Cameroun compte encore des vétérans de la Seconde Guerre mondiale, des guerres d'Indochine et d'Algérie, ainsi que d'anciens militaires engagés dans les opérations de maintien de la paix (MINUSCA, BIR, etc.). La \textbf{Journée nationale du vétéran} (instituée en 2019, célébrée le 8 mars) et la \textbf{Caisse nationale de prévoyance sociale (CNPS)} tentent de répondre à l'ingratitude dénoncée par Oyono. Pourtant, des reportages montrent encore des anciens combattants dans la précarité, médailles sur la poitrine, quémandant une pension ou des soins. Le roman reste \textbf{prophétique} : la médaille ne paie toujours pas le pain.
\end{savaistu}

\begin{exercice}[Sujet de dissertation type Bac -- Entraînement]
\textbf{Sujet :} « La médaille est un leurre ; la vraie reconnaissance est dans le partage. » À partir de \emph{Le Vieux nègre et la médaille} et de votre connaissance du monde contemporain, discutez cette affirmation.

\textbf{Pistes de problématisation (à faire en classe) :}
\begin{enumerate}
\item En quoi la médaille coloniale est-elle un \emph{leurre} (illusion, manipulation, outil de domination) ?
\item Qu'est-ce que la \emph{vraie reconnaissance} pour Meka ? Pour Kelara ? Pour le village ?
\item Le « partage » (matériel, symbolique, mémoriel) peut-il remplacer la distinction honorifique individuelle ?
\item Ouverture : politiques publiques actuelles (pensions, santé, mémoire) -- sont-elles à la hauteur ?
\end{enumerate}
\end{exercice}

\begin{corrige}[Exercice n°1 -- Éléments de corrigé type]
\textbf{Introduction} : Accroche sur l'actualité (vétéran médaillé mendiant). Définition des termes : leurre = tromperie ; reconnaissance = acte par lequel on admet la valeur de quelqu'un ; partage = distribution équitable. Problématique : La distinction honorifique individuelle suffit-elle à réparer l'injustice structurelle ? Plan dialectique : I. La médaille comme leurre colonial (instrument de pouvoir, vide matériel). II. La quête d'une reconnaissance véritable (besoins vitaux, dignité collective). III. Le partage comme seule justice effective (solidarité villageoise, politiques publiques).
\end{corrige}

\coursec{Maîtrise de la langue : Les outils de liaison pour une cohérence textuelle (2 séances)}

Après avoir exploré l'univers romanesque de \emph{Le Vieux nègre et la médaille} de Ferdinand Oyono et appréhendé les enjeux thématiques de l'œuvre, nous devons désormais outiller notre expression. Une dissertation, fût-elle riche d'idées, demeure inopérante si le fil conducteur se rompt. Ces deux séances visent à transformer une suite de phrases juxtaposées en un texte fluide, logique et persuasif, en maîtrisant l'arsenal des connecteurs et des outils de reprise.

\courssub{Typologie et hiérarchie des outils de liaison}

Pour choisir l'outil adapté, il faut d'abord identifier la nature du lien logique que l'on souhaite établir (cause, conséquence, opposition, concession, énumération, temps...) et le niveau syntaxique auquel on opère (à l'intérieur de la phrase simple, entre les phrases, entre les paragraphes).

\begin{popfigure}
\begin{tikzpicture}[
    level distance=2.1cm,
    every node/.style={font=\small, align=center},
    edge from parent/.style={draw=popblue, thick, -{Stealth[length=3mm]}},
    level 1/.style={sibling distance=4.6cm},
    level 2/.style={sibling distance=2.6cm}
]
\node[draw=popblue, fill=popblueL, rounded corners, thick] {Outils de liaison}
    child {node[draw=poporange, fill=poporangeL, rounded corners] {Coordination}
        child {node[align=center]{\textbf{mais, ou, et,}\\ \textbf{donc, or, ni, car}}}
        child {node[align=center]{Éléments de\\ \textbf{même fonction}}}
    }
    child {node[draw=popgreen, fill=popgreenL, rounded corners] {Subordination}
        child {node[align=center]{\textbf{que, quand, comme,}\\ \textbf{si, puisque, afin que...}}}
        child {node[align=center]{Lien de \textbf{dépendance} :\\ proposition subordonnée}}
    }
    child {node[draw=poppurple, fill=poppurpleL, rounded corners] {Adverbes et locutions}
        child {node[align=center]{\textbf{donc, cependant,}\\ \textbf{en outre, en effet...}}}
        child {node[align=center]{\textbf{Organisation du discours} :\\ lien entre les phrases}}
    }
    child {node[draw=poppink, fill=poppinkL, rounded corners] {Pronoms relatifs}
        child {node[align=center]{\textbf{qui, que, dont,}\\ \textbf{où, lequel...}}}
        child {node[align=center]{\textbf{Reprise de l'antécédent},\\ évite la répétition}}
    };
\end{tikzpicture}
\legende{Classification fonctionnelle des principaux outils de liaison en français.}
\end{popfigure}

\begin{definition}[Connecteur logique]
Un \cle{connecteur logique} (ou connecteur textuel) est un mot ou un groupe de mots (locution) qui marque explicitement la relation sémantique (cause, conséquence, opposition, concession, ajout, illustration, conclusion...) entre deux unités de sens (propositions, phrases, paragraphes). Il guide le lecteur dans l'enchaînement de la pensée.
\end{definition}

\begin{propriete}[Polyfonctionnalité et variation des catégories]
Une même relation logique (par exemple la \textbf{cause}) peut s'exprimer par des outils de nature grammaticale différente, offrant des nuances de style et de registre :
\begin{itemize}
    \item \textbf{Conjonction de subordination :} \emph{parce que, comme, puisque, dès que} (lien standard, intégré à la phrase).
    \item \textbf{Préposition / Locution prépositive :} \emph{à cause de, en raison de, du fait de, grâce à} (suivies d'un groupe nominal, registre soutenu).
    \item \textbf{Adverbe / Locution adverbiale :} \emph{en effet, c'est pourquoi} (liaison entre les phrases ; marque l'explication ou la justification).
    \item \textbf{Participe présent / Gérondif :} \emph{la situation l'exigeant, en refusant la médaille} (concision, registre littéraire).
    \item \textbf{Nominalisation :} \emph{c'est la raison pour laquelle..., l'origine de ce phénomène réside dans...} (fortement structurant, typique de la dissertation).
\end{itemize}
\textbf{Exigence au Probatoire / Baccalauréat technique :} savoir passer de l'une à l'autre (reformulation) est une compétence évaluée en contraction de texte et en dissertation (note de langue).
\end{propriete}

\courssub{La coordination : l'égalité syntaxique}

Les conjonctions de coordination (\textbf{mais, ou, et, donc, or, ni, car} — moyen mnémotechnique : \emph{Mais Où Et Donc Or Ni Car}) lient des éléments de \textbf{même fonction} (sujet + sujet, verbe + verbe, proposition + proposition) et de \textbf{même nature}. Elles n'établissent pas de dépendance hiérarchique.

\begin{exemplebox}[Valeurs nuancées de \og mais\fg{}, \og car\fg{} et \og donc\fg]
\begin{itemize}
    \item \textbf{Mais} : opposition nette. \emph{Meka est pauvre, \textbf{mais} il garde sa dignité.}
    \item \textbf{Car} : justification, explication (cause) ; toujours en position médiane, jamais en tête de phrase à l'écrit standard. \emph{Meka refuse la médaille, \textbf{car} elle symbolise l'aliénation.}
    \item \textbf{Donc} : conséquence logique (souvent déductive). \emph{Il a servi la France, \textbf{donc} il mérite reconnaissance.}
\end{itemize}
\end{exemplebox}

\begin{attention}[Pièges de la coordination]
\begin{itemize}
    \item \textbf{Car en tête de phrase :} \emph{\og Car il était fatigué, il dormit.\fg} $\rightarrow$ \textbf{incorrect} à l'écrit standard. Utiliser \textbf{puisque}, \textbf{comme}, \textbf{parce que} (subordination) ou \textbf{en effet} (adverbe).
    \item \textbf{Accumulation de \og et\fg :} \emph{Il vint, et il vit, et il vainquit.} Style parataxique (effet de liste). En dissertation, privilégier la subordination pour hiérarchiser : \emph{Après être venu et avoir vu, il vainquit.}
    \item \textbf{Ponctuation :} devant \og et\fg{} ou \og ou\fg{} qui coordonnent deux propositions de sujets différents, la virgule est possible mais non obligatoire ; devant \og mais\fg{} et \og car\fg{}, elle est la norme.
\end{itemize}
\end{attention}

\courssub{La subordination : la dépendance hiérarchique}

La subordination permet d'\textbf{imbriquer} l'information : une proposition principale (idée forte) et une ou plusieurs propositions subordonnées (circonstances, compléments, précisions). C'est le moteur de la phrase complexe, indispensable à l'argumentation fine.

\begin{methode}[Identifier et construire une subordination]
\begin{enumerate}
    \item \textbf{Repérer le mot subordonnant :} conjonction de subordination (\emph{que, si, quand, comme, puisque, afin que, bien que}) ou pronom relatif (\emph{qui, que, dont, où}).
    \item \textbf{Identifier la fonction de la proposition subordonnée :}
    \begin{itemize}
        \item \emph{Que / Si} $\rightarrow$ complétive (sujet ou COD) : \emph{Je pense \textbf{qu'il a raison}.}
        \item \emph{Quand / Lorsque / Afin que / Bien que / Puisque} $\rightarrow$ circonstancielle (temps, but, concession, cause) : \emph{\textbf{Bien qu'il soit vieux}, il résiste.}
        \item \emph{Qui / Que / Dont / Où / Lequel} $\rightarrow$ relative (épithète de l'antécédent) : \emph{La médaille \textbf{qu'il a reçue} est en toc.}
    \end{itemize}
    \item \textbf{Vérifier la concordance des temps et les modes :} si la principale est au passé, la subordonnée se met souvent à l'imparfait ou au plus-que-parfait ; après \emph{afin que, bien que, pour que}, le \textbf{subjonctif} est obligatoire.
\end{enumerate}
\end{methode}

\begin{exemplebox}[Subordination et argumentation]
Phrases simples : \emph{Le commandant offre une médaille. Meka la refuse. Il veut garder sa liberté.}
$\rightarrow$ Phrase complexe subordonnée : \emph{\textbf{Alors que} le commandant lui offre une médaille, Meka la refuse \textbf{afin de} préserver sa liberté.}
\emph{Gain :} hiérarchisation (opposition + but), fluidité, densité informative.
\end{exemplebox}

\courssub{Adverbes et locutions adverbiales : les organisateurs du discours}

Contrairement aux conjonctions, les adverbes de liaison (\emph{donc, cependant, par conséquent, en outre, d'ailleurs, en effet, par ailleurs, ensuite, enfin}) ne lient pas syntaxiquement deux propositions \emph{à l'intérieur} d'une même phrase. Ils opèrent \textbf{entre les phrases} (ou entre les paragraphes) pour structurer le raisonnement global. Ils sont mobiles (début, milieu, fin de phrase).

\begin{center}
\begin{tabularx}{\linewidth}{|l|X|X|}
\hline
\rowcolor{popblueL}
\textbf{Relation logique} & \textbf{Registre courant} & \textbf{Registre soutenu / Dissertation} \\ \hline
\textbf{Ajout / Énumération} & et puis, aussi, en plus & \textbf{en outre, de surcroît, par ailleurs, qui plus est, de plus} \\ \hline
\textbf{Conséquence / Conclusion} & donc, du coup, c'est pour ça que & \textbf{par conséquent, dès lors, il s'ensuit que, c'est pourquoi, en définitive} \\ \hline
\textbf{Opposition / Concession} & mais, pourtant, par contre & \textbf{cependant, néanmoins, toutefois, en revanche, certes... mais} \\ \hline
\textbf{Explication / Justification} & parce que, en fait & \textbf{en effet, en réalité, c'est que, force est de constater que} \\ \hline
\textbf{Illustration / Exemple} & par exemple, comme & \textbf{notamment, à titre d'illustration, ainsi, c'est le cas de} \\ \hline
\textbf{Résumé / Conclusion} & bref, en résumé & \textbf{en somme, pour conclure, en dernière analyse, finalement} \\ \hline
\end{tabularx}
\end{center}

\begin{attention}[Erreurs fréquentes sur les adverbes de liaison]
\begin{itemize}
    \item \textbf{Confusion \og par conséquent\fg{} / \og par suite\fg{} :} \emph{par conséquent} marque une déduction logique forte ; \emph{par suite} (ou \emph{par la suite}) marque une succession chronologique ou narrative.
    \item \textbf{Abus de \og de plus\fg{} / \og aussi\fg{} :} débuter trois phrases sur cinq par \og de plus\fg{} alourdit le style. Varier : \emph{par ailleurs, d'un autre côté, qui plus est, il convient d'ajouter que...}
    \item \textbf{Ponctuation :} l'adverbe de liaison en tête de phrase est suivi d'une virgule : \emph{Cependant, il persiste.} Au milieu de la phrase, il est encadré de virgules : \emph{Il persiste, cependant, dans son erreur.}
\end{itemize}
\end{attention}

\courssub{Les pronoms relatifs : fluidité et économie syntaxique}

Le pronom relatif (\textbf{qui, que, dont, où, lequel, auquel, duquel}) introduit une \textbf{proposition subordonnée relative} qui précise un antécédent (nom ou pronom). Il évite la répétition lourde et crée la cohésion interne de la phrase.

\begin{definition}[Antécédent et fonction du pronom relatif]
L'\cle{antécédent} est le mot (généralement un nom) que le pronom relatif remplace et dont il reprend le sens. La fonction du pronom relatif \textbf{dans sa propre proposition} détermine sa forme :
\begin{itemize}
    \item \textbf{Qui} = sujet : \emph{L'homme \textbf{qui} parle.}
    \item \textbf{Que} (\emph{qu'} devant voyelle) = COD : \emph{La médaille \textbf{qu'}il reçoit.}
    \item \textbf{Dont} = complément introduit par \textbf{de} (verbe + de, adjectif + de, nom + de) : \emph{Le vieux nègre \textbf{dont} on parle. / La médaille \textbf{dont} il est fier.}
    \item \textbf{Où} = complément circonstanciel de \textbf{lieu} ou de \textbf{temps} : \emph{Le village \textbf{où} il vit. / L'année \textbf{où} il est né.}
    \item \textbf{Lequel, laquelle, lesquels, lesquelles} (et leurs formes contractées \emph{auquel, auxquels, auxquelles, duquel, desquels, desquelles}) = complément introduit par une préposition (\textbf{à, sur, avec, pour, dans...}) autre que \emph{de} seul : \emph{La médaille \textbf{à laquelle} il tient. / Le pays \textbf{dans lequel} il vit.}
\end{itemize}
\end{definition}

\begin{exoresolu}[Choix du pronom relatif : cas pratiques]
\textbf{Énoncé :} Complétez les phrases en choisissant le pronom relatif correct (\emph{qui, que, dont, où, avec lequel, pour laquelle}).

1. Voici l'auteur \_\_\_\_\_\_\_\_ a écrit ce roman.
2. C'est le livre \_\_\_\_\_\_\_\_ je t'ai parlé. (parler \textbf{de})
3. La ville \_\_\_\_\_\_\_\_ je suis né est loin.
4. Le stylo \_\_\_\_\_\_\_\_ j'écris est bleu. (écrire \textbf{avec})
5. Voici la raison \_\_\_\_\_\_\_\_ il est parti. (partir \textbf{pour} une raison)
\tcblower
\textbf{Solution détaillée :}
\begin{enumerate}
    \item \textbf{qui} (sujet de \emph{a écrit}). \emph{Voici l'auteur \textbf{qui} a écrit ce roman.}
    \item \textbf{dont} (\emph{parler de} $\rightarrow$ complément introduit par \emph{de}). \emph{C'est le livre \textbf{dont} je t'ai parlé.}
    \item \textbf{où} (complément circonstanciel de lieu). \emph{La ville \textbf{où} je suis né est loin.}
    \item \textbf{avec lequel} (préposition \emph{avec} + \emph{lequel}). \emph{Le stylo \textbf{avec lequel} j'écris est bleu.}
    \item \textbf{pour laquelle} (préposition \emph{pour} + \emph{laquelle}, accord avec \emph{raison}). \emph{Voici la raison \textbf{pour laquelle} il est parti.}
\end{enumerate}
\end{exoresolu}

\begin{attention}[Pièges classiques : \og dont\fg{}, \og où\fg{}, \og lequel\fg]
\begin{itemize}
    \item \textbf{Dont} $\leftrightarrow$ préposition \textbf{de} seule : \emph{se souvenir \textbf{de} quelque chose $\rightarrow$ \textbf{dont} je me souviens.} On n'écrit \textbf{jamais} \emph{de dont}.
    \item \textbf{Où} $\leftrightarrow$ \textbf{lieu / temps} : \emph{le jour \textbf{où}..., la maison \textbf{où}...} Mais après une préposition (\emph{dans, sur, vers}), on revient à \emph{lequel} : \emph{la maison \textbf{dans laquelle} il vit} (et non \emph{dans où}).
    \item \textbf{Lequel / auquel / duquel} $\leftrightarrow$ prépositions \textbf{à, de, pour, avec, sur, sous, dans...} (sauf \emph{de} seul, qui donne \emph{dont}) : \emph{le livre \textbf{sur lequel} je travaille}.
    \item \textbf{Accord :} \emph{lequel} s'accorde en genre et en nombre avec l'antécédent : \emph{la table sur \textbf{laquelle}}, \emph{les livres \textbf{auxquels}}.
\end{itemize}
\end{attention}

\courssub{Atelier de transformation : de la phrase hachée au paragraphe fluide}

C'est l'exercice clé pour la dissertation. Un paragraphe \og haché\fg{} utilise des phrases courtes, juxtaposées, souvent débutées par le sujet ou par \og Et / Mais / Donc\fg. Le paragraphe \og fluide\fg{} utilise la subordination, les relatifs et les adverbes de liaison pour rendre visible la logique interne du raisonnement.

\begin{methode}[Réécrire un paragraphe haché (démarche en 4 étapes)]
\begin{enumerate}
    \item \textbf{Segmenter et étiqueter :} repérer les idées unitaires (une phrase = une idée) et la relation logique entre elles (cause ? conséquence ? opposition ? illustration ? précision ?).
    \item \textbf{Hiérarchiser :} décider quelle idée est la \textbf{principale} (l'idée directrice) et lesquelles sont \textbf{secondaires} (arguments, exemples, nuances).
    \item \textbf{Choisir les outils :}
    \begin{itemize}
        \item idée secondaire = précision sur un nom $\rightarrow$ \textbf{proposition relative} (\emph{qui, que, dont, où}) ;
        \item idée secondaire = circonstance (cause, but, concession, temps) $\rightarrow$ \textbf{subordonnée conjonctive} (\emph{parce que, afin que, bien que, quand}) ;
        \item idée principale suivante = conséquence ou opposition par rapport à la précédente $\rightarrow$ \textbf{adverbe de liaison} (\emph{par conséquent, cependant, en outre}).
    \end{itemize}
    \item \textbf{Recomposer et relire :} vérifier la concordance des temps, les accords (participes, relatifs), la ponctuation et le rythme (alterner phrases longues et courtes).
\end{enumerate}
\end{methode}

\begin{exoresolu}[Transformation d'un paragraphe de dissertation « haché »]
\textbf{Sujet :} \og L'école peut-elle tout changer ?\fg

\textbf{Paragraphe haché (élève) :}
\emph{L'école transmet des savoirs. Elle forme les citoyens. C'est sa mission principale. Mais l'école ne peut pas tout changer. La famille a un rôle. La société a un rôle. L'élève a sa responsabilité. Donc l'école n'est pas la seule responsable.}
\tcblower
\textbf{Analyse des relations logiques :}
\begin{itemize}
    \item Phrases 1 + 2 + 3 : énumération, définition (mission principale).
    \item Phrase 4 : opposition, nuance (\emph{mais}).
    \item Phrases 5 + 6 + 7 : arguments complémentaires (les autres acteurs).
    \item Phrase 8 : conclusion partielle (\emph{donc}).
\end{itemize}

\textbf{Paragraphe réécrit (fluide, niveau Terminale) :}
\emph{\textbf{Si} l'école assure, par sa mission première, la transmission des savoirs et la formation du citoyen, \textbf{elle ne saurait pour autant} porter l'entière responsabilité du changement social. \textbf{En effet}, la famille, la société dans son ensemble, \textbf{ainsi que} la responsabilité personnelle de l'élève, constituent des leviers \textbf{sans lesquels} l'action scolaire reste inopérante. \textbf{Par conséquent}, l'école apparaît comme un acteur nécessaire mais non suffisant.}

\textbf{Gains observés :}
\begin{itemize}
    \item \textbf{Subordination concessive} (\emph{Si... elle ne saurait pour autant}) : remplace le \og mais\fg{} initial et nuance l'opposition.
    \item \textbf{Groupes nominaux développés} (\emph{par sa mission première, la transmission des savoirs...}) : condensent les trois premières phrases.
    \item \textbf{Adverbe \og en effet\fg} : justifie l'opposition.
    \item \textbf{Relative \og sans lesquels\fg} (pronom relatif composé après préposition) : lie élégamment les acteurs à la condition d'efficacité.
    \item \textbf{Adverbe \og par conséquent\fg} : marque la conclusion logique du paragraphe.
\end{itemize}
\end{exoresolu}

\courssub{Application directe : entraînement à la réécriture}

\begin{exercice}[Réécriture de paragraphes « hachés »]
Reprenez les deux paragraphes suivants en appliquant la méthode en 4 étapes. Visez deux à trois phrases complexes maximum par paragraphe, en variant les outils de liaison (subordination, relatifs, adverbes). Soulignez les connecteurs utilisés.

\textbf{Paragraphe A (thème : le progrès technique) :}
\emph{Le progrès technique améliore la vie. Il facilite le travail. Il soigne les maladies. Mais il crée de la pollution. Il détruit la nature. Il peut aliéner l'homme. Donc il faut le contrôler.}

\textbf{Paragraphe B (thème : la révolte dans \emph{Le Vieux nègre et la médaille}) :}
\emph{Meka est un vieux tirailleur. Il a combattu pour la France. Il reçoit une médaille. Il la refuse. Il veut montrer sa révolte. Il refuse l'assimilation. C'est un acte de dignité.}
\end{exercice}

\begin{corrige}[Exercice n°1 : propositions de réécriture]
\textbf{Paragraphe A :}
\emph{\textbf{Bien que} le progrès technique améliore indéniablement le quotidien \textbf{en facilitant} le travail et \textbf{en soignant} les maladies, il engendre \textbf{parallèlement} pollution et destruction de la nature, au point de risquer d'aliéner l'homme. \textbf{Dès lors}, un contrôle rigoureux s'impose.}

\emph{(Outils : subordonnée concessive \emph{bien que} + subjonctif ; gérondifs \emph{en facilitant / en soignant} ; adverbe \emph{parallèlement} pour l'opposition interne ; adverbe conclusif \emph{dès lors}.)}

\textbf{Paragraphe B :}
\emph{Meka, \textbf{vieux tirailleur ayant combattu} pour la France, refuse la médaille \textbf{qu'on lui décerne} \textbf{afin de} manifester sa révolte contre l'assimilation ; \textbf{c'est en cela} que réside son acte de dignité.}

\emph{(Outils : participe passé composé \emph{ayant combattu} pour condenser ; relative \emph{qu'on lui décerne} ; subordonnée de but \emph{afin de} + infinitif ; structure emphatique \emph{c'est en cela que} pour conclure.)}
\end{corrige}

\begin{aretenir}[Bilan : les 4 piliers de la cohérence textuelle]
\begin{enumerate}
    \item \textbf{Coordination} (\emph{mais, ou, et, donc, or, ni, car}) : pour l'addition simple ou l'opposition nette \textbf{à l'intérieur} de la phrase (éléments de même niveau).
    \item \textbf{Subordination} (\emph{que, si, comme, puisque, afin que, bien que...}) : pour \textbf{hiérarchiser} (cause, but, concession, condition, temps) \textbf{à l'intérieur} de la phrase.
    \item \textbf{Adverbes et locutions} (\emph{par conséquent, cependant, en outre, en effet...}) : pour \textbf{structurer le raisonnement} \textbf{entre} les phrases et les paragraphes.
    \item \textbf{Pronoms relatifs} (\emph{qui, que, dont, où, lequel...}) : pour la \textbf{cohésion nominale} (précision, évitement de la répétition).
\end{enumerate}
\textbf{Règle d'or de la dissertation :} \og Une idée directrice = un paragraphe = une architecture subordonnée (arguments et exemples imbriqués) + des adverbes de liaison pour chaîner les paragraphes.\fg
\end{aretenir}

\begin{savaistu}[Le connecteur dans l'histoire de la rhétorique]
Les Anciens (Aristote, Cicéron, Quintilien) distinguaient déjà les \emph{conjonctions} qui lient grammaticalement les mots, de l'\emph{asyndète} (absence de connecteur : style télégraphique, effet d'urgence) et de la \emph{polysyndète} (multiplication des \emph{et} : style solennel, biblique). Au XVII\textsuperscript{e} siècle, Malherbe puis Boileau imposent la \textbf{clarté logique} : chaque relation de sens doit avoir son mot propre. C'est la naissance du \og bon usage\fg{} français, où \emph{car} ne commence pas la phrase et où \emph{dont} remplace \emph{de qui}. En dissertation aujourd'hui, on attend de vous cet \og esprit géométrique\fg : ni asyndète (incohérence), ni polysyndète (lourdeur), mais \textbf{précision des connecteurs}.
\end{savaistu}

\begin{exemplebox}[Repères chiffrés : la densité des connecteurs dans une bonne copie]
Dans une copie de dissertation bien notée en expression écrite, on observe en moyenne, pour environ 300 mots :
\begin{itemize}
    \item \textbf{18 à 22 connecteurs explicites} au total ;
    \item une répartition équilibrée, par exemple : environ 30\,\% de subordinations (\emph{que, si, comme, afin que}), 25\,\% de relatifs (\emph{qui, que, dont, où}), 25\,\% d'adverbes de discours (\emph{par conséquent, cependant, en outre}) et 20\,\% de coordinations (\emph{mais, et, or}) ;
    \item \textbf{aucune occurrence} de \emph{et} ou de \emph{car} en début de phrase, et pas plus de deux \emph{de plus} dans tout le devoir.
\end{itemize}
La \textbf{variété} des connecteurs employés est un critère déterminant de la note de langue : un correcteur sanctionne la répétition du même connecteur autant que son absence.
\end{exemplebox}

\coursec{Le texte argumentatif : Stratégies et analyse (1 séance)}

Dans la section précédente, nous avons appris à relier les phrases entre elles grâce aux conjonctions, aux adverbes de liaison et aux pronoms relatifs : nous savons désormais construire un texte \cle{cohérent}. Mais un texte cohérent n'est pas forcément un texte \cle{convaincant}. Or, au Probatoire comme au Baccalauréat technique, l'épreuve de dissertation et le commentaire de texte exigent de comprendre et de produire des textes qui \emph{argumentent}. Cette séance nous donne les clés pour analyser les stratégies par lesquelles un auteur cherche à nous persuader, avant de les utiliser nous-mêmes dans nos propres devoirs.

\courssub{Qu'est-ce qu'argumenter ? Argumentation directe et argumentation indirecte}

Commençons par une intuition simple. Quand un élève veut convaincre son père de lui acheter un téléphone, il peut dire franchement : « J'ai besoin d'un téléphone pour mes recherches scolaires » : c'est une démarche directe. Mais il peut aussi raconter l'histoire de son voisin qui a raté son concours faute d'accès à Internet : il fait passer le même message par un récit, de manière indirecte. L'argumentation, c'est toujours un \cle{choix de stratégie}.

\begin{definition}[Le texte argumentatif]
Un \cle{texte argumentatif} est un texte dont le but est de \textbf{convaincre} (faire adhérer à une idée par la raison) ou de \textbf{persuader} (faire adhérer en touchant les émotions) un destinataire à propos d'une \cle{thèse}, c'est-à-dire un point de vue défendu sur une question. On distingue :
\begin{itemize}
\item l'\cle{argumentation directe} : l'auteur expose ouvertement sa thèse et la défend en son nom propre (essai, discours, éditorial, pamphlet, lettre ouverte) ;
\item l'\cle{argumentation indirecte} : l'auteur fait passer son message à travers une fiction, des personnages, une histoire (fable, conte, roman, théâtre, apologue). Le lecteur doit alors \emph{déduire} la thèse.
\end{itemize}
\end{definition}

\begin{exemplebox}[Deux formes pour une même idée]
\textbf{Argumentation directe} : « L'État camerounais doit honorer ses dettes envers les anciens combattants. » (éditorial)

\textbf{Argumentation indirecte} : dans \emph{Le Vieux nègre et la médaille}, Ferdinand Oyono raconte l'histoire de Meka, vieillard décoré pour avoir donné ses terres et ses fils à la France. À travers les déconvenues de ce personnage naïf, le romancier dénonce la colonisation sans jamais écrire « la colonisation est injuste ». C'est le lecteur qui tire lui-même la conclusion : la médaille est une récompense dérisoire.
\end{exemplebox}

\begin{savaistu}[L'ironie, arme des écrivains colonisés]
L'\cle{ironie} est l'arme majeure de Ferdinand Oyono. Elle consiste à dire le contraire de ce que l'on pense, ou à faire l'éloge de ce que l'on veut condamner, pour mieux dénoncer. Sous la colonisation, critiquer ouvertement l'administration était dangereux : l'ironie et la fiction romanesque permettaient de dénoncer tout en se protégeant de la censure coloniale. Le rire devient ainsi une forme de résistance.
\end{savaistu}

\courssub{Identifier la thèse et le point de vue de l'énonciateur}

Analyser un texte argumentatif commence toujours par une question : \emph{que veut me faire croire ou ressentir l'auteur ?}

\begin{definition}[Thèse et point de vue]
La \cle{thèse} est l'idée principale que l'auteur défend. Elle peut être :
\begin{itemize}
\item \textbf{explicite} : formulée clairement dans le texte (« Je soutiens que... », « Il faut que... ») ;
\item \textbf{implicite} : suggérée, non formulée ; c'est le cas fréquent dans l'argumentation indirecte.
\end{itemize}
Le \cle{point de vue de l'énonciateur} est la position (sociale, politique, affective) depuis laquelle il parle : un colon, un chef traditionnel et un vétéran ne diront pas la même chose de la même médaille.
\end{definition}

Dans \emph{Le Vieux nègre et la médaille}, la thèse est \textbf{implicite} : Oyono ne dit jamais « la colonisation est une imposture ». Mais le récit de la cérémonie de la médaille, présentée comme un grand honneur alors qu'elle humilie Meka, impose au lecteur cette conclusion. Le point de vue adopté est celui des colonisés : le narrateur épouse souvent le regard naïf de Meka, ce qui rend l'ironie encore plus mordante, car le lecteur, lui, comprend ce que le personnage ne comprend pas.

\begin{attention}[Piège fréquent]
Ne pas confondre la thèse de l'\textbf{auteur} avec celle d'un \textbf{personnage}. Quand le commandant dans le roman vante la « générosité » de la France, ce n'est pas Oyono qui le pense : c'est précisément ce discours qu'Oyono ridiculise. Relever qui parle est donc indispensable avant d'attribuer une opinion.
\end{attention}

\courssub{Les types d'arguments et leur organisation}

Une fois la thèse identifiée, il faut inventorier les \cle{arguments}, c'est-à-dire les raisons avancées pour la soutenir.

\begin{definition}[Les trois grands types d'arguments]
\begin{itemize}
\item L'\cle{argument logique} (ou rationnel) : il s'adresse à la raison (faits, chiffres, comparaisons, causes et conséquences). Exemple : « Trente ans après, aucune pension n'a été versée. »
\item L'\cle{argument affectif} : il s'adresse aux émotions (pitié, peur, fierté, indignation). Exemple : le récit de la nuit de Meka en prison, qui éveille la compassion du lecteur.
\item L'\cle{argument d'autorité} : il s'appuie sur une référence prestigieuse (loi, coutume, expert, texte sacré). Exemple : « Le décret de 19.. stipule que... »
\end{itemize}
\end{definition}

\begin{attention}[Argument ou exemple ?]
Ne pas confondre l'\cle{argument} (la raison générale donnée pour défendre la thèse) et l'\cle{exemple} (le cas concret qui illustre cet argument). L'exemple seul ne prouve rien : il \emph{soutient} et \emph{incarne} l'argument. Dans une copie, chaque argument doit être accompagné d'un exemple, mais un exemple sans argument est une simple anecdote.
\end{attention}

Ces arguments ne sont pas jetés au hasard : ils sont \textbf{ordonnés}. C'est là qu'intervient la notion centrale de cette séance.

\begin{definition}[Stratégie argumentative]
Une \cle{stratégie argumentative} est le choix délibéré d'un ordonnancement des arguments pour convaincre le destinataire. On distingue trois organisations principales :
\begin{itemize}
\item l'organisation \textbf{déductive} : on va du général au particulier (on pose un principe, puis on l'applique) ;
\item l'organisation \textbf{inductive} : on va du particulier au général (on accumule des faits, puis on conclut) ;
\item l'organisation \textbf{dialectique} : on confronte une thèse et une antithèse avant de proposer une synthèse (c'est le plan classique de la dissertation).
\end{itemize}
On observe aussi souvent une gradation : les arguments sont rangés \emph{du plus faible au plus fort} pour terminer par l'argument décisif.
\end{definition}

\courssub{Les procédés rhétoriques au service de l'argumentation}

Les arguments forment le squelette du texte ; les \cle{procédés rhétoriques} en sont la chair : ils donnent force, relief et émotion au propos.

\begin{center}
\begin{tabularx}{\linewidth}{|l|X|X|}
\hline
\rowcolor{popblueL} \textbf{Procédé} & \textbf{Définition} & \textbf{Exemple} \\
\hline
\cle{Ironie} & Dire le contraire de ce qu'on pense pour dénoncer & « Quel honneur immense que cette médaille en échange de deux fils morts ! » \\
\hline
\cle{Antiphrase} & Figure d'ironie : énoncé dont le sens est inverse du sens voulu & « Bravo, belle récompense ! » (pour une injustice) \\
\hline
\cle{Hyperbole} & Exagération pour frapper l'esprit & « Ils ont attendu mille ans leur pension. » \\
\hline
\cle{Énumération} & Accumulation de termes de même nature & « Ils ont donné leurs terres, leurs fils, leur jeunesse. » \\
\hline
\cle{Question rhétorique} & Question qui n'attend pas de réponse : elle affirme & « Peut-on payer une vie avec un ruban ? » \\
\hline
\cle{Parallélisme} & Répétition de la même structure syntaxique (souvent avec anaphore) & « Ils ont donné leurs terres. Ils ont donné leurs fils. Ils ont donné leur vie. » \\
\hline
\end{tabularx}
\end{center}

Chaque procédé produit un \cle{effet de sens} précis : l'ironie crée une complicité entre l'auteur et le lecteur contre la cible visée ; l'hyperbole dramatise ; la question rhétorique implique le lecteur en le forçant à répondre intérieurement ; le parallélisme donne un rythme de martèlement qui rend l'argumentation inoubliable.

\begin{exemplebox}[Barème et attentes au Bac]
Au Baccalauréat technique, l'explication d'un texte argumentatif vaut souvent 4 à 6 points. Pour obtenir la totalité des points, il ne suffit pas de repérer un procédé : il faut le \textbf{citer exactement}, le \textbf{nommer correctement} et \textbf{expliquer son effet}. « L'auteur utilise l'ironie » rapporte peu ; « L'auteur qualifie ironiquement la médaille de "grande récompense", ce qui souligne le contraste entre le sacrifice de Meka et la dérisoire contrepartie, et dénonce ainsi l'hypocrisie coloniale » rapporte le maximum.
\end{exemplebox}

\courssub{Application : analyse d'un éditorial de presse camerounaise}

Transposons ces acquis sur un texte de presse. Lisons l'extrait suivant, inspiré d'un éditorial intitulé « La dette de la nation » :

\begin{exemplebox}[Extrait étudié]
« Ils ont donné leurs terres. Ils ont donné leurs fils. Ils ont donné leur jeunesse. Soixante ans plus tard, que reste-t-il à nos vétérans ? Une médaille en chocolat et des promesses gravées dans le sable. La loi existe pourtant, votée, promulguée, solennelle. Mais une loi qui dort dans les tiroirs vaut-elle mieux qu'une loi qui n'existe pas ? Pendant que l'on décore les poitrines, on oublie les ventres. Il est temps que la nation paie sa dette. »
\end{exemplebox}

\begin{methode}[Analyser un texte argumentatif en cinq étapes]
\begin{enumerate}
\item \textbf{Situation d'énonciation} : qui parle ? À qui ? Dans quel but ? (Ici : un éditorialiste s'adresse aux lecteurs et aux pouvoirs publics pour dénoncer une injustice.)
\item \textbf{La thèse} : quelle idée est défendue ? Est-elle explicite ou implicite ?
\item \textbf{Les arguments} : quelles raisons soutiennent la thèse ? De quel type sont-elles (logique, affectif, autorité) ?
\item \textbf{Les procédés} : quelles figures renforcent les arguments, et avec quel effet ?
\item \textbf{L'organisation} : comment les arguments sont-ils ordonnés (déductif, inductif, dialectique, gradation) ?
\end{enumerate}
\end{methode}

\begin{popfigure}
\begin{tikzpicture}
\node[draw=popgreen, rectangle, rounded corners, fill=popgreenL, text width=11cm, align=left, inner sep=8pt] (grille) {\textbf{Grille d'analyse argumentative}\\[2pt]
\textbf{Sujet :} extrait éditorial « La dette de la nation »\\
\textbf{Thèse (explicite à la fin) :} l'État doit réparer l'injustice faite aux vétérans.\\
\textbf{Argument 1 (éthique) :} les promesses faites n'ont pas été tenues. \emph{Exemple :} témoignages des vétérans, « promesses gravées dans le sable ».\\
\textbf{Argument 2 (juridique / autorité) :} la loi existe mais n'est pas appliquée. \emph{Exemple :} « la loi votée, promulguée » qui « dort dans les tiroirs ».\\
\textbf{Procédés :} ironie (« médaille en chocolat »), anaphore et parallélisme (« Ils ont donné... »), question rhétorique, antithèse (« décorer les poitrines / oublier les ventres »).\\
\textbf{Ton :} indigné et solennel.\\
\textbf{Organisation :} inductive (constats $\rightarrow$ conclusion) avec gradation.};
\end{tikzpicture}
\legende{Grille d'analyse argumentative complétée pour l'éditorial « La dette de la nation ».}
\end{popfigure}

\begin{exoresolu}[Analyse guidée de l'extrait]
\textbf{Énoncé.} Relevez dans l'extrait « La dette de la nation » : (a) la thèse ; (b) deux arguments avec leur type ; (c) trois procédés rhétoriques en précisant leur effet ; (d) le type d'organisation.
\tcblower
\textbf{Solution détaillée.}
\begin{enumerate}
\item \textbf{Thèse} : elle est explicite dans la dernière phrase (« Il est temps que la nation paie sa dette ») : l'État doit réparer l'injustice faite aux anciens combattants. Le reste du texte la préparait de manière implicite.
\item \textbf{Arguments} : (1) argument \emph{éthique et affectif} : les promesses faites aux vétérans n'ont pas été tenues, ce qui provoque l'indignation ; (2) argument \emph{d'autorité} : la loi existe (référence juridique) mais reste inappliquée, ce qui rend l'injustice encore plus grave.
\item \textbf{Procédés} : l'\emph{anaphore} « Ils ont donné... » crée un rythme de martèlement qui mesure l'ampleur du sacrifice ; l'\emph{ironie} de « médaille en chocolat » ridiculise la récompense et dénonce la dérision ; la \emph{question rhétorique} « une loi qui dort... vaut-elle mieux... ? » force le lecteur à répondre « non » et l'associe au jugement de l'éditorialiste.
\item \textbf{Organisation} : inductive — l'auteur accumule les constats (sacrifices, promesses, loi inappliquée) avant de conclure —, avec une gradation qui mène au mot final « dette », le plus fort.
\end{enumerate}
\end{exoresolu}

\courssub{La grille d'analyse : un outil à réutiliser}

\begin{aretenir}[La grille d'analyse argumentative]
Pour tout texte argumentatif, remplis ce tableau à cinq entrées :
\begin{center}
\begin{tabularx}{\linewidth}{|l|X|}
\hline
\rowcolor{popblueL} \textbf{Rubrique} & \textbf{Questions à se poser} \\
\hline
\cle{Thèse} & Quelle idée est défendue ? Explicite ou implicite ? \\
\hline
\cle{Arguments} & Quelles raisons ? Logiques, affectifs, d'autorité ? \\
\hline
\cle{Exemples} & Quels cas concrets illustrent chaque argument ? \\
\hline
\cle{Procédés} & Quelles figures ? Quel effet de sens ? \\
\hline
\cle{Ton} & Sérieux, ironique, indigné, polémique, didactique ? \\
\hline
\end{tabularx}
\end{center}
Cette grille servira à la fois pour le commentaire de texte et pour préparer vos propres dissertations : ce que l'on sait analyser, on sait le produire.
\end{aretenir}

\begin{exercice}[À vous de jouer]
Choisissez un éditorial d'un journal camerounais (sur les retraites, la jeunesse, l'environnement urbain, etc.). (1) Précisez la situation d'énonciation. (2) Formulez la thèse en une phrase. (3) Relevez deux arguments et classez-les. (4) Identifiez deux procédés rhétoriques et expliquez leur effet. (5) Indiquez le type d'organisation et le ton du texte. Vous présenterez vos résultats dans la grille d'analyse.
\end{exercice}

En résumé, argumenter, c'est faire des choix : choisir une voie directe ou indirecte, choisir des arguments, les ordonner selon une stratégie, et les habiller de procédés rhétoriques. Oyono, dans \emph{Le Vieux nègre et la médaille}, a choisi la fiction et l'ironie ; l'éditorialiste choisit la dénonciation frontale. La prochaine étape sera de mettre ces stratégies au service de votre propre plume : construire, de l'introduction à la conclusion, une dissertation complète.

\coursec{La dissertation : Architecture de l'introduction à la conclusion (3 séances)}

Dans la section précédente, nous avons analysé les stratégies argumentatives du texte : comment un auteur convainc, persuade, organise ses arguments et ses exemples. Nous allons maintenant passer de la \emph{réception} à la \emph{production} : c'est à votre tour de construire un texte argumentatif complet et codifié, la \cle{dissertation}, exercice roi du Probatoire et du Baccalauréat technique. Une dissertation n'est pas un texte libre : c'est un édifice. Comme une maison, elle repose sur un plan précis : une porte d'entrée (l'introduction), des pièces organisées (le développement) et une sortie (la conclusion). Cette section vous donne les plans de cet édifice, pièce par pièce.

\courssub{L'introduction : les quatre étapes obligatoires}

\textbf{Intuition.} Imaginez que vous invitez un ami chez vous dans un quartier qu'il ne connaît pas. Vous ne lui dites pas seulement « viens » : vous lui indiquez d'abord la ville, puis le quartier, puis la rue, puis la maison. L'introduction fonctionne exactement ainsi : elle conduit le lecteur, pas à pas, du général vers le sujet précis. C'est ce qu'on appelle la méthode de l'\emph{entonnoir}.

\begin{methode}[L'introduction « entonnoir » : du général au particulier]
L'introduction se rédige en \textbf{quatre étapes obligatoires}, dans cet ordre :
\begin{enumerate}
\item \textbf{L'accroche (ou amorce)} : une phrase générale qui situe le sujet dans un contexte large (littéraire, historique, social). Pour un sujet sur \emph{Le Vieux Nègre et la médaille}, on peut partir du contexte colonial au Cameroun, du rapport entre l'écrivain et la société, ou du rôle des distinctions honorifiques dans toutes les sociétés.
\item \textbf{Le sujet posé} : on reformule le sujet avec ses propres mots, en le reliant à l'accroche par une phrase de transition. On montre ainsi qu'on a compris de quoi il s'agit.
\item \textbf{Le sujet divisé (la problématique)} : on dégage les mots-clés du sujet et on formule \textbf{une seule question} qui exprime la tension, le problème intellectuel que le sujet soulève. C'est le nœud du devoir.
\item \textbf{L'annonce du plan} : on présente les grandes parties du développement à l'aide de mots-charnières : \emph{Dans un premier temps... Ensuite... Enfin...} (ou : \emph{Nous verrons d'abord... puis... enfin...}).
\end{enumerate}
Le mouvement va donc du général (le contexte : Cameroun, Afrique, monde) vers le particulier (le sujet précis et sa question centrale).
\end{methode}

\begin{attention}[Les interdits de l'introduction]
\begin{itemize}
\item \textbf{Jamais de définition de dictionnaire} (« Selon le Petit Robert, une médaille est... ») : c'est une accroche artificielle et mal vue.
\item \textbf{Jamais de copie du sujet} tel quel : il faut le \emph{reformuler}.
\item \textbf{Une seule problématique} : aligner trois ou quatre questions à la suite montre qu'on n'a pas dégagé le vrai problème. La question doit contenir une \emph{tension} (un doute, un paradoxe, un conflit d'idées).
\item \textbf{Pas de première personne envahissante} : évitez « Je vais montrer que... » ; préférez « Nous verrons que... ».
\end{itemize}
\end{attention}

\begin{exemplebox}[Introduction rédigée — sujet : « L'écrivain est-il la conscience de sa nation ? »]
\emph{Accroche} : Depuis toujours, l'écrivain occupe une place à part dans la société : par la parole et par la plume, il dit ce que le peuple pense tout bas. \emph{Sujet posé} : Au Cameroun comme ailleurs en Afrique, des auteurs comme Ferdinand Oyono ont dénoncé dans leurs œuvres les injustices de leur époque, au point d'apparaître comme les porte-parole de leurs compatriotes. \emph{Problématique} : Mais peut-on réellement considérer l'écrivain comme la conscience de sa nation, ou son rôle se limite-t-il à distraire et à faire rêver ? \emph{Annonce du plan} : Dans un premier temps, nous montrerons que l'écrivain dénonce les maux de sa société ; ensuite, nous verrons qu'il éveille et guide les consciences ; enfin, nous nuancerons en rappelant que l'écrivain reste aussi un artiste soucieux de beauté.
\end{exemplebox}

\courssub{Le paragraphe type du développement : la structure I.D.E.E.}

\textbf{Intuition.} Un paragraphe de dissertation ressemble à une déposition devant un juge : on énonce d'abord ce qu'on affirme, on explique pourquoi, on apporte une preuve, et on conclut. Sans preuve, l'affirmation ne vaut rien ; sans affirmation, la preuve ne sert à rien.

\begin{definition}[L'idée directrice]
L'\cle{idée directrice} est la phrase unique qui résume le point principal de la sous-partie. Elle doit être \textbf{argumentative}, c'est-à-dire contenir un verbe d'action ou d'état qui engage une thèse, et non être simplement descriptive.
\begin{itemize}
\item Descriptif (mauvais) : « Meka reçoit une médaille. » (simple résumé du récit)
\item Argumentatif (bon) : « La médaille de Meka illustre l'illusion d'une gloire accordée par le colonisateur. » (une thèse à défendre)
\end{itemize}
\end{definition}

\begin{methode}[Construire un paragraphe en quatre temps : I.D.E.E.]
\begin{enumerate}
\item \textbf{I — Idée directrice} : une phrase thèse qui annonce l'argument de la sous-partie.
\item \textbf{D — Développement / argumentation} : on explique l'idée, on la justifie par le raisonnement (deux à quatre phrases liées par des connecteurs : \emph{en effet, car, c'est pourquoi, ainsi}).
\item \textbf{E — Exemple(s) précis} : une preuve concrète tirée du texte étudié (une scène du \emph{Vieux Nègre et la médaille}), de l'histoire ou de l'actualité camerounaise, citée brièvement.
\item \textbf{E — Explication du lien / transition} : on montre en quoi l'exemple confirme l'idée et on répond partiellement à la problématique, puis on annonce la sous-partie suivante.
\end{enumerate}
\end{methode}

\begin{exoresolu}[Rédiger un paragraphe I.D.E.E. sur Le Vieux Nègre et la médaille]
\textbf{Énoncé} : Rédigez un paragraphe complet (structure I.D.E.E.) pour la sous-partie : « L'illusion de la gloire ».
\tcblower
\textbf{Solution détaillée} :
\begin{itemize}
\item \textbf{I — Idée directrice} : La médaille remise à Meka symbolise d'abord une gloire trompeuse, car elle flatte l'orgueil du vieil homme sans changer sa condition.
\item \textbf{D — Développement} : En effet, cette distinction ne récompense pas un mérite réel : elle sert les intérêts de l'administration coloniale, qui cherche à donner bonne conscience à son œuvre de « civilisation ». Meka, aveuglé par l'honneur inattendu, ne voit pas qu'il est instrumentalisé. Sa joie naïve révèle ainsi le décalage entre l'apparence de la reconnaissance et la réalité de la domination.
\item \textbf{E — Exemple} : Ainsi, dans le roman de Ferdinand Oyono, Meka, ivre de fierté, parcourt le village pour annoncer la nouvelle et fait préparer un festin somptueux, alors que ses terres ont été cédées aux missionnaires et que sa hutte reste misérable.
\item \textbf{E — Explication / transition} : Cet exemple montre que la médaille n'honore pas l'homme : elle masque sa dépossession. Mais cette illusion ne dure pas : c'est ce que révèle la nuit de la cérémonie, que nous allons maintenant examiner.
\end{itemize}
\end{exoresolu}

\courssub{L'équilibre du développement : parties et sous-parties}

\textbf{Intuition.} Un édifice bancal s'effondre. De même, un développement où la première partie occupe deux pages et la seconde trois lignes est un devoir déséquilibré, donc sanctionné.

\begin{aretenir}[La structure du développement]
\begin{itemize}
\item Le développement comporte \textbf{2 ou 3 grandes parties} (I, II, éventuellement III), selon le type de plan : thèse/antithèse (2 parties) ou thématique progressif (3 parties).
\item Chaque grande partie se divise en \textbf{2 sous-parties} (A et B), chacune correspondant à un paragraphe I.D.E.E.
\item Les parties doivent avoir un \textbf{volume équivalent} : comptez vos lignes au brouillon.
\item Chaque partie est séparée par une \textbf{transition} visible (une phrase qui fait le lien), et l'on saute une ligne entre les grandes parties.
\end{itemize}
\end{aretenir}

\begin{attention}[Erreur fatale : le plan « liste de courses »]
Ne construisez jamais un plan sur les personnages ou les chapitres : « I. Meka, II. Kelara, III. Le Commandant » n'est pas un plan, c'est un sommaire. Le plan doit être \textbf{thématique}, c'est-à-dire organisé autour d'idées : « I. L'illusion de la gloire, II. Le réveil douloureux, III. La révolte silencieuse ». Chaque intitulé doit exprimer une thèse, pas un nom propre.
\end{attention}

\begin{popfigure}
\begin{center}
\begin{tikzpicture}[node distance=1.2cm]
\node[draw, rectangle, rounded corners, fill=poporangeL, text width=0.85\linewidth, align=center] (intro) {\textbf{INTRODUCTION}\\ Accroche $\rightarrow$ Sujet posé $\rightarrow$ Sujet divisé $\rightarrow$ Problématique $\rightarrow$ Annonce du plan};
\node[draw, rectangle, rounded corners, fill=popblueL, text width=0.85\linewidth, align=center, below=of intro] (dev) {\textbf{DÉVELOPPEMENT}\\ I. Partie A / B \quad II. Partie A / B \quad (III. Partie A / B)\\ Chaque sous-partie = un paragraphe I.D.E.E.};
\node[draw, rectangle, rounded corners, fill=poppinkL, text width=0.85\linewidth, align=center, below=of dev] (conc) {\textbf{CONCLUSION}\\ Bilan (réponse à la problématique) $\rightarrow$ Ouverture};
\draw[->, thick, double, popdark] (intro) -- (dev);
\draw[->, thick, double, popdark] (dev) -- (conc);
\end{tikzpicture}
\end{center}
\legende{Schéma de l'architecture d'une dissertation : trois blocs reliés, du plus encadré (introduction) au plus bref (conclusion).}
\end{popfigure}

\courssub{La conclusion : deux temps plus un}

\textbf{Intuition.} La conclusion est la dernière impression laissée au correcteur : c'est elle qui reste dans sa mémoire au moment où il attribue la note. Elle ne doit ni s'éterniser, ni s'arrêter brutalement.

\begin{methode}[Rédiger la conclusion en deux temps plus un]
\begin{enumerate}
\item \textbf{Le bilan synthétique} : on répond \emph{directement} à la problématique posée dans l'introduction, en résumant le chemin parcouru (sans répéter mot pour mot l'annonce du plan). Commencez par exemple par : « Ainsi, ... » ou « En définitive, ... ».
\item \textbf{L'ouverture} : on élargit la réflexion vers une autre œuvre (par exemple \emph{Une vie de boy} du même Oyono), un autre contexte (une autre littérature engagée) ou une actualité brûlante au Cameroun (le rôle des artistes et des intellectuels dans le débat public).
\item \textbf{Le « plus un »} : une dernière phrase courte et percutante, qui donne au devoir une tonalité personnelle et maîtrisée.
\end{enumerate}
\end{methode}

\begin{exemplebox}[Conclusion rédigée — sujet : « La médaille honore-t-elle l'homme ou l'homme honore-t-il la médaille ? »]
\emph{Bilan} : Ainsi, l'analyse du \emph{Vieux Nègre et la médaille} montre que la distinction remise à Meka ne l'honore pas : elle révèle au contraire sa soumission et précipite sa chute. C'est donc l'homme, par sa dignité et ses actes, qui donne sa valeur à la médaille, et non l'inverse. \emph{Ouverture} : Cette leçon dépasse le cadre du roman : aujourd'hui encore, au Cameroun, les titres et les décorations ne valent que par l'intégrité de ceux qui les portent. \emph{Phrase finale} : La vraie médaille d'un homme, c'est finalement l'estime qu'il inspire.
\end{exemplebox}

\begin{savaistu}[Ce que le correcteur lit en premier]
Au Baccalauréat technique, les correcteurs lisent souvent l'introduction et la conclusion \emph{avant} le développement, pour se faire une première idée du candidat. Si ces deux passages sont faibles (problématique absente, bilan bâclé), la note plafonne fréquemment autour de 8 ou 9 sur 20, même lorsque le développement est correct. Soigner l'entrée et la sortie, c'est donc soigner sa note.
\end{savaistu}

\courssub{Entraînement sur sujets types Bac Tech et gestion du temps}

Mettons toute l'architecture en application sur deux sujets proches de ceux du Baccalauréat technique.

\begin{exoresolu}[Construire problématique et plan — « L'écrivain est-il la conscience de sa nation ? »]
\textbf{Énoncé} : Pour ce sujet, dégagez les mots-clés, formulez la problématique et proposez un plan détaillé en trois parties.
\tcblower
\textbf{Solution détaillée} :
\begin{itemize}
\item \textbf{Mots-clés} : « écrivain », « conscience », « nation ». La tension porte sur le rôle social de l'écrivain : dénonciateur engagé ou simple artiste ?
\item \textbf{Problématique} : L'écrivain africain, et en particulier Ferdinand Oyono, assume-t-il pleinement le rôle de conscience de sa nation, ou sa mission est-elle d'une autre nature ?
\item \textbf{Plan détaillé} :
\begin{itemize}
\item I. L'écrivain dénonce les injustices de son époque. A. La critique de la colonisation dans \emph{Le Vieux Nègre et la médaille} (la médaille, instrument de domination). B. La dénonciation des élites et de l'hypocrisie sociale.
\item II. L'écrivain éveille et guide les consciences. A. Il donne la parole aux humiliés (Meka, figure du peuple). B. Il invite le lecteur à la lucidité et à la résistance.
\item III. Mais l'écrivain reste d'abord un artiste. A. Il cherche aussi à plaire par l'humour et l'ironie (le rire chez Oyono). B. Son œuvre vise l'universel, au-delà de sa seule nation.
\end{itemize}
\end{itemize}
\end{exoresolu}

\begin{exercice}[Sujet type Bac : « La médaille honore-t-elle l'homme ou l'homme honore-t-il la médaille ? »]
En vous appuyant sur \emph{Le Vieux Nègre et la médaille}, rédigez au brouillon : 1) l'accroche et le sujet posé ; 2) la problématique (une seule question) ; 3) le plan détaillé en deux ou trois parties thématiques avec leurs sous-parties A et B ; 4) la première phrase du bilan de la conclusion. Temps imparti : 15 minutes.
\end{exercice}

\begin{corrige}[Exercice : sujet « La médaille »]
Éléments attendus :
\begin{itemize}
\item \textbf{Accroche} : partir du rôle des distinctions honorifiques dans toutes les sociétés, puis du contexte colonial au Cameroun.
\item \textbf{Problématique possible} : La médaille remise à Meka est-elle une véritable reconnaissance de sa valeur, ou révèle-t-elle que c'est l'homme qui, par sa dignité, donne seul son sens à la distinction ?
\item \textbf{Plan thématique possible} : I. L'illusion de la gloire (A. Une récompense imméritée ; B. L'orgueil aveugle de Meka) ; II. Le réveil douloureux (A. La nuit de la cérémonie et l'humiliation ; B. La découverte de la supercherie coloniale) ; III. La dignité retrouvée (A. La révolte silencieuse de Meka ; B. La valeur de l'homme au-delà des honneurs).
\item \textbf{Amorce du bilan} : « Ainsi, le roman d'Oyono démontre que la médaille n'honore pas celui qui la reçoit : c'est l'homme, par sa conduite, qui donne ou retire sa valeur à tout honneur. »
\end{itemize}
\end{corrige}

\begin{exemplebox}[Gestion du temps au Bac : les chiffres à retenir]
L'épreuve de dissertation dure \textbf{4 heures}. La répartition idéale est la suivante :
\begin{itemize}
\item \textbf{45 minutes de brouillon} : analyse du sujet (10 min), problématique (5 min), plan détaillé et recherche des exemples (30 min) ;
\item \textbf{2 h 30 de rédaction} au propre : introduction (25 min), développement (1 h 45), conclusion (20 min) ;
\item \textbf{30 minutes de relecture} : correction de l'orthographe, des accords, de la ponctuation, vérification des connecteurs de liaison.
\end{itemize}
Un entraînement régulier chronométré — par exemple \textbf{15 minutes} pour produire introduction complète et plan détaillé au brouillon — est la clé de la réussite le jour de l'examen.
\end{exemplebox}

\begin{aretenir}[L'essentiel de la section]
\begin{itemize}
\item L'introduction suit l'entonnoir en quatre étapes : accroche, sujet posé, problématique (une seule question), annonce du plan avec mots-charnières.
\item Chaque paragraphe du développement obéit à la structure I.D.E.E. : Idée directrice, Développement, Exemple, Explication-transition.
\item Le développement compte 2 ou 3 parties équilibrées, chacune divisée en A et B ; le plan est thématique, jamais une liste de personnages.
\item La conclusion répond directement à la problématique (bilan), s'élargit (ouverture) et se termine par une phrase percutante.
\item Au Bac : 45 min de brouillon, 2 h 30 de rédaction, 30 min de relecture.
\end{itemize}
\end{aretenir}

Cette architecture acquise, la prochaine section vous conduira du plan à la rédaction intégrale : vous apprendrez à formuler une problématique originale, à élaborer seul le plan détaillé du développement, puis à rédiger et améliorer un devoir complet, avant le bilan et la remédiation qui clôtureront la leçon.

\coursec{De la problématique à la rédaction finale : Processus complet d'élaboration (4 séances)}

Dans la section précédente, nous avons appris à construire chaque pièce de la dissertation séparément : l'introduction, le paragraphe de développement (idée directrice, argument, exemple) et la conclusion. Vous savez désormais fabriquer les briques. Il reste à bâtir la maison entière, dans les conditions réelles de l'examen : un sujet inconnu, une feuille blanche, une montre qui tourne. Cette cinquième section vous conduit, pas à pas, du sujet brut à la copie corrigée et améliorée. C'est ici que tout se joue au Probatoire et au Baccalauréat technique.

\courssub{Analyser le sujet : la première victoire se gagne au brouillon}

\textbf{L'intuition d'abord.} Imaginez un mécanicien qui commence à démonter un moteur avant d'avoir écouté le client décrire la panne : il risque de réparer ce qui n'est pas cassé. De même, le candidat qui se jette sur sa plume après une lecture rapide du sujet écrit souvent une belle copie... à côté de la question. L'analyse du sujet est le diagnostic : elle conditionne tout le traitement.

\begin{definition}[Le sujet de dissertation]
Le \cle{sujet} est l'énoncé posé par l'examinateur. Il comprend trois éléments à repérer systématiquement :
\begin{itemize}
\item les \cle{mots-clés} : les termes porteurs de sens (ex. \emph{médaille}, \emph{humiliation}, \emph{reconnaissance}) ;
\item les \cle{présupposés} : ce que le sujet admet sans le dire (ex. « La médaille suffit-elle à... ? » présuppose que la médaille prétend réparer quelque chose) ;
\item l'\cle{instruction} : le verbe de consigne qui fixe le type de travail attendu.
\end{itemize}
\end{definition}

\begin{center}
\begin{tabularx}{\linewidth}{|l|X|X|}\hline
\rowcolor{popblueL} \textbf{Instruction} & \textbf{Ce qu'elle exige} & \textbf{Plan adapté} \\ \hline
\emph{Discutez} & Examiner le pour et le contre, puis trancher & Dialectique (thèse / antithèse / synthèse) \\ \hline
\emph{Êtes-vous d'accord ?} & Prendre position clairement et la justifier & Dialectique ou thématique argumenté \\ \hline
\emph{Qu'en pensez-vous ?} & Donner un avis personnel nourri de références & Thématique ou dialectique \\ \hline
\emph{Commentez} & Expliquer, analyser en profondeur un jugement & Thématique \\ \hline
\end{tabularx}
\end{center}

\begin{methode}[La méthode QQOQCP pour analyser le sujet]
Avant toute écriture, interrogez le sujet au brouillon :
\begin{enumerate}
\item \textbf{Qui ?} Quels personnages, quels groupes sociaux sont concernés ? (Meka, les Blancs, l'administration coloniale...)
\item \textbf{Quoi ?} Quel est l'objet ou le thème central ? (la médaille, la reconnaissance...)
\item \textbf{Où ?} Dans quel espace ? (Doum, le village, l'Afrique coloniale...)
\item \textbf{Quand ?} À quelle époque ? (la période coloniale, mais la question dépasse-t-elle ce cadre ?)
\item \textbf{Comment ?} Par quels moyens le thème se manifeste-t-il ? (cérémonie, discours, ironie du narrateur...)
\item \textbf{Pourquoi ?} Quel est l'enjeu, le contexte, la visée de l'auteur ?
\end{enumerate}
Puis posez la question décisive : \textbf{« Quelle tension ? »} — c'est-à-dire quel conflit d'idées le sujet met-il en scène ? Cette tension, formulée en une question, devient la \cle{problématique}.
\end{methode}

\courssub{Construire la problématique : transformer le sujet en question vivante}

\textbf{L'intuition d'abord.} Un sujet fermé appelle une réponse fermée ; une problématique ouvre un débat. Dire « la médaille » est un thème ; demander « la médaille peut-elle réparer l'humiliation ? » est une problématique : on sent deux forces s'affronter, la reconnaissance officielle d'un côté, la blessure de la dignité de l'autre.

\begin{definition}[La problématique]
La \cle{problématique} est la question centrale, née d'une \cle{tension} entre deux pôles opposés, à laquelle tout le devoir répondra. Elle se formule en une phrase interrogative et doit être annoncée à la fin de l'introduction.
\end{definition}

\begin{exemplebox}[Du sujet à la problématique]
\textbf{Sujet :} « La médaille » dans \emph{Le Vieux Nègre et la médaille} de Ferdinand Oyono.\par
\textbf{Tension repérée :} honneur officiel $\leftrightarrow$ humiliation réelle.\par
\textbf{Problématique :} \emph{La reconnaissance officielle, symbolisée par la médaille, peut-elle réparer l'humiliation vécue par Meka et les siens ?}\par
Autre exemple : \textbf{Sujet :} « Le rire chez Oyono » $\rightarrow$ \textbf{Problématique :} \emph{L'ironie d'Oyono est-elle une simple arme de dérision ou un instrument de dénonciation de l'injustice coloniale ?}
\end{exemplebox}

\begin{attention}[Le piège du « hors-sujet déguisé »]
Le hors-sujet n'est pas toujours flagrant. Le plus dangereux est le \cle{hors-sujet déguisé} : traiter le thème général (la colonisation en Afrique) au lieu du sujet précis (la symbolique de la médaille chez Oyono). \textbf{Test de sécurité :} après avoir rédigé chaque idée directrice, demandez-vous : « Cette phrase répond-elle à ma problématique ? » Si la réponse est non, supprimez ou reformulez.
\end{attention}

\courssub{Élaborer le plan détaillé au brouillon : l'architecture avant les murs}

\textbf{L'intuition d'abord.} Aucun maçon sérieux ne monte un mur sans plan. Le plan détaillé est votre plan de construction : il se fait entièrement au brouillon, avant la rédaction, et il vous libère l'esprit pour soigner l'expression ensuite.

\begin{propriete}[Le poids du plan détaillé]
Un bon \cle{plan détaillé} au brouillon représente environ 50\,\% de la réussite du devoir. Il doit tenir sur une page A4, être structuré (I. A. B. / II. A. B.), et contenir les citations exactes avec leur référence (numéro de page et édition). Une citation approximative ou inventée discrédite tout le paragraphe.
\end{propriete}

Pour \textbf{chaque sous-partie} (A et B de chaque grande partie), notez au brouillon cinq éléments :
\begin{enumerate}
\item l'\cle{idée directrice} : une phrase complète qui répond à la problématique ;
\item un \cle{argument} précis qui la justifie ;
\item un \cle{exemple textuel} : une citation courte d'Oyono, avec la page ;
\item un \cle{exemple hors texte} : histoire du Cameroun, actualité, fait de société ;
\item un \cle{connecteur de transition} vers la sous-partie suivante.
\end{enumerate}

\begin{exemplebox}[Extrait de plan détaillé (sujet : la médaille)]
\textbf{I. La médaille comme instrument de domination déguisée}\par
\textbf{A.} \emph{Idée directrice :} la médaille flatte Meka pour mieux le maintenir dans la soumission. \emph{Argument :} l'honneur accordé ne change rien à sa condition réelle. \emph{Exemple textuel :} la cérémonie de remise, où Meka, écrasé par le soleil et la fatigue, sert de décor au discours du commandant. \emph{Exemple hors texte :} les décorations coloniales distribuées aux chefs traditionnels au Cameroun pour asseoir l'autorité administrative. \emph{Connecteur :} « Mais cet honneur apparent cache une réalité plus cruelle... »\par
\textbf{B.} \emph{Idée directrice :} la médaille ne répare ni la pauvreté ni l'humiliation subies. \emph{Argument :} après la fête, Meka reste démuni, et la pluie détruit même le symbole. \emph{Exemple textuel :} la médaille perdue dans la boue, la nuit de la tornade. \emph{Exemple hors texte :} les distinctions honorifiques actuelles qui ne s'accompagnent d'aucune amélioration concrète de la vie du récipiendaire. \emph{Connecteur :} « Ainsi, loin d'être une réparation, la médaille apparaît comme... »
\end{exemplebox}

\begin{popfigure}
\begin{center}
\begin{tikzpicture}[node distance=1.1cm and 0.8cm, every node/.style={font=\small, align=center}]
\node[draw, ellipse, fill=popgoldL, thick, align=center] (sujet){Sujet\\ Bac};
\node[draw, rectangle, rounded corners, fill=popblueL, right=of sujet, align=center] (prob){Problématique\\ (question-tension)};
\node[draw, rectangle, rounded corners, fill=popgreenL, right=of prob, align=center] (plan){Plan détaillé\\ I.A/B, II.A/B};
\node[draw, rectangle, rounded corners, fill=poporangeL, right=of plan, align=center] (redac){Rédaction\\ Intro, Dev, Conc};
\node[draw, rectangle, rounded corners, fill=poppinkL, right=of redac, align=center] (relect){Relecture\\ (grille 10 pts)};
\draw[->, very thick, popdark] (sujet) -- (prob);
\draw[->, very thick, popdark] (prob) -- (plan);
\draw[->, very thick, popdark] (plan) -- (redac);
\draw[->, very thick, popdark] (redac) -- (relect);
\end{tikzpicture}
\end{center}
\legende{La chaîne de production complète de la dissertation : cinq étapes obligatoires, du sujet à la relecture. Aucune flèche ne peut être sautée.}
\end{popfigure}

\courssub{La rédaction suivie : la séance « bac blanc » (2 heures)}

\textbf{L'intuition d'abord.} Un athlète ne court pas son premier 100 mètres le jour de la finale : il s'entraîne dans les conditions exactes de la compétition. Le bac blanc est votre entraînement chronométré.

\begin{methode}[Gérer les 2 heures de l'épreuve]
\begin{enumerate}
\item \textbf{0 h 00 – 0 h 20 :} analyse du sujet (QQOQCP), repérage de la tension, formulation de la problématique.
\item \textbf{0 h 20 – 0 h 45 :} plan détaillé complet au brouillon (idées directrices, arguments, citations, connecteurs).
\item \textbf{0 h 45 – 1 h 45 :} rédaction au propre : introduction (15 min), développement (40 min), conclusion (5 min).
\item \textbf{1 h 45 – 2 h 00 :} relecture complète avec la grille (voir 5.5).
\end{enumerate}
\textbf{Présentation :} écriture lisible et soignée, un alinéa pour l'introduction, un pour chaque sous-partie, un pour la conclusion ; sauter une ligne entre les grandes parties ; ne pas raturer massivement (un trait fin suffit).
\end{methode}

\begin{exoresolu}[Sujet de bac blanc et début de traitement]
\textbf{Énoncé.} Sujet : « La médaille de Meka : un honneur ou une insulte ? Qu'en pensez-vous ? » Vous rédigerez une dissertation complète en deux heures.\par
\tcblower
\textbf{Solution détaillée (étapes 1 et 2).}\par
\emph{Analyse QQOQCP :} Qui ? Meka, vieillard de Doum. Quoi ? la médaille, symbole de reconnaissance. Où/Quand ? village africain sous administration coloniale. Comment ? par une cérémonie officielle pleine d'ironie. Pourquoi ? Oyono dénonce l'hypocrisie du système colonial.\par
\emph{Tension :} honneur affiché $\leftrightarrow$ mépris réel.\par
\emph{Problématique :} la médaille décernée à Meka constitue-t-elle une véritable reconnaissance de sa dignité, ou n'est-elle qu'un instrument supplémentaire de son humiliation ?\par
\emph{Plan (dialectique) :} I. Une reconnaissance en apparence généreuse (A. un honneur officiel rare pour un colonisé ; B. une fierté sincèrement ressentie par Meka et le village). II. Une insulte déguisée (A. un honneur qui ne change aucune condition de vie ; B. un symbole détruit par la réalité : la médaille perdue dans la boue). Synthèse : la médaille révèle, par son échec, la profondeur de l'injustice coloniale.\par
\emph{Conclusion attendue :} réponse nuancée — la médaille est un honneur retourné en insulte, et c'est précisément ce retournement qui fait la force critique du roman.
\end{exoresolu}

\begin{exemplebox}[Objectif « zéro faute » : le compte exact]
Sur les 10 points d'expression, chaque faute de langue coûte entre 0,25 et 0,5 point. Visez : \textbf{0 faute d'accord majeur} (sujet-verbe, participe passé), \textbf{0 solécisme} (tournure incorrecte), \textbf{0 anglicisme}. Cinq fautes d'accord peuvent vous priver de 2,5 points : c'est la différence entre la moyenne et l'échec.
\end{exemplebox}

\courssub{Réécriture et amélioration : l'auto-correction guidée}

\textbf{L'intuition d'abord.} Un tailleur repasse sa couture avant de livrer le vêtement. La relecture n'est pas un luxe : c'est la dernière étape de fabrication du devoir, celle qui transforme une copie moyenne en bonne copie.

\begin{methode}[La grille de relecture en six questions]
Relisez votre copie en vous posant, dans l'ordre, ces six questions :
\begin{enumerate}
\item \textbf{Hors-sujet ?} Chaque idée directrice répond-elle à la problématique ?
\item \textbf{Plan respecté ?} La copie suit-elle le plan du brouillon (I.A/B, II.A/B) ?
\item \textbf{Connecteurs ?} Y a-t-il des outils de liaison (\emph{d'abord, en effet, cependant, ainsi, par conséquent}) entre et dans les paragraphes ?
\item \textbf{Accords ?} Sujet-verbe, adjectifs, participes passés vérifiés un à un ?
\item \textbf{Ponctuation ?} Phrases trop longues ? Virgules et points bien placés ?
\item \textbf{Registre soutenu ?} Pas de familiarité, pas de « je pense que » répétitif, pas d'abréviations ?
\end{enumerate}
\end{methode}

\begin{attention}[Les trois fautes qui coûtent le plus cher]
(1) L'accord sujet-verbe quand le sujet est éloigné du verbe ; (2) le participe passé avec \emph{avoir} accordé à tort ; (3) les homophones (\emph{a/à}, \emph{ou/où}, \emph{ce/se}, \emph{la/l'a}). Relisez une dernière fois en ne traquant que ces trois cibles.
\end{attention}

\courssub{Les échanges pairs : la lecture croisée avec la grille officielle}

Après l'auto-correction, échangez votre copie avec celle d'un camarade. Le regard extérieur voit ce que le vôtre ne voit plus. Utilisez la grille d'évaluation officielle, sur 20 points :

\begin{center}
\begin{tabularx}{\linewidth}{|l|X|c|}\hline
\rowcolor{popblueL} \textbf{Critère} & \textbf{Ce qui est évalué} & \textbf{Points} \\ \hline
Contenu & Pertinence par rapport au sujet, qualité de la problématique, solidité des arguments, exactitude des exemples et citations & 10 \\ \hline
Expression & Correction de la langue (orthographe, accords, syntaxe), richesse du vocabulaire, clarté, présentation matérielle & 10 \\ \hline
\end{tabularx}
\end{center}

\begin{methode}[Conduire la lecture croisée]
\begin{enumerate}
\item Lisez la copie de votre pair en entier, sans annoter, pour juger l'impression globale.
\item Relisez en attribuant les points du contenu : la problématique est-elle visible ? Chaque paragraphe y répond-il ?
\item Relisez en attribuant les points d'expression : comptez les fautes, soulignez les lourdeurs.
\item Rédigez un commentaire constructif : deux points forts, deux axes d'amélioration.
\item Récupérez votre copie, confrontez les remarques, puis \textbf{réécrivez} le passage le plus faible.
\end{enumerate}
\end{methode}

\begin{exercice}[Entraînement complet]
Sujet : « Le rire est l'arme des faibles face aux puissants. Cette affirmation s'applique-t-elle au \emph{Vieux Nègre et la médaille} ? Discutez. » (1) Formulez la problématique. (2) Rédigez le plan détaillé complet au brouillon, avec pour chaque sous-partie une idée directrice, un argument, une citation d'Oyono et un exemple hors texte. (3) Rédigez l'introduction et le premier paragraphe du développement.
\end{exercice}

\begin{corrige}[Exercice]
(1) \emph{Problématique possible :} l'ironie et le rire chez Oyono permettent-ils aux colonisés de renverser le rapport de force avec le colonisateur ?\par
(2) \emph{Plan :} I. Le rire comme arme de résistance (A. l'ironie du narrateur démasque le ridicule des administrateurs — exemple textuel : le portrait satirique du commandant lors de la cérémonie ; exemple hors texte : la tradition camerounaise des contes moqueurs ; B. le rire collectif du village soude la communauté humiliée). II. Les limites de cette arme (A. le rire ne change pas la condition matérielle de Meka ; B. la médaille perdue montre que la dérision seule ne suffit pas à réparer l'injustice).\par
(3) \emph{Introduction attendue :} amorce sur le rire comme arme sociale, présentation de l'œuvre et de l'auteur, annonce du sujet, problématique, annonce du plan en deux temps.
\end{corrige}

\begin{savaistu}[Les sujets de dissertation au Cameroun]
Au Probatoire et au Baccalauréat technique camerounais, les sujets de dissertation croisent souvent des thèmes sociétaux — la jeunesse, l'emploi, la technologie, la culture, l'éducation — avec les œuvres littéraires au programme. S'entraîner sur \emph{Le Vieux Nègre et la médaille}, c'est donc aussi s'entraîner à penser les problèmes de la société camerounaise d'aujourd'hui : la reconnaissance, la dignité, la justice.
\end{savaistu}

\begin{aretenir}[L'essentiel de la section]
\begin{itemize}
\item Analysez le sujet avec QQOQCP, repérez les mots-clés, les présupposés et l'instruction.
\item Transformez le sujet en problématique : une question née d'une tension entre deux pôles.
\item Construisez un plan détaillé au brouillon : idée directrice, argument, citation, exemple hors texte, connecteur — c'est 50\,\% de la réussite.
\item Gérez les 2 heures : 20 min d'analyse, 25 min de plan, 60 min de rédaction, 15 min de relecture.
\item Relisez avec la grille en six questions, puis faites relire par un pair avec la grille officielle (contenu 10 pts / expression 10 pts).
\end{itemize}
\end{aretenir}

\coursec{Bilan, Compte-rendu et Remédiation individualisée (1 séance)}

Après quatre séances intenses consacrées à l'élaboration complète d'une dissertation — de la formulation de la problématique à la rédaction finale en passant par la révision —, cette séance de clôture joue un rôle charnière. Elle ne saurait se réduire à une simple correction de devoirs ; elle constitue le moment méta-cognitif où l'élève, futur candidat au Probatoire puis au Baccalauréat technique, prend la mesure exacte de ses acquis et de ses lacunes pour orienter sa révision autonome. Le passage de l'apprentissage guidé (séances 1 à 5) à l'autonomie examinée (épreuve certificative) se joue ici, dans la lucidité partagée entre l'enseignant et la classe.

\courssub{Compte-rendu synthétique de l'unité : les quatre piliers consolidés}

Avant toute remédiation, il est impératif de dresser un bilan global. Cette unité a reposé sur quatre piliers indissociables, chacun ayant fait l'objet d'un apprentissage progressif.

\begin{aretenir}[Les quatre piliers de l'unité « Dissertation »]
\textbf{1. L'œuvre ancrage : \emph{Le Vieux Nègre et la Médaille} de Ferdinand Oyono.}
Maîtrise de l'intrigue, des personnages (Meka, Kelara, le Commandant, le Prêtre), des thèmes majeurs (colonisation, aliénation, dignité, résistance silencieuse, choc des cultures) et des procédés stylistiques (ironie, perspective narrative, symbolisme de la médaille). L'élève doit pouvoir mobiliser des citations précises comme arguments d'autorité ou exemples illustratifs.

\textbf{2. La langue outillée : Les connecteurs logiques.}
Maîtrise des outils de liaison (conjonctions de coordination/subordination, adverbes de liaison, locutions adverbiales, pronoms relatifs) pour assurer la \cle{cohérence} (liens logiques entre idées) et la \cle{cohésion} (liens grammaticaux/lexicaux entre phrases) du texte dissertationnaire.

\textbf{3. L'argumentation structurée : Stratégies du texte argumentatif.}
Identification des thèses, arguments, exemples, concessions, réfutations. Distinction entre argument d'autorité, argument logique, argument pathétique. Capacité à construire un \cle{raisonnement dialectique} (thèse/antithèse/synthèse) ou \cle{analytique} (causes/conséquences) selon le sujet.

\textbf{4. La méthode dissertationnaire : Architecture standardisée.}
\textbf{Introduction :} Accroche $\rightarrow$ Présentation sujet/œuvre $\rightarrow$ Problématique $\rightarrow$ Annonce du plan.
\textbf{Développement :} 2 ou 3 parties équilibrées, paragraphes types (Idée directrice + Argument + Exemple/Analyse + Transition).
\textbf{Conclusion :} Bilan synthétique $\rightarrow$ Réponse à la problématique $\rightarrow$ Ouverture pertinente.
\end{aretenir}

Pour objectiver ce bilan, l'enseignant présente à la classe une analyse statistique simulée des résultats obtenus lors du « Devoir complet » de la séance 5. Ce miroir chiffré permet de cibler collectivement les zones de fragilité.

\begin{popfigure}
\begin{tikzpicture}
\begin{axis}[
    ybar,
    enlargelimits=0.15,
    ylabel={Nombre d'élèves (sur 30)},
    xlabel={Compétences évaluées},
    symbolic x coords={Œuvre, Langue, Argumentation, Dissertation},
    xtick=data,
    nodes near coords,
    ymin=0,
    ymax=30,
    bar width=1.2cm,
    tick label style={font=\small},
    label style={font=\small},
    title style={font=\normalsize\bfseries},
    title={Répartition des réussites (note $\ge$ 10/20) par compétence -- Simulation classe}
]
\addplot[fill=popblue!70, draw=popdark] coordinates {(Œuvre, 22) (Langue, 18) (Argumentation, 20) (Dissertation, 15)};
\end{axis}
\end{tikzpicture}
\legende{Graphique simulé de résultats de classe pour cibler la remédiation. La compétence « Dissertation » (structure globale, gestion du temps) apparaît comme le point le plus critique.}
\end{popfigure}

\courssub{Remise des copies corrigées : Grille d'évaluation et commentaires personnalisés}

La remise des copies est un acte pédagogique fort. Elle ne doit pas se faire « en vrac ». Chaque élève reçoit sa copie annotée \textbf{ET} une grille d'évaluation détaillée, dédoublée selon les deux grands critères du Baccalauréat : \textbf{Contenu (10 pts)} et \textbf{Expression (10 pts)}.

\begin{exemplebox}[Extrait type de Grille d'évaluation Bac -- Dissertation]
\begin{center}
\begin{tabularx}{\linewidth}{|>{\bfseries}l|X|X|}\hline
\textbf{Critères CONTENU (10 pts)} & \textbf{Indicateurs de réussite} & \textbf{Note} \\ \hline
\bfseries Compréhension du sujet / Problématique & Sujet analysé, termes définis, tension identifiée, problématique explicite et pertinente. & /2 \\ \hline
\bfseries Connaissance de l'œuvre / Exemples & Citations précises, références justes (personnages, scènes, thèmes), exemples variés (œuvre + culture générale). & /3 \\ \hline
\bfseries Qualité de l'argumentation & Arguments logiques, variés, pertinents, articulés (connecteurs), concessions/réfutations maîtrisées. & /3 \\ \hline
\bfseries Structure / Plan & Plan annoncé respecté, équilibre des parties, transitions fluides, introduction/conclusion conformes. & /2 \\ \hline\hline
\textbf{Critères EXPRESSION (10 pts)} & \textbf{Indicateurs de réussite} & \textbf{Note} \\ \hline
\bfseries Correction linguistique & Orthographe, grammaire, syntaxe, ponctuation (peu d'erreurs ne nuisant pas au sens). & /4 \\ \hline
\bfseries Qualité de la langue / Style & Vocabulaire précis/riche, phrases bien construites, connecteurs variés/adaptés, ton soutenu. & /3 \\ \hline
\bfseries Présentation / Lisibilité & Écriture soignée, paragraphes distincts, marges respectées, ratures minimes. & /1 \\ \hline
\bfseries Gestion du temps / Complétude & Devoir complet (Intro + Développement + Conclusion) dans le temps imparti. & /2 \\ \hline
\end{tabularx}
\end{center}
\legende{Grille simplifiée pour l'élève. L'enseignant surligne les items acquis/non acquis et ajoute un commentaire global manuscrit : « Point fort : ... / Point de vigilance : ... / Objectif prioritaire pour la révision : ... »}
\end{exemplebox}

\begin{attention}[Le piège de l'incomplétude]
Ne pas négliger la \textbf{gestion du temps} en remédiation. Beaucoup d'élèves techniques échouent par \textbf{incomplétude} (absence de conclusion, partie III bâclée ou inexistante) plutôt que par incompétence pure. La case « Gestion du temps / Complétude » (2 pts faciles) fait souvent basculer la note au-dessus de 10/20. Entraînez-vous à écrire une conclusion en 5 minutes chrono.
\end{attention}

\courssub{Ateliers de remédiation ciblés : Groupes de besoins (40 min)}

Fort du diagnostic (grille + graphique), la classe est répartie en quatre ateliers tournants (10 min par atelier) ou en groupes fixes selon l'effectif. L'enseignant circule ; les élèves « experts » sur un item peuvent tutorer.

\begin{methode}[Remédiation efficace : protocole en 4 temps]
1. \textbf{Diagnostic précis (grille items).} L'élève identifie \emph{lui-même} son item le plus faible (ex: « Problématique floue » ou « Connecteurs répétitifs »).
2. \textbf{Exercice ciblé court (15 min max).} Un micro-exercice unique, typé Bac, portant uniquement sur cet item.
3. \textbf{Correction collective explicite.} Mise en commun au tableau : pourquoi cette réponse ? quel connecteur ici ? quelle reformulation ?
4. \textbf{Nouvel essai immédiat (transfert).} Réécriture d'un paragraphe de sa propre copie corrigée en appliquant la règle travaillée.
\end{methode}

\begin{center}
\begin{tabularx}{\linewidth}{|>{\bfseries}l|X|X|}\hline
\textbf{Groupe} & \textbf{Profil / Item faible (Grille)} & \textbf{Activité de remédiation (15 min)} \\ \hline
\textbf{Groupe A} & \textbf{Problématique / Plan} (Items C1, C4) & \textit{Sujet :} « La médaille honore-t-elle vraiment Meka ? » \\
& (Problématique absente, plan binaire simple, pas d'annonce) & \textbf{Tâche :} 1. Écrire 3 problématiques possibles. 2. Choisir la meilleure. 3. Construire le plan détaillé (I, A, B / II, A, B / III, A, B) avec idées directrices. \\ \hline
\textbf{Groupe B} & \textbf{Connecteurs / Syntaxe} (Items E1, E2) & \textit{Support :} Paragraphe « troué » (sans connecteurs) sur le thème de l'aliénation. \\
& (Phrases juxtaposées, connecteurs erronés, répétitions « et/mais ») & \textbf{Tâche :} 1. Insérer 5 connecteurs imposés (\emph{cependant, par conséquent, en effet, de surcroît, en définitive}). 2. Réécrire 2 phrases mal construites (relatives, subordonnées). \\ \hline
\textbf{Groupe C} & \textbf{Exemples / Connaissance œuvre} (Item C2) & \textit{Support :} Tableau à double entrée \emph{Thème / Scène précise / Citation courte}. \\
& (Exemples vagues : « dans le livre », « Meka est triste », pas de citations) & \textbf{Tâche :} Compléter 5 lignes du tableau (ex: \emph{Humiliation / Scène remise médaille / « Il tendit la main... »}). Apprendre 3 citations « passe-partout » par cœur. \\ \hline
\textbf{Groupe D} & \textbf{Gestion temps / Présentation} (Items E3, E4) & \textit{Support :} Chronomètre + Sujet type Bac. \\
& (Copie sale, pas de conclusion, temps dépassé) & \textbf{Tâche :} \textbf{Entraînement « Flash » 20 min.} Rédiger \textbf{uniquement} Introduction + Plan détaillé + Conclusion. Focus : lisibilité, saut de lignes, respect du timing (5 min Intro / 10 min Plan / 5 min Concl). \\ \hline
\end{tabularx}
\end{center}

\courssub{Exercice de consolidation : « Sujet Surprise » (1h) -- Validation de l'autonomie}

Pour clore la séance et valider le transfert des acquis, un sujet inédit est distribué. L'instruction est volontairement réduite pour tester l'autonomie méthodologique : \textbf{« Rédigez l'Introduction complète et le Plan détaillé (avec idées directrices et exemples prévus) uniquement. Temps : 1 heure. »}

\begin{exoresolu}[Sujet Surprise -- Simulation]
\textbf{Sujet :} \emph{« Dans \emph{Le Vieux Nègre et la Médaille}, la cérémonie de remise de la médaille apparaît comme le point de bascule du roman. Analysez cette scène et montrez en quoi elle révèle la tragédie de Meka. »}

\textbf{Production attendue (Introduction + Plan) :}

\textbf{INTRODUCTION}
\textbf{[Accroche]} La remise d'une distinction honorifique est, dans l'imaginaire commun, un moment de gloire et de reconnaissance sociale.
\textbf{[Présentation]} C'est ce moment que met en scène Ferdinand Oyono dans \emph{Le Vieux Nègre et la Médaille} (1956), roman phare de la littérature africaine d'expression française. Meka, vieux chef de village, reçoit la médaille de la Légion d'honneur des mains du Commandant.
\textbf{[Problématique]} Pourtant, loin de consacrer une réussite, cette cérémonie ne révèle-t-elle pas, par un cruel retournement ironique, la vacuité de la collaboration et l'aliénation définitive d'un homme que le colonisateur croit honorer mais qu'il humilie en réalité ?
\textbf{[Annonce du plan]} Nous montrerons d'abord que la cérémonie est une mise en scène coloniale grotesque (I), puis qu'elle expose au grand jour la tragédie intime de Meka, pris entre deux mondes (II), pour enfin démontrer qu'elle scelle l'échec de la « politique des cadres » et ouvre sur la prise de conscience lucide de Kelara (III).

\textbf{PLAN DÉTAILLÉ}
\textbf{I. La cérémonie : une mascarade coloniale orchestrée}
    A. Le décor et le protocole : singeage de la République (ex: discours du Commandant, « Monsieur Meka », uniforme).
    B. L'ironie narrative : le décalage entre la solennité attendue et la réalité burlesque (ex: ivresse de Meka, indifférence des notables).
    C. La médaille comme symbole vide : objet clinquant sans pouvoir réel (citation : « une petite chose brillante »).

\textbf{II. Meka : la figure tragique du collaborateur aliéné}
    A. L'attente fébrile et l'investissement identitaire (ex: achat de l'habit, fierté naïve).
    B. Le choc de la réception : le corps humilié (ex: la main tendue, le médaillon qui ne tient pas).
    C. La solitude du héros : incompréhension de Kelara, distance du Prêtre (citation : « Il avait l'air d'un enfant »).

\textbf{III. Portée symbolique : l'effondrement d'un monde et l'aube d'une conscience}
    A. L'échec de la médiation : Meka, trait d'union brisé entre tradition et modernité imposée.
    B. Le regard lucide de Kelara : refus de l'illusion, conscience de la servitude (ex: silence final, regard sur la médaille).
    C. Ouverture : la médaille comme « médaille de la trahison » annonçant les luttes pour l'indépendance.
\tcblower
\textbf{Analyse corrective (Auto-évaluation élève) :}
\begin{itemize}
    \item \textbf{Problématique :} Contient les termes du sujet, identifie la tension (honneur/humiliation), annonce une thèse (ironie/aliénation). \checkmark
    \item \textbf{Plan :} Dialectique (Apparence / Réalité / Portée). Équilibré. Idées directrices claires. Exemples précis (scènes, citations). \checkmark
    \item \textbf{Connecteurs :} « Pourtant », « loin de », « pour enfin démontrer », « puis », « pour enfin ». Variés. \checkmark
\end{itemize}
\end{exoresolu}

\courssub{Conseils pour la révision Bac : Les 4 Fiches Mémo « Prêtes à l'emploi »}

L'élève technique doit optimiser son temps de révision. Quatre fiches synthétiques (format A5, recto-verso) sont distribuées ou construites collectivement en fin de séance. Elles constituent le « kit de survie » pour l'épreuve de dissertation.

\begin{definition}[Fiche Mémo 1 : Structure de l'Introduction (Modèle universel)]
\textbf{ÉTAPE 1 -- ACCROCHE (1 phrase) :} Citation / Proverbe / Fait de société / Définition / Question rhétorique. \textit{(Lien avec le thème général du sujet).}
\textbf{ÉTAPE 2 -- PRÉSENTATION (2-3 phrases) :} Auteur + Titre (italique) + Date + Genre + Contexte bref + Résumé très court (1 ligne) centré sur le sujet.
\textbf{ÉTAPE 3 -- PROBLÉMATIQUE (1 phrase, point d'interrogation) :} Tension entre 2 termes du sujet. Contient la thèse implicite. \textit{Test : « Est-ce que ma réponse à cette question fera mon devoir ? »}
\textbf{ÉTAPE 4 -- ANNONCE DU PLAN (1 phrase) :} « Nous analyserons d'abord... / puis... / pour enfin montrer... » (Mots-clés des parties I, II, III).
\end{definition}

\begin{definition}[Fiche Mémo 2 : Connecteurs utiles par fonction logique]
\begin{center}
\begin{tabularx}{\linewidth}{|l|X|}\hline
\textbf{Fonction} & \textbf{Connecteurs (variez !)} \\ \hline
\bfseries Ajout / Enchaînement & De plus, par ailleurs, en outre, de surcroît, également, ensuite, puis. \\ \hline
\bfseries Opposition / Concession & Cependant, pourtant, néanmoins, toutefois, au contraire, certès... mais, si... cependant. \\ \hline
\bfseries Cause / Conséquence & Car, parce que, puisque, comme, en effet, donc, par conséquent, c'est pourquoi, par suite. \\ \hline
\bfseries Illustration / Exemple & Par exemple, notamment, c'est le cas de, ainsi, en particulier, tel que. \\ \hline
\bfseries Explication / Reformulation & C'est-à-dire, en d'autres termes, autrement dit, plus précisément. \\ \hline
\bfseries Conclusion / Synthèse & En définitive, en conclusion, pour conclure, au total, bref, en somme. \\ \hline
\end{tabularx}
\end{center}
\end{definition}

\begin{definition}[Fiche Mémo 3 : Citations clés \emph{Le Vieux Nègre et la Médaille} (à connaître par cœur -- 3 minimum)]
\begin{enumerate}
    \item \textbf{Sur l'aliénation / la médaille :} « Il avait l'air d'un enfant qu'on aurait habillé pour une fête. » (Incipit / Réception médaille)
    \item \textbf{Sur le choc des cultures / l'ironie :} « La médaille brillait sur sa poitrine comme un œil de chat dans la nuit. » (Symbolisme)
    \item \textbf{Sur la lucidité de Kelara / résistance :} « Elle ne dit rien. Elle regarda la médaille. Elle regarda son mari. » (Final, silence éloquent)
    \item \textbf{Sur le discours colonial :} « La France est généreuse... » (Discours Commandant, ironie dramatique).
    \item \textbf{Sur la condition de Meka :} « Meka n'était plus un chef, il était un vieux nègre décoré. » (Analyse critique possible).
\end{enumerate}
\end{definition}

\begin{definition}[Fiche Mémo 4 : Check-list Relecture (5 min avant de rendre la copie)]
\begin{enumerate}
    \item \textbf{Structure :} Intro (4 étapes) ? Plan annoncé respecté (I, II, III) ? Transitions ? Conclusion (Bilan + Réponse + Ouverture) ?
    \item \textbf{Paragraphes :} 1 idée directrice par § ? Argument + Exemple (citation/précis) ? Connecteur en tête de § ?
    \item \textbf{Langue :} Accords (sujet/verbe, GN) ? Homophones (à/a, et/est, ou/où, ce/se, son/sont) ? Majuscules (début phrase, noms propres) ? Ponctuation (virgules avant mais/car/donc) ?
    \item \textbf{Présentation :} Saut de ligne clair entre Intro/Parties/Concl ? Paragraphes indentés ? Lisible ? Nom/Numéro candidat ?
    \item \textbf{Sujet :} Hors-sujet ? Contre-sens ? Tous les termes du sujet traités ?
\end{enumerate}
\end{definition}

\courssub{Évaluation formative finale : Auto-positionnement de l'élève}

La séance se termine par un moment de méta-cognition individuel. Chaque élève remplit anonymement (ou nominalement pour le suivi) une grille d'auto-positionnement sur les 8 objectifs opérationnels de la leçon. Cela ancre la responsabilité de l'apprenant.

\begin{exemplebox}[Grille d'Auto-positionnement Élève -- Leçon Dissertation]
\textbf{Consigne :} Pour chaque objectif, cochez votre niveau actuel. \textbf{1 = Non acquis} | \textbf{2 = Fragile (besoin révision)} | \textbf{3 = Acquis (sait faire seul)} | \textbf{4 = Expert (peut expliquer à un pair)}.

\begin{center}
\begin{tabularx}{\linewidth}{|>{\bfseries}X|c|c|c|c|}\hline
\textbf{Objectifs de la leçon (Programme MINESEC)} & \textbf{1} & \textbf{2} & \textbf{3} & \textbf{4} \\ \hline
1. Analyser un sujet de dissertation et formuler une problématique pertinente. & $\square$ & $\square$ & $\square$ & $\square$ \\ \hline
2. Élaborer un plan détaillé (dialectique ou analytique) équilibré et annoncer le plan. & $\square$ & $\square$ & $\square$ & $\square$ \\ \hline
3. Rédiger une introduction complète et normée (4 étapes). & $\square$ & $\square$ & $\square$ & $\square$ \\ \hline
4. Construire un paragraphe argumenté (Idée directrice + Argument + Exemple précis + Connecteur). & $\square$ & $\square$ & $\square$ & $\square$ \\ \hline
5. Mobiliser connaissances précises et citations de \emph{Le Vieux Nègre et la Médaille}. & $\square$ & $\square$ & $\square$ & $\square$ \\ \hline
6. Utiliser les connecteurs logiques variés pour assurer cohérence et cohésion. & $\square$ & $\square$ & $\square$ & $\square$ \\ \hline
7. Rédiger une conclusion synthétique (Bilan + Réponse + Ouverture). & $\square$ & $\square$ & $\square$ & $\square$ \\ \hline
8. Gérer son temps (4h) : produire une copie complète, lisible, relue. & $\square$ & $\square$ & $\square$ & $\square$ \\ \hline
\end{tabularx}
\end{center}
\textbf{Action immédiate :} Identifiez vos \textbf{deux} items les plus bas (1 ou 2). Ce sont vos \textbf{objectifs prioritaires de révision} pour les prochaines semaines. Reportez-les sur votre Fiche Mémo 4.
\end{exemplebox}

\begin{exemplebox}[Statistiques MINESEC Bac Session 2023 -- Corrélations Dissertation]
L'analyse des copies du Baccalauréat technique (Session 2023) révèle une corrélation forte ($r > 0{,}8$) entre la note finale de dissertation ($> 10/20$) et la réussite simultanée à deux items « seuil » :
\begin{itemize}
    \item \textbf{Item « Problématique claire et explicite » :} 87\% des reçus ($>10$) l'avaient validé vs 32\% des recalés.
    \item \textbf{Item « Plan annoncé respecté » :} 84\% des reçus vs 28\% des recalés.
\end{itemize}
\textbf{Enseignement :} La \textbf{structure visible} (Intro/Plan/Concl) pèse plus lourd que la brillance stylistique pour franchir la barre de la moyenne. La remédiation Groupe A est donc la plus « rentable » en points.
\end{exemplebox}

\begin{savaistu}[La Fiche de Progression Harmonisée : un outil de pilotage]
La \textbf{Fiche de progression harmonisée} n'est pas une simple formalité administrative. C'est le document cadre national (MINESEC) qui garantit l'équité entre établissements. L'enseignant y reporte, pour chaque séance : la \textbf{date réelle}, l'\textbf{effectif présent}, le \textbf{contenu effectivement traité} (vs prévisionnel) et des \textbf{observations} (difficultés notées, remédiation faite). Lors du contrôle pédagogique (Inspection), ce document prouve la \textbf{réalité de l'enseignement} et justifie les décisions de remédiation. Ne la remplissez pas à la va-vite en fin d'année ; c'est votre « boîte noire » pédagogique.
\end{savaistu}

\courssub{Bilan de la leçon complète}

Cette leçon 02, étalée sur 15 séances, a conduit l'élève de la découverte de l'œuvre \emph{Le Vieux Nègre et la Médaille} à la production autonome d'une dissertation structurée, argumentée et linguistiquement maîtrisée. Le passage par l'analyse des outils de liaison (langue), la dissection des stratégies argumentatives (rhétorique) et l'entraînement progressif par modules (introduction, paragraphe, conclusion, devoir complet) a visé à transformer une compétence scolaire souvent perçue comme abstraite en une \cle{procédure internalisée}, reproductible en situation d'examen. La séance de bilan/remédiation (cette section 6) est le point d'ancrage qui consolide cet apprentissage en le rendant conscient et transférable pour l'échéance du Probatoire puis du Baccalauréat.

\newpage

\coursec{Activité d'intégration}

\coursec{Leçon 02 : Dissertation -- Produire une dissertation}

\courssub{Objectifs de l'activité}
\noindent Cette activité d'intégration vise à vérifier la capacité de l'élève à mobiliser, de manière articulée et autonome, l'ensemble des savoirs et savoir-faire travaillés au cours de la leçon n°2 dans une situation d'examen simulée (épreuve de dissertation sur corpus, coefficient fort au Probatoire/Baccalauréat technique). Plus précisément, l'élève doit être capable de :
\begin{itemize}
    \item \textbf{Analyser} un corpus hétérogène (littéraire, journalistique, iconique, statistique) en repérant les thèses, arguments et stratégies argumentatives (\cle{réception de textes}, \cle{texte argumentatif}).
    \item \textbf{Formuler} une problématique pertinente et \textbf{construire} un plan détaillé dialectique ou thématique (\cle{production de textes} : dissertation).
    \item \textbf{Rédiger} une introduction complète, deux paragraphes de développement structurés (idée directrice, argument, exemple, connecteurs) et une conclusion synthétique et ouverte (\cle{production de textes} : introduction, paragraphe, conclusion).
    \item \textbf{Maîtriser} la cohérence textuelle par l'usage pertinent des outils de liaison (conjonctions, adverbes, pronoms relatifs, locutions) (\cle{langue française} : morphosyntaxe).
    \item \textbf{S'auto-évaluer} sur des critères objectifs (grille d'évaluation) pour identifier ses axes de progression (\cle{remédiation}).
\end{itemize}

\courssub{Situation}
\noindent \textbf{Contexte :} Nous sommes le 20 mai 2024, à Yaoundé, sur le boulevard de la Réunification. La parade militaire et civile bat son plein. Au premier rang de la tribune officielle, un groupe d'une dizaine d'hommes, le dos voûté, l'uniforme kaki usé mais soigneusement repassé, porte fièrement des médailles qui brillent au soleil : la \emph{Médaille commémorative française}, la \emph{Croix de guerre}, la \emph{Médaille de la Résistance}. Ce sont les \cle{derniers vétérans} camerounais de la Seconde Guerre mondiale (1939-1945) et des guerres d'Indochine/Algérie. Le plus âgé, M. Essomba, 98 ans, peine à se lever pour saluer le Chef de l'État. Il touche sa médaille du bout des doigts tremblants~: un geste machinal, presque sacré.

\noindent En coulisses, loin des caméras, une tout autre réalité. L'association \emph{« Mémoire et Droits des Anciens Combattants »} vient de publier un communiqué cinglant. Selon leurs données (recensement 2023), il ne reste que \textbf{47 vétérans} de la Seconde Guerre mondiale vivants au Cameroun. Leur pension mensuelle moyenne s'élève à \textbf{18~500 FCFA} (environ 28 €), alors que le panier de biens de base pour une personne âgée (médicaments, riz, loyer) est estimé à \textbf{120~000 FCFA} par mois par l'INS (Institut National de la Statistique). Le ratio pension/coût de la vie est donc de \textbf{15,4~\%}. La plupart survivent grâce à la solidarité familiale ou vendent leurs médailles aux antiquaires du marché Mokolo pour 5~000 à 15~000 FCFA l'unité.

\noindent Ce contraste saisissant entre la \emph{visibilité protocolaire} (un jour par an) et l'\emph{invisibilité sociale} (364 jours par an) rappelle tragiquement le destin de \cle{Meka}, le protagoniste du \emph{Vieux nègre et la médaille} de Ferdinand Oyono. Dans la scène finale de l'œuvre, Meka, après avoir tout sacrifié (terre, chèvre, dignité) pour obtenir la médaille de la Légion d'honneur, la reçoit enfin des mains du Commandant. Mais au lieu de la joie attendue, il ressent un vide abyssal~: « \emph{La médaille brillait sur sa poitrine... Mais Meka ne la voyait pas. Il voyait le trou béant de sa case, il voyait sa femme qui mourait, il voyait sa terre vendue.} » La médaille, symbole de la reconnaissance officielle, devient le marqueur cruel de l'abandon réel.

\noindent \textbf{Corpus fourni pour l'épreuve~:}
\begin{enumerate}
    \item \textbf{Texte 1 (Littéraire) :} Extrait de \emph{Le Vieux nègre et la médaille}, Chapitre 10 (Remise de la médaille). « \emph{Le Commandant épingla la médaille... Meka ne bougeait pas... Il pensait à sa case... à sa femme... à sa terre.} »
    \item \textbf{Texte 2 (Journalistique) :} Article \emph{« Cameroun : les vétérans de la Seconde Guerre mondiale réclament leurs droits »}, \emph{Le Jour}, 20 mai 2023. Extrait~: « \emph{La Nation ne saurait se contenter de dépoussiérer les médailles pour les 20 mai. La reconnaissance est un devoir quotidien, pas un décor. Nos aînés meurent dans l'indifférence générale, leurs pensions n'ayant pas été revalorisées depuis 1994.} »
    \item \textbf{Document 3 (Iconique) :} Photographie légendée~: \emph{« Cérémonie du 20 mai 2023, Yaoundé. Vétérans en uniforme, médailles brillantes, visages fatigués. Au second plan, des officiels en costumes sombres, sourires figés. »}
    \item \textbf{Document 4 (Statistique) :} Tableau de données (Source~: MINAS / INS, 2023).
    \begin{center}
    \begin{tabularx}{\linewidth}{|l|X|X|}\hline
    \rowcolor{popblueL} \textbf{Indicateur} & \textbf{Valeur 2023} & \textbf{Évolution vs 2010} \\ \hline
    Nombre de vétérans 1939-45 vivants & 47 & -89\% (étaient 430) \\ \hline
    Pension mensuelle moyenne (FCFA) & 18~500 & Stable (0\% revalorisation) \\ \hline
    Coût vie mensuel senior (FCFA) & 120~000 & +45\% (inflation) \\ \hline
    Taux de couverture pension/vie & 15,4\% & En chute libre (était 28\%) \\ \hline
    Vétérans ayant vendu décorations & 12\% (estimé) & Phénomène nouveau \\ \hline
    \end{tabularx}
    \end{center}
\end{enumerate}

\noindent \textbf{Sujet de dissertation :} \og \textbf{La reconnaissance de la Nation se mesure-t-elle aux décorations officielles ?} \fg

\courssub{Consignes}
\noindent Durée~: 2 heures (séance double). L'élève rend une copie unique structurée en quatre parties distinctes.
\begin{enumerate}
    \item \textbf{Analyse du corpus (30 min) :} Compléter la \emph{Grille d'analyse argumentative} ci-dessous pour les 4 documents. Identifier pour chacun~: \textit{Thèse soutenue}, \textit{Arguments principaux}, \textit{Stratégies/Procédés} (ex~: pathos, logos, éthos, ironie, chiffre, image), \textit{Lien avec le sujet}.
    \begin{center}
    \begin{tabularx}{\linewidth}{|l|X|X|X|X|}\hline
    \rowcolor{popblueL} \textbf{Document} & \textbf{Thèse} & \textbf{Arguments} & \textbf{Stratégies/Procédés} & \textbf{Lien avec sujet} \\ \hline
    Texte 1 (Oyono) & & & & \\ \hline
    Texte 2 (Le Jour) & & & & \\ \hline
    Doc 3 (Photo) & & & & \\ \hline
    Doc 4 (Stats) & & & & \\ \hline
    \end{tabularx}
    \end{center}
    \item \textbf{Construction du devoir (30 min) :} Sur brouillon~:
    \begin{itemize}
        \item Formuler la \textbf{problématique} (une question unique, issue de la tension du corpus).
        \item Élaborer le \textbf{plan détaillé} (2 ou 3 parties, 2 sous-parties chacune, avec idées directrices et exemples précis tirés du corpus ou des connaissances personnelles).
    \end{itemize}
    \item \textbf{Rédaction (50 min) :} Rédiger au propre~:
    \begin{itemize}
        \item L'\textbf{Introduction} complète (accroche, présentation corpus, sujet, problématique, annonce de plan).
        \item Les \textbf{deux premiers paragraphes du développement} (I.A et I.B) : respect strict de la structure \cle{Idée Directrice} -- \cle{Argument} -- \cle{Exemple/Illustration} -- \cle{Connecteurs logiques}.
        \item La \textbf{Conclusion} (bilan, réponse à la problématique, ouverture).
    \end{itemize}
    \item \textbf{Auto-évaluation (10 min) :} Compléter la grille ci-dessous en vous notant sur 4 (1=Insuffisant, 2=Fragile, 3=Satisfaisant, 4=Maîtrisé) et en justifiant d'un mot.
    \begin{center}
    \begin{tabularx}{\linewidth}{|l|c|X|}\hline
    \rowcolor{popblueL} \textbf{Critère} & \textbf{Note /4} & \textbf{Justification / Action corrective} \\ \hline
    Problématique pertinente / Plan cohérent & & \\ \hline
    Introduction complète (4 étapes) & & \\ \hline
    Structure paragraphes (ID / Arg / Ex) & & \\ \hline
    Connecteurs variés et justes & & \\ \hline
    Exemples corpus / connaissances précis & & \\ \hline
    Conclusion (Bilan / Ouverture) & & \\ \hline
    Orthographe / Syntaxe / Présentation & & \\ \hline
    \end{tabularx}
    \end{center}
\end{enumerate}

\courssub{Ressources à mobiliser}
\noindent Rappel synthétique des outils méthodologiques et linguistiques vus en séances 1 à 5.

\begin{methode}[Structure type de l'Introduction (4 étapes)]
\begin{enumerate}
    \item \textbf{Accroche (A) :} Citation, fait d'actualité, chiffre choc, définition, référence littéraire (Oyono). \textit{Ne pas paraphraser le sujet.}
    \item \textbf{Présentation du corpus / Sujet posé (P) :} « Le corpus rassemblé... met en lumière... » + Reformulation du sujet (\cle{La reconnaissance de la Nation...}).
    \item \textbf{Problématique (Pr) :} Question unique, ouverte, issue de la contradiction corpus/sujet. \textit{Ex~: « Au-delà du symbole protocolaire, la médaille ne masque-t-elle pas l'absence de reconnaissance effective ? »}
    \item \textbf{Annonce de plan (Pl) :} « Nous analyserons d'abord... (I), puis nous montrerons... (II), enfin nous verrons... (III) ». \textbf{Verbes d'action} : Analyser, Démontrer, Mettre en évidence, Interroger.
\end{enumerate}
\end{methode}

\begin{methode}[Structure du Paragraphe de développement (Modèle A.E.I. ou I.A.E.)]
\begin{enumerate}
    \item \textbf{Idée Directrice (ID) :} Phrase affirmative, répondant à la problématique, annonçant l'argument principal du paragraphe. (Début de paragraphe, alinéa).
    \item \textbf{Argumentation (Arg) :} Explication, raisonnement, nuance. Utiliser \cle{connecteurs logiques} (cause, conséquence, opposition, concession, exemple).
    \item \textbf{Exemple / Illustration (Ex) :} Renvoi précis au corpus (Texte 1, l. 5 ; Doc 4, chiffre 47) ou connaissance personnelle (histoire, droit, actualité). \textbf{Analyser} l'exemple, ne pas le citer seulement.
    \item \textbf{Mini-synthèse / Transition (optionnel) :} Lien vers l'idée suivante.
\end{enumerate}
\end{methode}

\begin{propriete}[Principaux connecteurs logiques (Outils de liaison) -- Exigibles au Bac]
\begin{center}
\begin{tabularx}{\linewidth}{|l|X|}\hline
\rowcolor{popblueL} \textbf{Relation logique} & \textbf{Outils (Conjonctions, Adverbes, Locutions, Pronoms)} \\ \hline
\cle{Addition / Enchaînement} & et, de plus, furthermore, ensuite, par ailleurs, non seulement... mais encore \\ \hline
\cle{Opposition / Concession} & mais, cependant, pourtant, toutefois, néanmoins, bien que + subj., alors que, malgré \\ \hline
\cle{Cause / Conséquence} & car, parce que, comme, puisque, donc, par conséquent, c'est pourquoi, si bien que \\ \hline
\cle{Exemple / Illustration} & par exemple, notamment, c'est le cas de, ainsi, tel que, comme en témoigne \\ \hline
\cle{Condition / Hypothèse} & si, au cas où, à condition que, pourvu que + subj. \\ \hline
\cle{But} & pour, afin que, de peur que + subj. \\ \hline
\cle{Comparaison} & comme, ainsi que, de même que, de la même manière \\ \hline
\end{tabularx}
\end{center}
\textbf{Règle d'or~:} Varier les connecteurs (ne pas répéter \emph{de plus} 5 fois). Placer le connecteur en tête de phrase ou de proposition pour guider le lecteur.
\end{propriete}

\begin{methode}[Structure de la Conclusion (3 étapes)]
\begin{enumerate}
    \item \textbf{Bilan synthétique (Synthèse) :} Reprendre les idées directrices des grandes parties (I, II, III) sans copier-coller. « Au terme de cette analyse, il apparaît que... »
    \item \textbf{Réponse à la problématique (Réponse) :} Réponse nuancée, définitive. « La reconnaissance ne se mesure donc pas \emph{uniquement} aux décorations... »
    \item \textbf{Ouverture (Ouverture) :} Perspective élargie (autre œuvre, autre pays, actualité brûlante, question philosophique). « Cette réflexion invite à repenser le \emph{devoir de mémoire} non comme un rite, mais comme une politique publique... »
\end{enumerate}
\end{methode}

\courssub{Démarche et corrigé commenté}
\noindent Voici une proposition de résolution complète, étape par étape, servant de \cle{modèle de référence} pour l'enseignant et de support de remédiation pour l'élève.

\begin{exoresolu}[Atelier Intégration : 'Plaidoirie pour les oubliés de l'Histoire']
\textbf{Énoncé de l'activité (rappel) :} À partir d'un corpus de 4 documents (extrait romanesque \emph{Le Vieux nègre et la médaille}, article \emph{Le Jour} 2023, photo 20 mai 2023, statistiques MINAS/INS 2023), traiter le sujet~: \og La reconnaissance de la Nation se mesure-t-elle aux décorations officielles ? \fg en respectant la méthode dissertationnaire (introduction, développement en 2 paragraphes minimum, conclusion) et en mobilisant les outils de liaison.

\tcblower

\textbf{ÉTAPE 1 : Analyse du corpus (Grille commentée)}

\begin{center}
\begin{tabularx}{\linewidth}{|l|X|X|X|X|}\hline
\rowcolor{popblueL} \textbf{Document} & \textbf{Thèse} & \textbf{Arguments principaux} & \textbf{Stratégies / Procédés} & \textbf{Lien avec sujet} \\ \hline
\textbf{Texte 1 : Oyono (Extrait Ch. 10)} & La médaille officielle est un leurre, une coquille vide qui ne répare pas les sacrifices réels. & 1. Sacrifices consentis (terre, chèvre, femme). 2. Cérémonie solennelle mais froide. 3. Pensées intérieures de Meka~: le vide, la misère restante. & \textbf{Ironie dramatique} (lecteur sait, Meka subit). \textbf{Contraste} visuel (médaille brillante / case vide). \textbf{Focalisation interne} (accès à la pensée). \textbf{Champ lexical}~: misère, vide, trou, mort. & Illustre la \textbf{dissociation} entre le symbole (médaille) et la réalité vécue (reconnaissance nulle). Le \emph{Vieux nègre} est l'archétype du vétéran trahi. \\ \hline
\textbf{Texte 2 : Le Jour (2023)} & La reconnaissance nationale est un devoir quotidien (social, financier), non un décor annuel. & 1. Critique du « 20 mai » (événementiel). 2. Absence de revalorisation des pensions depuis 1994. 3. Mort dans l'indifférence. & \textbf{Argumentation directe} (logos~: dates, faits). \textbf{Pathos}~: « aînés meurent », « indifférence ». \textbf{Anaphore}~: « La Nation ne saurait... ». \textbf{Ton polémique / Plaidoyer}. & Apporte la \textbf{dimension juridique et sociale}~: la médaille sans pension = reconnaissance formelle vide. Prolonge Oyono dans le réel actuel. \\ \hline
\textbf{Doc 3 : Photo 20 mai 2023} & Le cérémonial officiel met en scène l'honneur mais laisse voir la fatigue et la précarité. & 1. Uniformes/médailles = mise en scène de la gloire. 2. Visages fatigués, dos voûtés = réalité du corps usé. 3. Contraste arrière-plan (officiels confortables). & \textbf{Composition}~: Premier plan (vétérans) vs Second plan (pouvoir). \textbf{Lumière}~: reflets médailles vs ombres visages. \textbf{Non-verbal}~: regard baissé vs sourires figés. & \textbf{Preuve visuelle} du fossé. La médaille \emph{brille} (officiel) mais le visage \emph{fatigué} (réel). Illustration parfaite du « trou béant » de Meka. \\ \hline
\textbf{Doc 4 : Stats MINAS/INS 2023} & La reconnaissance matérielle est dérisoire (15,4\% couverture) et en régression. & 1. Disparition démographique (47 vs 430). 2. Gel pensions (0\% reval. 13 ans). 3. Inflation (+45\%). 4. Vente médailles (symbole ultime de détresse). & \textbf{Logos pur}~: Chiffres, pourcentages, variations. \textbf{Tableau} = objectivité scientifique. \textbf{Choc}~: 15,4\% / Vente médailles. & \textbf{Preuve quantitative} irréfutable. La médaille devient \emph{monnaie de survie} (vente), perdant toute valeur symbolique. Mesure exacte de l'abandon. \\ \hline
\end{tabularx}
\end{center}

\textbf{ÉTAPE 2 : Problématique et Plan détaillé}

\noindent \textbf{Problématique~:} \og \emph{Dans quelle mesure la remise de décorations officielles, loin de sceller une véritable reconnaissance nationale, ne risque-t-elle pas de masquer, par son apparat symbolique, l'abandon matériel et mémoriel des acteurs de l'Histoire ?} \fg

\noindent \textbf{Plan détaillé (Dialectique~: Thèse / Antithèse / Synthèse dépassement)}

\textbf{I. La décoration officielle~: un symbole nécessaire mais insuffisant de la reconnaissance nationale (Thèse~: la fonction symbolique).}
\begin{itemize}
    \item \textbf{I.A. L'institutionnalisation de la mémoire~: le rôle unificateur du rite.} (ID) La médaille matérialise la dette de la Nation et inscrit l'individu dans l'Histoire collective. (Arg) \emph{Ex~:} Texte 1 (cérémonie solennelle, Commandant), Doc 3 (défilé 20 mai, cohésion nationale), connaissance~: création de l'Ordre de la Valeur (Cameroun).
    \item \textbf{I.B. La valorisation du mérite individuel~: une justice symbolique.} (ID) Elle distingue le sacrifice personnel et offre une fierté identitaire. (Arg) \emph{Ex~:} Texte 1 (fierté initiale de Meka, « il voulait la médaille »), Doc 3 (port de l'uniforme, dignité), connaissance~: rôle des décorations dans l'armée (motivation, hiérarchie).
\end{itemize}

\textbf{II. L'illusion du symbole~: quand la médaille masque l'abandon réel et la dette impayée (Antithèse~: la réalité matérielle).}
\begin{itemize}
    \item \textbf{II.A. Le décalage abyssal entre honneur protocolaire et précarité vitale.} (ID) La médaille ne nourrit pas, ne soigne pas, ne loge pas. (Arg) \emph{Ex~:} Texte 1 (Meka~: « trou béant de sa case », femme qui meurt), Doc 4 (Pension 18~500 FCFA vs Vie 120~000 FCFA = 15,4\%), Doc 2 (gel pensions 1994).
    \item \textbf{II.B. La marchandisation de la mémoire~: le vendeur de médailles.} (ID) L'extrême détresse transforme le symbole sacré en vulgaire marchandise. (Arg) \emph{Ex~:} Doc 4 (12\% vendent médailles à Mokolo 5~000 FCFA), Texte 2 (communiqué association), Connaissance~: vente décorations vétérans français années 1950 (Indochine).
\end{itemize}

\textbf{III. Vers une reconnaissance effective~: dépasser le symbole par la réparation juste et la mémoire vivante (Synthèse~: l'éthique du devoir).}
\begin{itemize}
    \item \textbf{III.A. La réparation matérielle comme condition de la reconnaissance vraie.} (ID) Revaloriser les pensions, soins gratuits, logement = traduction concrète de la dette. (Arg) \emph{Ex~:} Doc 2 (revendications « devoir quotidien »), Connaissance~: Loi de programmation militaire 2019-2025 (France) revalorisant pensions, modèle à adapter.
    \item \textbf{III.B. La mémoire active~: éducation, transmission, histoire orale.} (ID) Faire vivre l'Histoire par les programmes scolaires, musées, témoignages, non par le défilé unique. (Arg) \emph{Ex~:} Texte 1 (Oyono écrit pour transmettre), Doc 3 (photo comme archive), Connaissance~: Journée du 11 novembre / 8 mai en France, devoir de mémoire scolaire.
\end{itemize}

\textbf{ÉTAPE 3 : Rédaction de l'Introduction}

\begin{quote}
\textbf{[Accroche]} « \emph{La médaille brillait sur sa poitrine... Mais Meka ne la voyait pas. Il voyait le trou béant de sa case.} » Par cette phrase finale du \emph{Vieux nègre et la médaille}, Ferdinand Oyono condense en 1956 la tragédie intemporelle du soldat africain~: l'honneur officiel comme paravent de la misère réelle. 

\textbf{[Présentation corpus / Sujet]} Le corpus rassemblé — extrait romanesque, article de presse \emph{Le Jour} (2023), photographie du 20 mai, données statistiques MINAS/INS — dresse un constat convergent~: au Cameroun comme ailleurs, la reconnaissance nationale s'exhibe dans l'éclat des décorations le jour de la fête nationale, mais s'efface le reste de l'année derrière des pensions gelées depuis 1994 (18~500 FCFA pour un coût de vie de 120~000 FCFA). Le sujet invite à interroger la valeur de ce symbole~: \textbf{La reconnaissance de la Nation se mesure-t-elle aux décorations officielles ?}

\textbf{[Problématique]} Dans quelle mesure la remise de décorations officielles, loin de sceller une véritable reconnaissance nationale, ne risque-t-elle pas de masquer, par son apparat symbolique, l'abandon matériel et mémoriel des acteurs de l'Histoire ?

\textbf{[Annonce de plan]} Nous analyserons d'abord la fonction symbolique irremplaçable de la décoration comme ciment de la mémoire collective (I). Nous démontrerons ensuite que ce symbole devient mensonge lorsqu'il occulte l'abandon social et la précarité indigne des vétérans (II). Enfin, nous montrerons que la vraie reconnaissance exige de traduire le symbole en actes concrets~: réparation matérielle et mémoire vivante (III).
\end{quote}

\textbf{ÉTAPE 4 : Rédaction de deux paragraphes de développement (I.A et I.B)}

\noindent \textbf{I. La décoration officielle~: un symbole nécessaire mais insuffisant de la reconnaissance nationale}

\noindent \textbf{I.A. L'institutionnalisation de la mémoire~: le rôle unificateur du rite.}
\noindent \textbf{[ID]} La décoration officielle remplit une fonction anthropologique majeure~: elle institue la dette de la Nation et inscrit le sacrifice individuel dans la grande Histoire collective. \textbf{[Arg]} \emph{En effet}, par le rituel solennel de la remise, l'État transforme un acte de guerre personnel en patrimoine commun~; la médaille devient alors un \cle{lieu de mémoire} (Pierre Nora) visible, lisible par tous, qui dit au citoyen~: « Celui-ci a donné sa jeunesse pour que tu vives en paix. » \textbf{[Ex]} \emph{C'est ce que montre} le Texte 1 lors de la cérémonie où le Commandant épingle la Légion d'honneur sur la poitrine de Meka~: la solennité de l'uniforme, le protocole militaire, le silence respectueux créent un moment de \emph{communion nationale} où le « vieux nègre » redevient, l'espace d'un instant, un héros à part entière. \emph{De même,} le Document 3 capture cette dimension~: les vétérans défilent le 20 mai, médailles au cou, sous les drapeaux, incarnant devant la tribune officielle l'unité de la patrie. \textbf{[Trans]} \emph{Toutefois,} si le rite fonde le lien social, il ne saurait en être l'unique fondement.

\noindent \textbf{I.B. La valorisation du mérite individuel~: une justice symbolique.}
\noindent \textbf{[ID]} Au-delà du collectif, la médaille rend justice au mérite personnel et restaure la dignité de l'acteur de l'Histoire. \textbf{[Arg]} \emph{En effet,} le soldat n'est pas un numéro~; la distinction individuelle (Croix de guerre, Médaille militaire) reconnaît sa bravoure spécifique, son risque consenti, sa douleur propre. Elle est une \emph{réparation symbolique} de l'invisible. \textbf{[Ex]} \emph{Ainsi,} dans le Texte 1, Meka \emph{lui-même}, avant la désillusion finale, a tout sacrifié — sa terre, sa chèvre, l'honneur de sa femme — \emph{précisément} pour obtenir cette médaille~: « \emph{Il la voulait, cette médaille. Il la voulait de toutes ses forces.} » Ce désir prouve que le symbole a une valeur propre, qu'il comble un besoin de reconnaissance identitaire. \emph{Par ailleurs,} la photographie (Doc 3) montre des visages droits, des uniformes impeccables~: la fierté d'avoir « fait son devoir » se lit sur les traits fatigués. \textbf{[Mini-synthèse]} La décoration est donc le \emph{certificat moral} de la Nation~; mais un certificat sans contrepartie matérielle risque de devenir un chèque sans provision.

\textbf{ÉTAPE 5 : Rédaction de la Conclusion}

\begin{quote}
\textbf{[Bilan synthétique]} Au terme de cette analyse, il apparaît que la décoration officielle joue un double rôle contradictoire. Elle est d'abord le \emph{lien sacré} (I) qui unit la Nation à ses enfants par le rite et l'honneur, restaurant la dignité du combattant (I.A, I.B). Mais elle devient, dans le contexte camerounais actuel, un \emph{leurre cruel} (II) lorsque l'éclat du métal masque la rouille des pensions (II.A) et pousse jusqu'à la vente du symbole même de l'honneur (II.B).

\textbf{[Réponse à la problématique]} La reconnaissance de la Nation ne se mesure \emph{donc pas} aux décorations officielles \emph{seules}~: celles-ci sont la \emph{condition nécessaire} (le signe), mais non la \emph{condition suffisante} (la substance). Une médaille sur une poitrine qui a faim est une insulte à l'Histoire.

\textbf{[Ouverture]} Cette réflexion nous oblige à repenser le \emph{devoir de mémoire} non comme une cérémonie annuelle, mais comme une \textbf{politique publique continue}~: revalorisation immédiate des pensions (alignement sur le SMIC ou indice vieillesse), couverture santé intégrale (carte vétéran), et surtout inscription systématique de leurs témoignages dans les programmes scolaires (Français, Histoire, Éducation à la citoyenneté). Comme l'écrivait Oyono, « \emph{Il faut que les jeunes sachent.} » La vraie médaille, celle qui ne s'use pas, c'est la mémoire vivante transmise aux générations futures.
\end{quote}

\textbf{ÉTAPE 6 : Auto-évaluation type (Grille remplie modèle)}

\begin{center}
\begin{tabularx}{\linewidth}{|l|c|X|}\hline
\rowcolor{popblueL} \textbf{Critère} & \textbf{Note /4} & \textbf{Justification / Action corrective} \\ \hline
Problématique pertinente / Plan cohérent & 4 & Tension claire symbole/réel. Plan dialectique équilibré (3 parties). \textit{OK.} \\ \hline
Introduction complète (4 étapes) & 4 & Accroche Oyono intégrée, corpus cité, sujet reformulé, probl. unique, plan annoncé avec verbes d'action. \textit{OK.} \\ \hline
Structure paragraphes (ID / Arg / Ex) & 3 & ID claires, arguments développés, exemples précis (texte + doc). \textit{Attention~:} Mini-synthèse I.B un peu courte. \\ \hline
Connecteurs variés et justes & 4 & \emph{En effet, C'est ce que montre, De même, Toutefois, Ainsi, Par ailleurs, Par conséquent.} Variété respectée. \textit{OK.} \\ \hline
Exemples corpus / connaissances précis & 4 & Citations exactes (l. Texte 1), chiffres Doc 4 (15,4\%), référence Nora, loi française. \textit{OK.} \\ \hline
Conclusion (Bilan / Ouverture) & 4 & Synthèse dialectique, réponse nuancée (« nécessaires mais pas suffisantes »), ouverture concrète (politique publique, école). \textit{OK.} \\ \hline
Orthographe / Syntaxe / Présentation & 3 & 2 fautes d'inattention (accords participes). \textit{Action~:} Relire à l'envers (dernière phrase $\to$ première). \\ \hline
\end{tabularx}
\end{center}
\end{exoresolu}

\courssub{Bilan de l'activité}
\noindent Cette activité d'intégration, conçue comme une \emph{simulation d'examen} (épreuve unique de 2h), a permis de mettre en évidence les acquis transversaux suivants, conformément aux attendus du Probatoire/Baccalauréat technique~:

\begin{aretenir}[Compétences validées par l'activité]
\begin{itemize}
    \item \textbf{Compétence analytique (Réception) :} L'élève a démontré sa capacité à croiser des sources hétérogènes (littérature, presse, image, statistiques) pour en extraire une thèse commune et des nuances, en utilisant le vocabulaire précis de l'analyse argumentative (\cle{thèse}, \cle{argument}, \cle{stratégie}, \cle{procédé}, \cle{pathos/logos/éthos}).
    \item \textbf{Compétence méthodologique (Production) :} Le passage obligé par le brouillon (problématique + plan détaillé) a sanctionné la maîtrise de l'architecture dissertationnaire~: introduction en 4 temps, paragraphes structurés ID/Arg/Ex, conclusion en 3 temps. L'usage de la \emph{grille d'auto-évaluation} a ancré la méta-cognition.
    \item \textbf{Compétence linguistique (Langue) :} La contrainte explicite d'utiliser des connecteurs variés (addition, opposition, conséquence, illustration) a forcé l'élève à sortir du « \emph{et / mais / donc } » répétitif pour produire un texte \cle{cohérent} et \cle{cohesif}, critère majeur de notation au Bac.
    \item \textbf{Compétence contextuelle (Ancrage) :} Le corpus 100\% camerounais (Oyono, \emph{Le Jour}, 20 mai Yaoundé, données MINAS/INS, marché Mokolo) a permis de vérifier la capacité à mobiliser des \cle{connaissances du contexte} (histoire coloniale, actualité sociale, géographie locale) pour illustrer et argumenter, sortant de l'abstraction pure.
\end{itemize}
\end{aretenir}

\begin{attention}[Points de vigilance pour la remédiation (Séance 6)]
\begin{itemize}
    \item \textbf{Problématique « hors-sujet » ou « fermée » :} Certains élèves formulent encore une question à réponse Oui/Non (ex~: « La médaille est-elle une reconnaissance ? »). \textbf{Remédiation~:} Exercices de transformation de sujets en questions ouvertes (« Dans quelle mesure... », « Comment... », « Pourquoi... »).
    \item \textbf{Paragraphe « catalogue d'exemples » :} L'élève aligne citations/documents sans \cle{idée directrice} unificatrice ni \cle{raisonnement} (argument). \textbf{Remédiation~:} Entraînement systématique au modèle \textbf{ID $\to$ Arg $\to$ Ex $\to$ Analyse Ex} sur des micro-sujets (1 paragraphe = 15 min).
    \item \textbf{Connecteurs « décoratifs » :} Utilisation de \emph{par ailleurs} ou \emph{de plus} sans logique réelle (ex~: opposition marquée par \emph{de plus}). \textbf{Remédiation~:} Travail sur la \cle{sémantique des connecteurs} (tableau de correspondance Relation $\leftrightarrow$ Outils) + exercices de complétion de textes à trous.
    \item \textbf{Gestion du temps :} 2h est court. Beaucoup n'ont pas fini la conclusion. \textbf{Remédiation~:} Chronométrage imposé aux étapes (30/30/50/10) lors des prochains entraînements.
\end{itemize}
\end{attention}

\begin{savaistu}[Le Vieux Nègre et la Médaille : Contexte historique camerounais]
\emph{Le Vieux nègre et la médaille} (1956) de Ferdinand Oyono (1929--2010) est le premier roman d'une trilogie (suivi de \emph{Une vie de boy}, \emph{Chemin d'Europe}). Il décrit l'enrôlement forcé et l'exploitation des \cle{tirailleurs sénégalais} (terme générique pour soldats africains) par l'administration coloniale française. Le personnage de Meka est inspiré de la réalité des \textbf{« vétérans oubliés »}~: sur les \textbf{18~000 Camerounais} enrôlés en 1939-1940 (chiffres historiens), beaucoup n'ont jamais reçu leurs droits. La loi française de 1959 (puis 2002, 2010) a tenté de régler la question des \emph{pensions de cristallisation} (gelées au taux 1959), mais le contentieux persiste. Au Cameroun, le \textbf{Ministère des Affaires Sociales (MINAS)} gère le fichier des anciens combattants. Le 20 mai (Fête nationale) et le 11 novembre (Armistice) sont les deux moments de visibilité protocolaire. L'activité ancre ainsi la littérature dans une \cle{réalité citoyenne} brûlante.
\end{savaistu}

\newpage

\coursec{Méthodes et exercices résolus}

\coursec{Dissertation : produire une dissertation}

\courssub{Lecture méthodique de « Le Vieux nègre et la médaille » : appropriation de l'œuvre}

\begin{methode}[Lecture méthodique en quatre temps (programme MINESEC)]
\textbf{Objectif :} Maîtriser le protocole de lecture analytique imposé au Probatoire pour l'œuvre au programme.
\begin{enumerate}
    \item \textbf{Lecture 1 : Découverte et contextualisation.} Lire le texte sans s'arrêter. Identifier : auteur (Ferdinand Oyono), genre (roman court / nouvelle), date (1956), cadre spatio-temporel (Cameroun, période coloniale), personnage principal (Meka). Noter l'\cle{incipit} et l'\cle{explicit}.
    \item \textbf{Lecture 2 : Analyse du récit et des personnages.} Repérer la structure narrative (schéma actanciel : quête de la médaille, épreuves, dénouement). Étudier Meka (évolution : naïveté $\to$ lucidité), le Commandant, Kelara, les fils. Relever les \cle{thèmes majeurs} : colonisation, aliénation, dignité, tradition vs modernité, ironie.
    \item \textbf{Lecture 3 : Étude de la langue et du style.} Identifier les \cle{procédés ironiques} (antiphrase, sous-entendu, décalage narratif), le \cle{point de vue} (focalisation interne sur Meka, narrateur omniscient ironique), le \cle{registre} (satirique, pathétique, réaliste). Analyser le \cle{lexique} (champ lexical de l'administration, de la tradition, du corps).
    \item \textbf{Lecture 4 : Synthèse et problématisation.} Construire une fiche de lecture structurée : résumé détaillé, analyse des personnages, thèmes, style, portée. Formuler une \cle{problématique} type dissertation (ex. : « En quoi ce roman est-il une critique de l'assimilation coloniale ? »).
\end{enumerate}
\textbf{Reflexe Bac :} Toujours citer le texte (numéro de chapitre) pour étayer toute affirmation. Distinguer \emph{ce que dit le texte} et \emph{ce que le texte fait dire}.
\end{methode}

\begin{exoresolu}[Analyse d'un extrait : La remise de la médaille (Chapitre 6)]
\textbf{Énoncé :} Lisez l'extrait suivant (chapitre 6, remise de la médaille). \textbf{1.} Identifiez le registre dominant et justifiez par deux procédés. \textbf{2.} Analysez la position du narrateur vis-à-vis de Meka. \textbf{3.} En quoi cet épisode illustre-t-il l'ironie du titre ?

\tcblower
\textbf{Solution détaillée :}
\textbf{1. Registre et procédés.}
Le registre dominant est \textbf{l'ironique (satirique)}.
\begin{itemize}
    \item \textbf{Antiphrase / Décalage :} La solennité administrative (« cérémonie officielle », « décoré de la médaille ») contraste avec la réalité vécue par Meka (vêtements neufs qui le gênent, chaleur, incompréhension du sens réel). Le texte dit « honneur » mais montre « humiliation ».
    \item \textbf{Hyperbole / Dérision :} La description de Meka « gonflé d'orgueil » puis « petit vieux noir » réduit par la médaille physique (poids sur la poitrine) et symbolique (poids de la soumission).
\end{itemize}
\textbf{2. Position du narrateur.}
Le narrateur adopte une \textbf{focalisation interne} sur Meka (on accède à ses pensées : « il se sentait grand », « il comprenait enfin »). Mais cette subjectivité est encadrée par une \textbf{ironie narrative distanciée} : le narrateur sait ce que Meka ignore (la médaille = marque de servilité). Le lecteur est en position de \cle{supériorité cognitive} (ironie dramatique).
\textbf{3. Illustration de l'ironie du titre.}
Le titre « Le Vieux nègre et la médaille » oppose la dignité humaine (« Vieux ») à la réification coloniale (« nègre », « médaille »). L'épisode montre Meka \emph{recevant} l'objet qui le \emph{définit} aux yeux du colonisateur comme un « bon indigène ». La médaille, censée honorer, devient le symbole visible de son aliénation. C'est le cœur de la \cle{thèse} d'Oyono : la colonisation vole l'identité pour donner une breloque.
\textbf{Conclusion :} L'extrait condense la thèse du roman : l'ironie narrative dévoile la violence symbolique de l'ordre colonial.
\end{exoresolu}

\courssub{Maîtrise des outils de liaison : cohérence textuelle et articulation argumentative}

\begin{methode}[Les connecteurs logiques : classification et emploi dans la dissertation]
\textbf{Objectif :} Assurer la fluidité et la rigueur logique du développement (exigence « cohérence textuelle » du Bac).
\begin{enumerate}
    \item \textbf{Conjonctions de coordination (mais, ou, et, donc, or, ni, car).} Lient propositions de même nature. \textbf{Attention :} « Car » n'ouvre pas une phrase (style soutenu). Préférer « Car » en milieu de phrase ou « Car » $\to$ « En effet / Parce que » en tête de phrase.
    \item \textbf{Conjonctions de subordination.} Introduisent une proposition subordonnée.
        \begin{itemize}
            \item \textbf{Cause :} \emph{parce que, comme, puisque, dès lors que.}
            \item \textbf{Conséquence :} \emph{si bien que, de sorte que, au point que.}
            \item \textbf{Condition :} \emph{si, pourvu que, à condition que.}
            \item \textbf{Concession :} \emph{bien que, quoique, même si, quelque ... que.}
            \item \textbf{Finalité :} \emph{afin que, pour que, de peur que.}
            \item \textbf{Comparaison :} \emph{comme, ainsi que, de même que.}
        \end{itemize}
    \item \textbf{Adverbes et locutions adverbiales de liaison (connecteurs phraséologiques).} Indispensables pour \cle{structurer les paragraphes} et \cle{relier les paragraphes} (début de paragraphe).
        \begin{center}
        \begin{tabularx}{\linewidth}{|l|X|}\hline
        \rowcolor{popblueL} \textbf{Fonction logique} & \textbf{Connecteurs types (Bac)} \\ \hline
        \textbf{Énumération / Ordre} & \emph{Premièrement, En premier lieu, D'une part, D'autre part, Ensuite, Enfin, Pour terminer} \\ \hline
        \textbf{Addition / Renforcement} & \emph{De plus, En outre, Par ailleurs, Qui plus est, D'ailleurs, Non seulement ... mais encore} \\ \hline
        \textbf{Opposition / Concession} & \emph{Cependant, Néanmoins, Pourtant, Certes ... mais, En revanche, Au contraire, Quoi qu'il en soit} \\ \hline
        \textbf{Cause / Explication} & \emph{En effet, C'est pourquoi, Effectivement, Parce que, Du fait que} \\ \hline
        \textbf{Conséquence / Conclusion partielle} & \emph{Par conséquent, Donc, Ainsi, C'est ainsi que, Il en résulte que, En définitive} \\ \hline
        \textbf{Illustration / Exemple} & \emph{Par exemple, Ainsi, Notamment, C'est le cas de, Tel que} \\ \hline
        \end{tabularx}
        \end{center}
    \item \textbf{Pronoms relatifs et reprise anaphorique.} \emph{Qui, que, dont, où, lequel, auquel...} + \emph{ceci, cela, ce dernier, celui-ci, tel}. Éviter les répétitions lourdes, assurer la \cle{cohérence référentielle}.
\end{enumerate}
\textbf{Formule à retenir :} \textbf{1 Idée directrice = 1 Paragraphe = 1 Connecteur d'ouverture + Arguments (connecteurs internes) + Exemples (connecteurs d'illustration) + Mini-conclusion (connecteur de conséquence).}
\end{methode}

\begin{exoresolu}[Réécriture et cohésion : d'un brouillon à un paragraphe type Bac]
\textbf{Énoncé :} Voici des notes de brouillon pour un paragraphe sur « La médaille, symbole d'aliénation ». Réécrivez en un paragraphe unique, cohérent, de niveau Bac (connecteurs variés, syntaxe soignée).
\textit{Notes :} Meka fier. Médaille lourde. Il comprend pas. Commandant dit "honneur". En fait soumission. Exemple : Meka porte médaille enterrement fils. Ironie. Conclusion : médaille = chaîne.

\tcblower
\textbf{Solution détaillée :}
\textbf{Paragraphe rédigé :}
\emph{« D'emblée, la médaille offerte à Meka apparaît comme le \textbf{symbole majeur de l'aliénation coloniale}. \textbf{En effet}, si le vieil homme l'accueille d'abord avec une fierté naïve — \emph{« il se sentait grand »} —, \textbf{cependant} la lourdeur physique de l'objet (\emph{« un poids sur sa poitrine »}) préfigure le fardeau symbolique. \textbf{Par ailleurs}, le discours du Commandant, empreint de paternalisme (\emph{« honneur », « loyauté »}), masque mal la réalité d'une soumission consentie. \textbf{C'est ainsi que}, lors de l'enterrement de son fils Kelara, Meka, paré de sa médaille, incarne tragiquement la contradiction : décoré par l'oppresseur, il pleure une victime de cet ordre. \textbf{En définitive}, la médaille, loin d'honorer, enchaîne Meka à une identité d'« indigène méritant », lui déniant sa qualité d'homme libre. »}

\textbf{Analyse des connecteurs utilisés :}
\begin{itemize}
    \item \emph{D'emblée} (Ouverture / Thématisation).
    \item \emph{En effet} (Justification / Cause).
    \item \emph{Cependant} (Concession / Opposition interne : fierté vs poids).
    \item \emph{Par ailleurs} (Argument supplémentaire / Changement de focale : le discours).
    \item \emph{C'est ainsi que} (Conséquence narrative / Illustration concrète).
    \item \emph{En définitive} (Synthèse / Conclusion du paragraphe).
\end{itemize}
\textbf{Point de vigilance :} La citation (\emph{« un poids sur sa poitrine »}) est intégrée à la phrase (incorporation), pas plaquée. Les connecteurs ne sont pas tous en tête de phrase (variété syntaxique).
\end{exoresolu}

\courssub{Le texte argumentatif : stratégies d'analyse et de construction}

\begin{methode}[Analyse des stratégies argumentatives (Réception) et Construction (Production)]
\textbf{Objectif :} Décrypter l'argumentation dans un texte (Bac : commentaire / contraction) et l'appliquer en dissertation.
\begin{enumerate}
    \item \textbf{Identifier la thèse (position de l'énonciateur).} Explicite (affirmée) ou implicite (à inférer). Question : « Que veut-il me faire admettre ? »
    \item \textbf{Repérer les arguments (raisons logiques).} Les classifier :
        \begin{itemize}
            \item \textbf{Arguments d'autorité} (référence à un expert, texte sacré, loi).
            \item \textbf{Arguments factuels / réalistes} (faits, statistiques, exemples historiques, observation).
            \item \textbf{Arguments logiques / rationnels} (déduction, analogie, \emph{a fortiori}, \emph{a contrario}, réduction à l'absurde).
            \item \textbf{Arguments axiologiques} (valeurs morales, esthétiques, patriotiques, humanistes).
            \item \textbf{Arguments émotionnels / pathétiques} (peur, pitié, indignation, fierté).
        \end{itemize}
    \item \textbf{Analyser l'énonciation (ethos).} Qui parle ? Quelle image de lui-même construit l'auteur ? (Expert, victime, citoyen, sage, ironique). \cle{Credibilité} (ethos) = force de persuasion.
    \item \textbf{Étudier la disposition (structure).} Ordre des arguments : \textbf{déductif} (thèse $\to$ arguments), \textbf{inductif} (arguments $\to$ thèse), \textbf{dialectique} (thèse adverse $\to$ réfutation $\to$ thèse propre), \textbf{analytique} (analyse des termes $\to$ thèse).
    \item \textbf{Repérer les procédés rhétoriques.} Questions rhétoriques, anaphore, parallélisme, antithèse, métaphore filée, ironie, euphémisme, hyperbole. Effet visé : \emph{captiver, convaincre, persuader, émouvoir}.
\end{enumerate}
\textbf{Transfert Dissertation :} Votre développement \emph{est} un texte argumentatif. Votre \cle{problématique} = votre thèse. Vos \cle{idées directrices} = vos arguments principaux. Vos \cle{exemples} = vos arguments factuels/illustratifs. Vos \cle{connecteurs} = votre disposition.
\end{methode}

\begin{exoresolu}[Analyse d'un paragraphe argumentatif type « Éditorial »]
\textbf{Énoncé :} Analysez le paragraphe suivant (sujet : « L'école technique au Cameroun »). Identifiez : Thèse, 3 arguments (type + citation), Structure, 2 procédés rhétoriques, Éthos.
\textit{Texte :} « \emph{Former des techniciens, ce n'est pas reléguer des élèves. C'est, au contraire, armer la nation. D'abord, les pays émergents (Corée, Allemagne) le prouvent : l'industrie a besoin de mains expertes, pas seulement de théoriciens. Ensuite, nos lycées techniques forment des créateurs d'emplois : 60\% des lauréats BTS s'installent à leur compte (MINESUP 2023). Enfin, refuser la voie technique, c'est accepter la dépendance technologique éternelle. Choisissons l'excellence pratique.} »

\tcblower
\textbf{Solution détaillée :}
\begin{itemize}
    \item \textbf{Thèse :} « Former des techniciens [...] c'est armer la nation » (Valorisation de l'enseignement technique).
    \item \textbf{Arguments :}
        \begin{enumerate}
            \item \textbf{Argument d'autorité / Exemple étranger (Factuel) :} « les pays émergents (Corée, Allemagne) le prouvent ». Type : \emph{Exemplum} / Preuve par l'exemple réussi.
            \item \textbf{Argument factuel / Statistique (Chiffré) :} « 60\% des lauréats BTS s'installent à leur compte (MINESUP 2023) ». Type : \emph{Donnée objective} / Preuve par le résultat local.
            \item \textbf{Argument logique / A contrario (Peur/Réalisme) :} « refuser la voie technique, c'est accepter la dépendance technologique éternelle ». Type : \emph{Argument des conséquences néfastes} (reductio ad absurdum / menace).
        \end{enumerate}
    \item \textbf{Structure :} \textbf{Inductive / Accumulation.} Ouverture par une antithèse (définition par le contraire), puis 3 arguments croissants en force (externe $\to$ interne $\to$ existentiel), conclusion impérative.
    \item \textbf{Procédés rhétoriques :}
        \begin{itemize}
            \item \textbf{Antithèse / Parallélisme :} « ce n'est pas reléguer... C'est, au contraire, armer ». Structure binaire valorisante.
            \item \textbf{Impératif exhortatif final :} « Choisissons l'excellence pratique ». Appel à l'action collective (inclusif « nous »).
        \end{itemize}
    \item \textbf{Éthos :} \textbf{Citoyen responsable, lucide, volontariste.} Il manie le chiffre (rigueur), la comparaison internationale (ouverture), le « nous » (solidarité). Il se pose en \cle{guide éclairé}.
\end{itemize}
\textbf{Conclusion :} Texte efficace car arguments variés (externe, interne, logique), structure claire, ton mobilisateur. Modèle pour une \cle{conclusion de dissertation} (appel à l'action / ouverture).
\end{exoresolu}

\courssub{Architecture de la dissertation : Introduction, Développement, Conclusion}

\begin{methode}[L'Introduction en 4 étapes (Modèle « Entonnoir » exigé au Bac)]
\textbf{Objectif :} Rédiger une introduction parfaite (environ 15-20 lignes), notée sur 4-5 points.
\begin{enumerate}
    \item \textbf{Amorce (Accroche) :} 1 à 2 phrases. \textbf{Interdiction :} « Depuis la nuit des temps... », « De tout temps... », définitions du dictionnaire. \textbf{Autorisé :} Citation d'auteur (si pertinente), fait d'actualité précis, référence à l'œuvre au programme (\emph{Le Vieux nègre...}), paradoxe, chiffre marquant, anecdote courte. \textit{But :} Susciter l'intérêt, ancrer le sujet.
    \item \textbf{Présentation du sujet / Définition des termes :} Reprendre les mots-clés du sujet, les définir (sens commun / sens précis), en montrer la tension ou la polysémie. Situer le \cle{champ conceptuel} (littéraire, sociologique, philosophique...).
    \item \textbf{Problématique :} \textbf{La question unique, ouverte, débattue, qui guide tout le devoir.} Elle ne doit pas être une simple reformulation du sujet. Elle contient souvent une opposition (ex. : « Dans quelle mesure X ? », « Comment concilier X et Y ? », « X est-il encore pertinent face à Y ? »).
    \item \textbf{Annonce du plan :} \textbf{Formule rituelle :} « Pour y répondre, nous montrerons d'abord que... (I). Nous analyserons ensuite... (II). Nous verrons enfin... (III). » \textbf{Interdiction :} « Nous ferons un plan en trois parties. » Les titres I, II, III doivent être des \textbf{affirmations} (réponses partielles), pas des noms de thèmes.
\end{enumerate}
\textbf{Reflexe Bac :} L'introduction se rédige \textbf{au brouillon}, se peaufine, puis se recopie au propre. Elle conditionne la note du développement (cohérence).
\end{methode}

\begin{exoresolu}[Rédaction d'une introduction type Bac : Sujet « La médaille, symbole d'honneur ou marque de servitude ? »]
\textbf{Énoncé :} Rédigez l'introduction complète pour ce sujet de dissertation (corpus : \emph{Le Vieux nègre et la médaille} + culture personnelle).

\tcblower
\textbf{Solution détaillée :}
\emph{« \textbf{[Amorce]} Lors de la remise de la médaille à Meka, le Commandant proclame : \emph{« C'est un honneur pour vous, Meka, et pour tout votre peuple »}. Pourtant, le roman de Ferdinand Oyono, \emph{Le Vieux nègre et la médaille} (1956), transforme cette cérémonie en une scène de dérisoire aliénation. \textbf{[Définition / Tension]} Le terme « médaille » désigne originellement une récompense honorifique, distinction glorieuse. Mais dans l'univers colonial dépeint par Oyono, il bascule vers son contraire : une marque physique de soumission, un « collier » invisible. \textbf{[Problématique]} \textbf{Dans quelle mesure la médaille, objet censé honorer le « bon indigène », devient-elle dans l'œuvre le symbole le plus cruel de sa déshumanisation et de sa servitude volontaire ?} \textbf{[Annonce du plan]} Pour y répondre, nous montrerons d'abord que \textbf{la médaille est le produit d'une manipulation symbolique coloniale} (I). Nous analyserons ensuite \textbf{comment l'ironie narrative dévoile l'intériorisation de l'infériorité par le colonisé} (II). Nous verrons enfin \textbf{que le refus final de Meka (le jet de la médaille) ouvre la voie à une reconquête de la dignité} (III). »}

\textbf{Justification des choix :}
\begin{itemize}
    \item \textbf{Amorce :} Citation textuelle précise (ancrage œuvre), contraste immédiat annoncé.
    \item \textbf{Définition :} Polysémie honorifique / aliénante. Champ lexical : récompense / collier.
    \item \textbf{Problématique :} Question unique, tension « honorer vs déshumaniser », « servitude volontaire » (thèse centrale).
    \item \textbf{Plan :} 3 affirmations dialectiques (Thèse : manipulation / Antithèse : intériorisation / Synthèse : révolte/réappropriation). Titres = réponses partielles.
\end{itemize}
\end{exoresolu}

\begin{methode}[Le Paragraphe Argumenté Standard (I.D. - A - E - M.C.)]
\textbf{Objectif :} Produire des paragraphes de développement rigoureux (notation : 2-3 pts / paragraphe).
\begin{enumerate}
    \item \textbf{Idée Directrice (Phrase thèse) :} 1 phrase affirmative, répondant à une étape du plan. \textbf{Connecteur d'ouverture obligatoire.} (Ex. : « \emph{En premier lieu}, la médaille incarne la manipulation symbolique... »).
    \item \textbf{Argumentation (Explication) :} 2 à 3 phrases. Développer la logique : \emph{Pourquoi ? Comment ?} Utiliser connecteurs \emph{car, en effet, puisque, dès lors que}. Mobiliser les concepts de cours (ironie, focalisation, aliénation, ethos...).
    \item \textbf{Exemple / Illustration (Preuve) :} 1 à 2 phrases. \textbf{Citation courte intégrée} (guillemets, chapitre) OU référence précise à un passage (scène, personnage, détail). Connecteur : \emph{Par exemple, ainsi, c'est le cas de, notamment}.
    \item \textbf{Mini-Conclusion / Transition (Synthèse partielle) :} 1 phrase. Résume l'idée, montre sa limite ou appelle l'argument suivant. Connecteur : \emph{En définitive, ainsi, partant de là, cela nous amène à...}.
\end{enumerate}
\textbf{Règle d'or :} \textbf{1 Paragraphe = 1 Idée Directrice.} Pas de liste de courses. Pas d'exemple sans explication. Pas d'explication sans exemple.
\end{methode}

\begin{exoresolu}[Rédaction du Paragraphe I : « La médaille, produit d'une manipulation symbolique »]
\textbf{Énoncé :} Rédigez le paragraphe I complet (niveau Bac) à partir de l'idée directrice annoncée dans l'introduction précédente.

\tcblower
\textbf{Solution détaillée :}
\emph{« \textbf{En premier lieu}, la médaille remise à Meka est le produit achevé d'une \textbf{manipulation symbolique coloniale} savamment orchestrée. \textbf{En effet}, l'administration française ne récompense pas un mérite objectif (bravoure, compétence), mais une \textbf{servilité performative} : Meka a « bien servi » en recrutant des porteurs, en payant ses impôts, en acceptant l'ordre établi. \textbf{Dès lors}, la cérémonie officielle, décrite avec une solennité grotesque (\emph{« le Commandant, très droit, très blanc »}), vise à \textbf{naturaliser la domination} : le colonisé reçoit une distinction occidentale, gage de sa « civilisation », à condition de rester à sa place. \textbf{Par exemple}, le discours du Commandant use d'une rhétorique paternaliste (\emph{« père », « enfant », « fidélité »}) qui infantilise Meka tout en le flattant, transformant l'obéissance en « honneur ». \textbf{Ainsi}, la médaille agit comme un \textbf{leurre institutionnel} : elle matérialise le contrat inégal « protection contre soumission », fondement de l'ordre colonial. \textbf{Partant de là}, il faut examiner comment cette manipulation s'intériorise dans la conscience même du récipiendaire. »}

\textbf{Check-list Bac :}
\begin{itemize}
    \item [✓] Connecteur ouverture (\emph{En premier lieu}).
    \item [✓] Idée Directrice claire (Manipulation symbolique).
    \item [✓] Argumentation (Mérite vs Servilité, Naturalisation, Infantilisation).
    \item [✓] Exemple précis (Discours Commandant, citation intégrée, analyse rhétorique).
    \item [✓] Mini-conclusion + Transition vers II (\emph{Partant de là... intériorisation}).
\end{itemize}
\end{exoresolu}

\begin{methode}[La Conclusion en 3 temps (Bilan - Ouverture - Dernière phrase choc)]
\textbf{Objectif :} Clore le devoir sans le répéter bêtement, montrer l'élévation de la réflexion.
\begin{enumerate}
    \item \textbf{Bilan synthétique (Réponse à la problématique) :} 2-3 phrases. Reprendre les 3 idées directrices (I, II, III) pour formuler la \textbf{réponse globale}. \emph{« Au terme de cette étude, il apparaît que... »} / \emph{« En définitive, la médaille... »}. \textbf{Ne pas recopier l'annonce du plan.}
    \item \textbf{Ouverture (Élargissement) :} 1-2 phrases. Dépasser le sujet précis sans faire du hors-sujet. Pistes : Autre œuvre (Oyono : \emph{Une vie de boy}, \emph{Le Prix du sang}), autre auteur (Césaire, Fanon, Beti), autre art (cinéma, peinture), actualité (restitution œuvres d'art, mémoire coloniale), question philosophique (liberté, identité). \emph{« Cette analyse invite à relire... »} / \emph{« Elle résonne avec l'actualité de... »}.
    \item \textbf{Phrase de chute (Dernière impression) :} 1 phrase forte, mémorable, souvent aphoristique ou interrogative. \emph{« La vraie médaille, c'est la liberté reconquise. »} / \emph{« Et si l'honneur véritable était de refuser toute médaille ? »}.
\end{enumerate}
\textbf{Interdictions :} « En conclusion, j'ai montré que... », « Pour conclure, je dirai que... », introduction d'un \textbf{nouvel argument}, « J'espère avoir convaincu. »
\end{methode}

\begin{exoresolu}[Rédaction de la conclusion pour le sujet « Médaille : honneur ou servitude ? »]
\textbf{Énoncé :} Rédigez la conclusion complète.

\tcblower
\textbf{Solution détaillée :}
\emph{« \textbf{[Bilan]} Au terme de cette analyse, la médaille apparaît bien moins comme une récompense que comme un \textbf{instrument de pouvoir colonial}. D'abord outil de manipulation symbolique (I), elle devient ensuite le miroir d'une aliénation intériorisée où le colonisé collabore à sa propre diminution (II), avant de se révéler, in fine, comme le catalyseur d'une prise de conscience lucide et d'un geste de révolte fondatrice (III). \textbf{[Ouverture]} Cette dynamique — du leurre à la lucidité — dépasse le cas Meka pour éclairer l'ensemble de la littérature négro-africaine d'expression française, de \emph{L'Aventure ambiguë} de Cheikh Hamidou Kane à \emph{Le Prix du sang} d'Oyono lui-même : l'objet colonial (médaille, école, uniforme) est toujours un piège identitaire dont il faut se déprendre. \textbf{[Chute]} \textbf{La seule médaille qui vaille, pour l'homme libre, est celle qu'il forge lui-même dans le refus de l'indignité.} »}

\textbf{Analyse :} Bilan synthétique (pas liste I/II/III mais articulation logique). Ouverture littéraire pertinente (Kane, Oyono autre œuvre). Chute forte, aphoristique, reprise mot « médaille » / « honneur » / « liberté ».
\end{exoresolu}

\courssub{Processus complet : De la problématique à la rédaction finale (Méthode « Brouillon structuré »)}

\begin{methode}[Élaboration du Plan Détaillé et Rédaction Gérée (4h type Bac)]
\textbf{Objectif :} Maîtriser le temps d'examen (4h) : Analyse (30') $\to$ Plan détaillé (45') $\to$ Rédaction (2h15') $\to$ Relecture (30').
\begin{enumerate}
    \item \textbf{Analyse du sujet (30').} \textbf{Surligner} : Mots-clés, verbes opératoires (Analyser, Commenter, Discuter, Dans quelle mesure), bornes (œuvre, période, corpus). \textbf{Questionner} : Qu'est-ce qui est acquis ? Qu'est-ce qui est discutable ? \textbf{Formuler} 3-4 problématiques possibles, choisir la plus féconde.
    \item \textbf{Plan Détaillé au Brouillon (45').} \textbf{Tableau 3 colonnes (Idée Directrice / Arguments / Exemples Précis).}
        \begin{center}
        \begin{tabularx}{\linewidth}{|>{\bfseries}X|X|X|}\hline
        \rowcolor{popblueL} \textbf{I. Idée Directrice (Titre Partie)} & \textbf{Arguments (Sous-parties A, B, C)} & \textbf{Exemples / Citations / Références} \\ \hline
        I. Manipulation symbolique & A. Mérite redéfini (servilité) & Scène cérémonie Ch.6 / Discours Commandant \\
         & B. Naturalisation domination & Description Commandant / Meka « enfant » \\
         & C. Leurre institutionnel & Médaille = contrat inégal \\ \hline
        II. Intériorisation / Aliénation & A. Fierté naïve de Meka & « Il se sentait grand » Ch.6 \\
         & B. Regard des autres (jalousie/respect) & Réaction village / Kelara (fils instruit) \\
         & C. Corps marqué (poids physique) & « Poids sur la poitrine » / Enterrement Kelara \\ \hline
        III. Lucidité / Révolte & A. Prise de conscience progressive & Mort Kelara / Injustice / Prison \\
         & B. Geste de rejet (jet médaille) & Final : « Il la jeta dans la boue » \\
         & C. Recouvrance dignité / Identité & Silence final / Regard neuf sur village \\ \hline
        \end{tabularx}
        \end{center}
        \textbf{Règle :} Chaque case « Arguments » = 1 paragraphe (donc 9 paragraphes au total). Vérifier l'équilibre (3 parties $\times$ 3 §).
    \item \textbf{Rédaction (2h15').} Suivre le plan \textbf{à la lettre}. Rédiger Intro $\to$ Développement (en sautant des lignes pour aérer) $\to$ Conclusion. \textbf{Ne pas réinventer le plan en rédigeant.}
    \item \textbf{Relecture Active (30').} 3 passages obligatoires :
        \begin{enumerate}
            \item \textbf{Sens / Cohérence :} Problématique répondue ? Plan respecté ? Transitions logiques ? Exemples pertinents ?
            \item \textbf{Langue / Style :} Accords (sujet/verbe, participe passé), conjugaison, ponctuation, \textbf{connecteurs variés} (pas 5 « De plus » de suite), vocabulaire précis (pas « chose », « truc », « bien »).
            \item \textbf{Propreté / Présentation :} Ratures propres, paragraphes indentés, titres I/II/III visibles (ou saut de ligne double), lisibilité.
        \end{enumerate}
\end{enumerate}
\textbf{Astuce temps :} Si blocage sur un paragraphe $\to$ passer au suivant, revenir à la fin. Mieux vaut 9 paragraphes inégaux que 6 parfaits et 3 manquants.
\end{methode}

\begin{exoresolu}[Simulation Bac : Sujet « L'ironie dans Le Vieux nègre et la médaille : une arme critique ? » — Plan détaillé + Intro]
\textbf{Énoncé :} Sujet de dissertation probable. Produisez : 1. La problématique retenue. 2. Le plan détaillé (tableau 3 colonnes comme ci-dessus, parties I, II, III). 3. L'introduction rédigée.

\tcblower
\textbf{Solution détaillée :}

\textbf{1. Problématique :} \textbf{En quoi l'ironie, par son double langage, constitue-t-elle l'arme critique principale d'Oyono pour déconstruire le mythe colonial et restaurer la dignité humaine ?}

\textbf{2. Plan Détaillé :}
\begin{center}
\begin{tabularx}{\linewidth}{|>{\bfseries}X|X|X|}\hline
\rowcolor{popblueL} \textbf{Partie / Idée Directrice} & \textbf{Arguments (Sous-parties)} & \textbf{Exemples / Citations Précis} \\ \hline
\textbf{I. L'ironie narrative : déconstruction du discours colonial officiel} & A. Ironie situationnelle : Cérémonie vs Réalité & Ch.6 : Solennité Commandant / Inconfort Meka / Chaleur / Mouches \\
 & B. Ironie verbale : Antiphrase du vocabulaire honoraire & « Honneur », « Fidélité », « Mérite » $\to$ Sens inversé par le contexte \\
 & C. Focalisation interne ironique : Meka dupe puis lucide & Pensées Meka : « grand » $\to$ « petit vieux noir » (Ch.10) \\ \hline
\textbf{II. L'ironie dramatique : le lecteur complice, le colonisé piégé} & A. Supériorité cognitive du lecteur (savoir vs ignorance Meka) & Lecteur sait médaille = marque bétail ; Meka croit honneur (début) \\
 & B. Personnages miroirs : Kelara (lucidité) vs Meka (naïveté) & Dialogues père/fils : Kelara refuse médaille, voit le piège \\
 & C. Ironie du sort : Médaille portée à l'enterrement du fils tué par l'ordre qu'elle honore & Scène enterrement Kelara : Meka décoré pleure victime du système \\ \hline
\textbf{III. L'ironie finale : retournement libérateur et portée universelle} & A. Geste final : Jet de la médaille = rejet ironique du symbole & Fin roman : « Il la jeta dans la boue » = acte performatif de liberté \\
 & B. Récupération de la parole / du silence & Meka ne parle plus « langue du colon », redevient ancien respecté \\
 & C. Portée méta-littéraire : Oyono écrivain ironiste = « Vieux nègre » qui rend sa médaille à la littérature coloniale & Titre ironique : « Vieux nègre » (réduit) + « Médaille » (leurre) $\to$ Roman = contre-médaille \\ \hline
\end{tabularx}
\end{center}

\textbf{3. Introduction rédigée :}
\emph{« \textbf{[Amorce]} « L'ironie est la tristesse qui a appris à sourire », écrit Fernando Pessoa. Dans \emph{Le Vieux nègre et la médaille}, Ferdinand Oyono fait de ce sourire amer une arme politique. \textbf{[Définition / Tension]} L'ironie, figure de rhétorique disant le contraire de ce qu'on pense, devient ici une \textbf{stratégie narrative globale} : le récit tout entier bascule entre le discours dominant (colonial, honorifique) et la réalité vécue (humiliante, absurde). \textbf{[Problématique]} \textbf{En quoi l'ironie, par son double langage, constitue-t-elle l'arme critique principale d'Oyono pour déconstruire le mythe colonial et restaurer la dignité humaine ?} \textbf{[Annonce du plan]} Pour y répondre, nous montrerons d'abord que \textbf{l'ironie narrative déconstruit méthodiquement le discours colonial officiel} (I). Nous analyserons ensuite \textbf{comment l'ironie dramatique place le lecteur en position de lucidité face à l'aveuglement du colonisé} (II). Nous verrons enfin \textbf{que l'ironie finale opère un retournement libérateur, fondateur d'une nouvelle subjectivité africaine} (III). »}
\end{exoresolu}

\courssub{Auto-évaluation, Remédiation et Préparation Finale}

\begin{methode}[Grille d'auto-évaluation « Prêt pour le Probatoire » (Check-list finale)]
\textbf{Objectif :} Vérifier son devoir avant de rendre la copie (ou s'auto-évaluer en entraînement).
\begin{center}
\begin{tabularx}{\linewidth}{|l|X|c|}\hline
\rowcolor{popblueL} \textbf{Critère (Barème indicatif)} & \textbf{Indicateurs de réussite (Oui / Non / Partiel)} & \textbf{Pts} \\ \hline
\textbf{INTRODUCTION (4-5 pts)} & Amorce pertinente (pas cliché) ? Termes définis ? Problématique unique, débattue, au cœur du sujet ? Plan annoncé en 3 affirmations (I, II, III) ? & /5 \\ \hline
\textbf{DÉVELOPPEMENT (10-12 pts)} & 3 Parties équilibrées (3 § chacune = 9 §) ? Chaque § : Idée Directrice (début) + Arguments + Exemple cité/précis + Mini-conclusion/Transition ? Problématique traitée (pas hors-sujet) ? Arguments variés (textuels, littéraires, contextuels) ? & /12 \\ \hline
\textbf{CONCLUSION (3-4 pts)} & Bilan synthétique (réponse probl.) ? Ouverture pertinente (littéraire/historique/actu) ? Phrase de chute forte ? Pas de nouvel argument ? & /4 \\ \hline
\textbf{LANGUE / EXPRESSION (4-5 pts)} & Orthographe / Grammaire / Conjugaison (0 faute = 5, 1-2 = 4, 3-5 = 3, $>$5 = 1-2) ? Connecteurs variés \& justes ? Vocabulaire précis (termes d'analyse) ? Phrases complexes maîtrisées ? & /5 \\ \hline
\textbf{PRÉSENTATION (1-2 pts)} & Lisible, aéré (sauts de lignes §), titres I/II/III visibles, ratures propres, encre bleue/noire ? & /2 \\ \hline
\textbf{TOTAL} &  & \textbf{/28} \\ \hline
\end{tabularx}
\end{center}
\textbf{Note sur 20 :} $\dfrac{\text{Total}}{28} \times 20$.
\textbf{Remédiation ciblée :} Si note $<$ 14/20 $\to$ Identifier LA faiblesse dominante (Problématique floue ? Exemples manquants ? Langue ? Structure ?). Refaire \textbf{uniquement} l'exercice correspondant (ex. : 5 intro en 30 min, 10 paragraphes types, 200 lignes de réécriture correction).
\end{methode}

\begin{exoresolu}[Correction type « Copie élève » : Diagnostic et Remédiation]
\textbf{Énoncé :} Voici l'introduction et le début du développement d'un élève (Sujet : « La médaille, honneur ou servitude ? »). \textbf{1.} Notez sur 5 l'intro (grille ci-dessus). \textbf{2.} Identifiez 3 erreurs majeures dans le §1 du développement. \textbf{3.} Proposez une réécriture corrigée de ce §1.

\textit{Copie élève :}
\textit{Intro : « Depuis toujours, les hommes aiment les médailles. Dans le livre de Oyono, Meka a une médaille. C'est un honneur mais c'est aussi une servitude. Est-ce que la médaille est un honneur ou une servitude ? Je vais dire que c'est les deux. D'abord je vais parler de l'honneur. Ensuite de la servitude. Enfin de la fin. »}
\textit{Développement I-A : « D'abord la médaille c'est un honneur. Meka il est content. Il la met sur sa poitrine. Le commandant il dit bravo. C'est bien pour lui. Il est fier. Par exemple page 45 il dit je suis content. Donc c'est un honneur. »}

\tcblower
\textbf{Solution détaillée :}
\textbf{1. Note Introduction : 1/5 (Insuffisant).}
\begin{itemize}
    \item \textbf{Amorce :} « Depuis toujours... » (\textbf{Interdit : cliché banal}).
    \item \textbf{Définition :} Absente. Pas d'analyse des termes « honneur », « servitude », « médaille ».
    \item \textbf{Problématique :} Reformulation plate du sujet (« Est-ce que... ? »), fermée (Oui/Non), pas débattue.
    \item \textbf{Plan :} « Je vais dire que c'est les deux » (pas d'affirmation thèse). « D'abord honneur, ensuite servitude, enfin la fin » (Thèmes, pas idées directrices ; « la fin » n'est pas un argument).
\end{itemize}

\textbf{2. Trois erreurs majeures au §1 (I-A) :}
\begin{enumerate}
    \item \textbf{Langue / Registre :} Familier oralisé (« Meka il est content », « Le commandant il dit », « C'est bien pour lui »). \textbf{Exigence Bac :} Français standard soutenu.
    \item \textbf{Structure argumentative absente :} Pas d'Idée Directrice claire (juste « c'est un honneur »). Pas d'argument explicatif (POURQUOI c'est un honneur dans le roman ?). Exemple plaqué sans analyse (« page 45 il dit... »).
    \item \textbf{Contenu / Compréhension :} Contredit la thèse du roman (l'ironie montre que ce n'est \emph{pas} un honneur vrai). L'élève prend le point de vue de Meka (naïf) pour thèse, au lieu du point de vue de l'auteur (critique).
\end{enumerate}

\textbf{3. Réécriture corrigée (Niveau Bac visé) :}
\emph{« \textbf{En premier lieu}, la médaille fonctionne comme une \textbf{simulation d'honneur} soigneusement mise en scène par le pouvoir colonial. \textbf{En effet}, la cérémonie du chapitre 6 déploie tous les attributs de la reconnaissance officielle : le Commandant « très droit, très blanc », le discours solennel qualifiant Meka de « fidèle serviteur », la remise solennelle de l'insigne. \textbf{Cependant}, cette honorabilité est \textbf{conditionnelle et intéressée} : elle récompense non l'excellence mais la \textbf{docilité} (recrutement porteurs, perception impôts). \textbf{Par exemple}, l'ironie narrative souligne le décalage entre la fierté naïve de Meka (\emph{« il se sentait grand, il était quelqu'un »}) et la réalité du contrat implicite : la médaille « pèse sur sa poitrine » comme le poids de sa soumission. \textbf{Ainsi}, l'honneur apparent n'est que le vernis doré d'une servitude structurelle, premier pas vers l'aliénation. »}
\end{exoresolu}

\begin{attention}[Les 5 erreurs qui coûtent des points au Bac (Français Terminale Tech)]
\begin{enumerate}
    \item \textbf{Problématique absente, fermée ou hors-sujet.} $\to$ Note max intro : 1/5. Le développement « flotte ». \textbf{Remède :} S'entraîner à formuler 3 probl. par sujet au brouillon, choisir celle qui génère le meilleur plan dialectique.
    \item \textbf{Plan « Thématique » au lieu de « Dialectique / Argumentatif ».} (Ex. : I. L'honneur / II. La servitude / III. La médaille). $\to$ Pas de démonstration, juste une liste. \textbf{Remède :} Titres de parties = \textbf{Affirmations argumentées} (réponses partielles à la probl.), pas noms de thèmes.
    \item \textbf{Paragraphes sans Idée Directrice ou sans Exemple textuel précis.} $\to$ « Blabla » général. \textbf{Remède :} Discipline du \textbf{1 § = 1 ID + 1 Cit/Ex précis + 1 Analyse}. Citer le chapitre.
    \item \textbf{Langue : Fautes d'accords de base (sujet/verbe, part. passé), connecteurs répétitifs (« Et », « Aussi », « De plus » x5), vocabulaire vague.} $\to$ Perte 3-5 pts sur « Expression ». \textbf{Remède :} Relecture \textbf{unique} focalisée « Accords » (lire à l'envers), liste perso connecteurs variés à côté de la copie.
    \item \textbf{Gestion du temps : Introduction parfaite mais développement incomplet (2 parties sur 3) ou conclusion absente.} $\to$ Pénalité lourde sur « Développement » et « Conclusion ». \textbf{Remède :} \textbf{Chronomètre impératif :} 30' Ana / 45' Plan / 2h15' Rédaction / 30' Relecture. Si blocage $\to$ Sauter, finir plan, revenir. Mieux vaut 9 § imparfaits que 6 parfaits.
\end{enumerate}
\end{attention}

\newpage

\coursec{Exercices}

\courssub{Je vérifie mes connaissances}

\begin{exercice}[QCM : « Le Vieux nègre et la médaille »]
Pour chacune des affirmations suivantes, cochez la bonne réponse.
\begin{enumerate}
\item L'auteur de \emph{Le Vieux nègre et la médaille} est : 
\begin{itemize}
\item[a)] Mongo Beti
\item[b)] Ferdinand Oyono
\item[c)] Francis Bebey
\item[d)] René Philombe
\end{itemize}
\item Le protagoniste principal, Meka, reçoit une médaille pour : 
\begin{itemize}
\item[a)] Son courage au combat
\item[b)] Sa fidélité à l'administration coloniale
\item[c)] Son rôle de chef traditionnel
\item[d)] Son aide aux missionnaires
\end{itemize}
\item La scène finale où Meka jette la médaille dans les latrines symbolise : 
\begin{itemize}
\item[a)] La folie du personnage
\item[b)] Le rejet de l'aliénation et la reconquête de la dignité
\item[c)] L'ingratitude envers la France
\item[d)] Un geste rituel traditionnel
\end{itemize}
\item Le roman met en scène principalement le choc entre : 
\begin{itemize}
\item[a)] Deux ethnies camerounaises
\item[b)] La tradition et la modernité imposée par le colonisateur
\item[c)] Le christianisme et l'islam
\item[d)] Les jeunes et les anciens du village
\end{itemize}
\end{enumerate}
\end{exercice}

\begin{exercice}[Vrai ou Faux justifié : Outils de liaison]
Indiquez si les affirmations sont vraies (V) ou fausses (F) et justifiez brièvement votre réponse.
\begin{enumerate}
\item La conjonction de coordination « mais » exprime toujours une opposition.
\item L'adverbe de liaison « par conséquent » marque une conséquence et se place toujours en début de phrase.
\item Le pronom relatif « dont » ne peut remplacer que des compléments d'objet indirect introduits par la préposition « de ».
\item Les locutions conjonctives comme « bien que », « pour que », « afin que » exigent systématiquement le subjonctif.
\end{enumerate}
\end{exercice}

\begin{exercice}[Définitions et concepts clés]
Donnez une définition précise pour chacun des termes suivants, en lien avec la dissertation :
\begin{enumerate}
\item \cle{Problématique}
\item \cle{Plan dialectique} (Thèse / Antithèse / Synthèse)
\item \cle{Idée directrice} (ou phrase-topic)
\item \cle{Connecteur logique} (donnez deux exemples de natures grammaticales différentes)
\end{enumerate}
\end{exercice}

\begin{exercice}[Identification de connecteurs]
Dans le paragraphe ci-dessous, soulignez tous les outils de liaison (conjonctions, adverbes, pronoms relatifs, locutions) et précisez pour chacun sa nature grammaticale et sa fonction logique (cause, conséquence, opposition, concession, but, etc.).

\textit{« Meka avait servi fidèlement l'administration coloniale ; \textbf{c'est pourquoi} il espérait une récompense. \textbf{Or}, la médaille qu'il reçut ne lui apporta que mépris de la part des siens. \textbf{Bien qu'}il fût vieux, il comprit \textbf{que} l'honneur ne s'achète pas. Il jeta donc le symbole de sa honte \textbf{afin de} retrouver sa liberté. »}
\end{exercice}

\courssub{Je m'entraîne}

\begin{exercice}[Analyse de stratégies argumentatives]
Relisez l'extrait suivant de \emph{Le Vieux nègre et la médaille} (chapitre final, monologue intérieur de Meka). 
\textit{« Il avait donné ses fils, ses terres, sa paix. On lui donnait un morceau de ferraille. [...] Il avait été le chien fidèle. On lui jetait un os. »}

\begin{enumerate}
\item Identifiez les figures de style présentes (au moins deux) et expliquez leur effet argumentatif.
\item Quel type de raisonnement (analogie, cause-effet, réductio ad absurdum, etc.) structure ce passage ?
\item Quel est le \cle{thème} et la \cle{thèse} défendue implicitement par le personnage/narrateur ?
\item Proposez un connecteur de concession pour relier les deux premières phrases en une seule phrase complexe.
\end{enumerate}
\end{exercice}

\begin{exercice}[Rédaction d'une introduction de dissertation]
Sujet type Probatoire : \textbf{« À travers l'étude du Vieux nègre et la médaille, montrez comment l'écrivain dénonce l'aliénation coloniale. »}

Rédigez une introduction complète (amorce, présentation du sujet/œuvre, problématique, annonce du plan) en respectant la méthode en entonnoir. Le plan attendu est thématique : 
I. Le mirage de la gloire coloniale ; 
II. La violence symbolique de l'ingratitude ; 
III. L'éveil de la conscience révoltée.
\end{exercice}

\begin{exercice}[Transformation linguistique : du style télégraphique au style académique]
Le paragraphe suivant est écrit dans un style « télégraphique » (phrases courtes, juxtaposées, manque de fluidité). 
\textit{« Meka est vieux. Meka est fier. Meka reçoit une médaille. La médaille est en fer. Les villageois rient. Meka est humilié. Meka jette la médaille. Meka retrouve sa dignité. »}

Réécrivez ce paragraphe en un seul paragraphe cohérent et fluide (style académique), en utilisant \textbf{au moins trois types de connecteurs différents} (ex: subordination conjonctive, adverbe de liaison, pronom relatif, participe présent/gérondif). Mettez en \textbf{gras} les connecteurs utilisés.
\end{exercice}

\begin{exercice}[Élaboration d'un plan détaillé]
Sujet : \textbf{« Discutez l'affirmation suivante : "Une médaille ne fait pas le héros." »}

\begin{enumerate}
\item Formulez une problématique pertinente pour ce sujet.
\item Choisissez le type de plan le plus adapté (dialectique, thématique, analytique) et justifiez votre choix en deux lignes.
\item Rédigez le plan détaillé (grandes parties I, II, III + 2 ou 3 arguments par partie avec exemples précis tirés de l'œuvre, de lectures personnelles ou de l'actualité camerounaise).
\end{enumerate}
\end{exercice}

\courssub{Je me prépare au Bac}

\begin{exercice}[Sujet type Probatoire : Dissertation complète (12 pts)]
\textbf{Sujet :} \textit{« À travers l'étude du Vieux nègre et la médaille, montrez comment l'écrivain dénonce l'aliénation coloniale. »}

\begin{enumerate}
\item \textbf{Problématique} (2 pts) : Formulez la problématique unique qui guidera votre devoir.
\item \textbf{Plan détaillé} (4 pts) : Présentez le plan thématique en trois parties (I, II, III) avec, pour chaque partie, deux arguments précis appuyés sur des références textuelles (citations ou scènes précises).
\item \textbf{Rédaction} (6 pts) : 
\begin{itemize}
\item[a)] Rédigez l'\textbf{introduction} complète.
\item[b)] Rédigez le \textbf{premier paragraphe de développement} (Partie I, premier argument) en respectant la structure : Idée directrice -- Argument -- Exemple/Analyse -- Mini-transition.
\end{itemize}
\end{enumerate}
\end{exercice}

\begin{exercice}[Sujet type Baccalauréat Technique : Dissertation complète (12 pts)]
\textbf{Sujet :} \textit{« Discutez l'affirmation suivante : "Une médaille ne fait pas le héros." Vous vous appuier sur Le Vieux nègre et la médaille, vos lectures personnelles et l'actualité camerounaise. »}

\begin{enumerate}
\item \textbf{Analyse du sujet} (2 pts) : Définissez les termes clés (« médaille », « héros ») et reformulez le sujet en question.
\item \textbf{Problématique et Plan dialectique} (4 pts) : Proposez une problématique dialectique et le plan correspondant (Thèse / Antithèse / Synthèse) avec arguments et exemples variés (œuvre au programme, lecture personnelle, actualité locale).
\item \textbf{Rédaction de la conclusion} (3 pts) : Rédigez une conclusion complète (bilan, ouverture) pour ce devoir.
\item \textbf{Qualité de la langue} (3 pts) : La correction linguistique (syntaxe, vocabulaire, connecteurs, orthographe) sera évaluée sur l'ensemble de la copie.
\end{enumerate}
\end{exercice}

\begin{exercice}[Exercices combinés : Méthodologie \& Langue (8 pts)]
\textbf{Partie A : Méthodologie de la problématique (3 pts)}
Pour le sujet : \textbf{« L'école est-elle le seul lieu de la réussite sociale ? »}, proposez \textbf{trois problématiques différentes} (une pour un plan dialectique, une pour un plan thématique, une pour un plan analytique). Justifiez en deux lignes pourquoi la problématique dialectique est la plus pertinente pour « discuter » ce sujet.

\textbf{Partie B : Expression écrite et maîtrise de la langue (5 pts)}
Transformez le paragraphe suivant (style oral / télégraphique) en un paragraphe argumentatif structuré, fluide et correct, adapté à une dissertation. Vous devez \textbf{impérativement} utiliser :
\begin{itemize}
\item Au moins une subordination conjonctive (cause, conséquence, but, concession, condition).
\item Au moins un adverbe de liaison ou locution adverbiale.
\item Au moins un pronom relatif (simple ou composé).
\item Un participe présent ou gérondif.
\end{itemize}
Mettez en \textbf{gras} les connecteurs/outils de liaison introduits.

\textit{Paragraphe à transformer :}
« Les jeunes cherchent du travail. Ils n'en trouvent pas. Ils ont des diplômes. Les diplômes ne servent à rien. Ils partent à l'étranger. Ils veulent gagner de l'argent. Ils laissent leur famille. C'est triste. Le pays perd ses forces vives. Il faut créer des entreprises ici. Les jeunes resteront. Le pays se développera. »
\end{exercice}

\newpage

\coursec{Corrigés des exercices}

\coursec{Corrigés des exercices — Leçon 02 : Dissertation et « Le Vieux nègre et la médaille »}

\begin{corrige}[Exercice 1 — QCM : « Le Vieux nègre et la médaille »]
\begin{enumerate}
\item \textbf{Réponse b) : Ferdinand Oyono.} \emph{Le Vieux nègre et la médaille} (1956) est le premier roman de Ferdinand Oyono, écrivain camerounais, auteur également d'\emph{Une vie de boy} et du \emph{Chemin d'Europe}. Mongo Beti a écrit \emph{Le Pauvre Christ de Bomba} ; Francis Bebey et René Philombe sont d'autres auteurs camerounais.
\item \textbf{Réponse b) : Sa fidélité à l'administration coloniale.} Meka est décoré de la médaille de l'amitié franco-africaine pour sa « fidélité » : il a donné ses terres aux missionnaires et ses deux fils, morts pour la France pendant la guerre.
\item \textbf{Réponse b) : Le rejet de l'aliénation et la reconquête de la dignité.} En jetant la médaille, Meka rejette le symbole de sa soumission au colonisateur et affirme sa liberté retrouvée : c'est un geste de décolonisation intérieure.
\item \textbf{Réponse b) : La tradition et la modernité imposée par le colonisateur.} Le roman oppose la communauté villageoise (ses valeurs, sa solidarité, sa sagesse) à l'ordre colonial (administration, mission, récompenses factices) qui bouleverse et humilie les traditions.
\end{enumerate}
\end{corrige}

\begin{attention}[Points de vigilance]
Ne pas confondre Oyono avec Mongo Beti : les deux dénoncent la colonisation, mais avec des œuvres et des styles différents. Retenir pour chaque auteur au programme au moins une œuvre et sa date.
\end{attention}

\begin{corrige}[Exercice 2 — Vrai ou Faux justifié : Outils de liaison]
\begin{enumerate}
\item \textbf{Vrai (avec nuance).} « Mais » est la conjonction de coordination qui exprime l'\cle{opposition} (ou la restriction). Exemple : « Meka était fier, \emph{mais} il fut humilié. » On accepte « vrai » car dans l'usage scolaire, « mais » introduit toujours une idée qui s'oppose ou limite la précédente.
\item \textbf{Faux.} « Par conséquent » marque bien la \cle{conséquence}, mais il ne se place pas \emph{toujours} en début de phrase : il peut être précédé de « et » (« il comprit, \emph{et par conséquent} il jeta la médaille ») ou être déplacé à l'intérieur de la phrase.
\item \textbf{Faux.} « Dont » remplace un complément introduit par « de » : complément du verbe (COI : « le livre \emph{dont} je parle »), mais aussi complément du nom (« le vieillard \emph{dont} le fils est mort »), complément d'agent ou complément d'adjectif (« fier de » : « la médaille \emph{dont} il était fier »).
\item \textbf{Faux.} « Bien que », « pour que », « afin que » exigent le subjonctif, mais ce n'est pas le cas de \emph{toutes} les locutions conjonctives : « parce que », « puisque », « dès que », « après que » se construisent avec l'indicatif. Il faut donc apprendre le mode exigé par chaque locution.
\end{enumerate}
\end{corrige}

\begin{attention}[Point de vigilance]
« Après que » prend l'indicatif (norme académique), bien que le subjonctif soit fréquent dans l'usage. « Avant que » exige, lui, le subjonctif.
\end{attention}

\begin{corrige}[Exercice 3 — Définitions et concepts clés]
\begin{enumerate}
\item \textbf{Problématique :} question centrale, née de l'analyse du sujet, qui pose le problème intellectuel à résoudre et guide l'ensemble du devoir. Elle se formule généralement par une phrase interrogative à l'introduction et commande le choix du plan.
\item \textbf{Plan dialectique :} plan en trois mouvements qui confronte deux points de vue opposés avant de les dépasser : la \cle{thèse} (premier point de vue, souvent l'opinion commune), l'\cle{antithèse} (point de vue contraire, qui nuance ou réfute), la \cle{synthèse} (dépassement, position équilibrée de l'élève). Il convient aux sujets introduits par « Discutez ».
\item \textbf{Idée directrice (phrase-topic) :} phrase placée en tête de paragraphe qui annonce l'idée principale du paragraphe. Elle est ensuite justifiée par un argument, illustrée par un exemple, puis conclue par une mini-transition.
\item \textbf{Connecteur logique :} mot ou groupe de mots qui relie les idées et exprime un rapport logique (cause, conséquence, opposition, concession, but, addition...). Exemples de natures différentes : une \emph{conjonction de subordination} (« parce que » — cause) ; un \emph{adverbe de liaison} (« cependant » — opposition) ; on peut citer aussi la conjonction de coordination (« donc ») ou le pronom relatif (« qui », outil de liaison phrastique).
\end{enumerate}
\end{corrige}

\begin{corrige}[Exercice 4 — Identification de connecteurs]
\begin{center}
\begin{tabularx}{\linewidth}{|l|l|X|}\hline
\rowcolor{popblueL} \textbf{Connecteur} & \textbf{Nature grammaticale} & \textbf{Fonction logique} \\ \hline
c'est pourquoi & locution adverbiale (adverbe de liaison) & conséquence : l'espérance de Meka résulte de sa fidélité \\ \hline
Or & adverbe de liaison & opposition / rupture : la réalité contredit l'attente \\ \hline
qu' (que) & conjonction de subordination & introduit une subordonnée complétive (COD de « comprit ») \\ \hline
Bien que & locution conjonctive (+ subjonctif « fût ») & concession : malgré son âge, il comprend \\ \hline
que (l'honneur...) & conjonction de subordination & subordonnée complétive \\ \hline
donc & conjonction de coordination & conséquence / conclusion logique \\ \hline
afin de & locution prépositive (+ infinitif) & but : retrouver sa liberté \\ \hline
\end{tabularx}
\end{center}
\end{corrige}

\begin{attention}[Point de vigilance]
Distinguer « donc » (conjonction de coordination) de « par conséquent » (adverbe de liaison) : tous deux expriment la conséquence mais n'ont pas la même nature. Ne pas confondre « que » conjonction (invariable, introduit une subordonnée) et « que » pronom relatif (représente un nom, a un antécédent).
\end{attention}

\begin{corrige}[Exercice 5 — Analyse de stratégies argumentatives]
\begin{enumerate}
\item \textbf{Figures de style :}
\begin{itemize}
\item \emph{La gradation descendante} : « ses fils, ses terres, sa paix » énumère des sacrifices de plus en plus intimes, ce qui amplifie l'ampleur du don de soi et rend l'injustice plus choquante.
\item \emph{L'antithèse} : l'immensité des sacrifices (« ses fils, ses terres, sa paix ») s'oppose à la misère de la récompense (« un morceau de ferraille », « un os »). Ce contraste dénonce l'ingratitude coloniale.
\item \emph{La métaphore animalière} : « le chien fidèle » à qui « on jetait un os » déshumanise Meka et dénonce le mépris du colonisateur, qui traite l'Africain comme une bête.
\item \emph{La litote} : « un morceau de ferraille » diminue la médaille pour souligner son néant symbolique.
\end{itemize}
Effet argumentatif : ces figures frappent l'imagination du lecteur, suscitent son indignation et rendent la dénonciation plus persuasive qu'un simple exposé.
\item \textbf{Type de raisonnement :} raisonnement par \cle{analogie} (comparaison implicite entre Meka et le chien fidèle récompensé d'un os) doublé d'un raisonnement de \cle{cause à effet} (les sacrifices consentis → la récompense dérisoire), qui aboutit à une forme d'absurdité : le rapport don/récompense est tellement disproportionné qu'il révèle la supercherie coloniale.
\item \textbf{Thème et thèse :} le \cle{thème} est l'ingratitude et le mépris de la colonisation envers ceux qui l'ont servie. La \cle{thèse} implicite : la fidélité au colonisateur n'est jamais récompensée à sa juste valeur ; la médaille est un leurre qui masque l'exploitation et l'aliénation.
\item \textbf{Connecteur de concession :} « \emph{Bien qu'}il eût donné ses fils, ses terres et sa paix, on ne lui donna qu'un morceau de ferraille. » (On accepte aussi : « Même s'il avait donné... », « Alors qu'il avait donné... ».)
\end{enumerate}
\end{corrige}

\begin{corrige}[Exercice 6 — Rédaction d'une introduction de dissertation]
\textbf{Corrigé rédigé (modèle) :}

« La littérature africaine d'expression française est née, pour une large part, du besoin de dire les blessures de la colonisation. \emph{(amorce)} C'est dans cette perspective que s'inscrit \emph{Le Vieux nègre et la médaille}, roman publié en 1956 par l'écrivain camerounais Ferdinand Oyono. L'œuvre raconte l'histoire de Meka, vieillard décoré de la médaille de l'amitié franco-africaine pour sa fidélité à l'administration coloniale, avant de comprendre, au prix d'une cruelle humiliation, la vanité de cette récompense. \emph{(présentation de l'œuvre et du sujet)} On peut alors se demander : par quels procédés romanesques et par quelles étapes du parcours de son héros Oyono dénonce-t-il l'aliénation coloniale ? \emph{(problématique)} Nous verrons d'abord le mirage de la gloire coloniale qui aveugle le vieillard, puis la violence symbolique de l'ingratitude qui le frappe, enfin l'éveil de sa conscience révoltée, signe d'une dignité reconquise. \emph{(annonce du plan)} »

\end{corrige}

\begin{methode}[Rappel de la méthode en entonnoir]
\begin{enumerate}
\item Amorce : phrase générale liée au thème (littérature engagée, colonisation).
\item Présentation : auteur, titre, date, résumé bref en lien avec le sujet.
\item Problématique : une question unique qui reprend les termes du sujet.
\item Annonce du plan : « Nous verrons d'abord..., puis..., enfin... »
\end{enumerate}
\end{methode}

\begin{attention}[Points de vigilance]
Interdiction de dire « je » ; pas de définition de dictionnaire en amorce ; l'annonce du plan ne doit pas être détaillée (pas d'arguments) ; la problématique doit être une vraie question, pas une affirmation.
\end{attention}

\begin{corrige}[Exercice 7 — Transformation linguistique : du style télégraphique au style académique]
\textbf{Corrigé rédigé (modèle) :}

« Meka, \textbf{qui} était un vieillard fier de sa longue vie, reçut une médaille \textbf{qu'}il crut d'abord être un honneur. \textbf{Cependant}, cette médaille de fer, \textbf{loin de} lui attirer le respect, ne suscita que les rires des villageois. \textbf{En constatant} l'humiliation \textbf{que} lui valait ce prétendu hommage, il comprit son erreur~; \textbf{c'est pourquoi} il jeta la médaille, \textbf{retrouvant ainsi} sa dignité perdue. »

\textbf{Connecteurs utilisés :}
\begin{itemize}
\item Pronoms relatifs : \textbf{qui}, \textbf{qu'}, \textbf{que} (liaison phrastique, évitent la répétition de « Meka » et de « médaille »).
\item Adverbe de liaison : \textbf{Cependant} (opposition).
\item Locution prépositive : \textbf{loin de} (opposition/nuance).
\item Gérondif / participe présent : \textbf{En constatant}, \textbf{retrouvant ainsi} (moyen, simultanéité, conséquence).
\item Locution adverbiale : \textbf{c'est pourquoi} (conséquence).
\end{itemize}
\end{corrige}

\begin{attention}[Point de vigilance]
Un paragraphe de dissertation évite les phrases juxtaposées : chaque phrase doit être reliée à la précédente par un connecteur, un pronom ou une subordination. Vérifier l'accord du participe présent (invariable) et du gérondif.
\end{attention}

\begin{corrige}[Exercice 8 — Élaboration d'un plan détaillé]
\begin{enumerate}
\item \textbf{Problématique :} « Les décorations et les honneurs officiels suffisent-ils à faire d'un homme un héros, ou la véritable héroïcité ne réside-t-elle pas ailleurs ? »
\item \textbf{Choix du plan :} le plan \cle{dialectique} est le plus adapté, car le sujet invite à « discuter » une affirmation : il faut d'abord examiner en quoi la médaille semble consacrer le mérite (thèse), puis montrer ses limites (antithèse), avant de définir la vraie héroïcité (synthèse).
\item \textbf{Plan détaillé :}
\begin{itemize}
\item \textbf{I. Thèse — La médaille semble récompenser et consacrer le mérite.}
\begin{itemize}
\item Arg. 1 : la décoration officialise le courage et la fidélité. Ex. : Meka est d'abord fier de sa médaille, le village entier se déplace pour la cérémonie.
\item Arg. 2 : les honneurs donnent un statut social. Ex. : les médaillés militaires et sportifs célébrés au Cameroun (lions indomptables décorés après 1990).
\end{itemize}
\item \textbf{II. Antithèse — Mais la médaille peut être un leurre qui masque l'ingratitude ou l'imposture.}
\begin{itemize}
\item Arg. 1 : la récompense peut être dérisoire au regard du sacrifice. Ex. : Meka a donné ses fils et ses terres ; il reçoit « un morceau de ferraille ».
\item Arg. 2 : les honneurs sont parfois instrumentalisés pour acheter la soumission. Ex. : la médaille vise à faire de Meka un exemple de fidélité coloniale ; actualité : distinctions parfois distribuées par clientélisme.
\end{itemize}
\item \textbf{III. Synthèse — Le vrai héros se définit par ses actes et sa dignité, non par ses décorations.}
\begin{itemize}
\item Arg. 1 : l'héroïsme véritable est dans le geste de lucidité et de liberté. Ex. : Meka jetant la médaille devient plus « héros » en la rejetant qu'en la recevant.
\item Arg. 2 : la postérité reconnaît les héros sans médaille. Ex. : les martyrs de l'indépendance camerounaise (Um Nyobè) longtemps non reconnus ; lecture personnelle : \emph{Une vie de boy}, où la dignité du colonisé prime sur les honneurs coloniaux.
\end{itemize}
\end{itemize}
\end{enumerate}
\end{corrige}

\begin{corrige}[Exercice 9 — Sujet type Probatoire : Dissertation complète (12 pts)]
\textbf{1. Problématique (2 pts) :} « Par quels moyens romanesques et à travers quelles étapes du parcours de Meka, Ferdinand Oyono dénonce-t-il l'aliénation coloniale ? » \emph{(1 pt pour la question unique et claire, 1 pt pour la reprise des termes du sujet : « dénonce », « aliénation coloniale ».)}

\textbf{2. Plan détaillé (4 pts) :}
\begin{itemize}
\item \textbf{I. Le mirage de la gloire coloniale.} Arg. 1 : Meka croit sincèrement à la grandeur de la France et à la valeur de la médaille (sa fierté, ses préparatifs, la fête organisée au village). Arg. 2 : le discours colonial flatte la fidélité pour mieux aliéner (le discours du commandant, la cérémonie grandiloquente).
\item \textbf{II. La violence symbolique de l'ingratitude.} Arg. 1 : la disproportion entre les sacrifices (ses deux fils morts à la guerre, ses terres données à la mission) et la récompense dérisoire. Arg. 2 : l'humiliation publique de Meka (l'arrestation, la nuit au poste, le mépris des Blancs après la cérémonie).
\item \textbf{III. L'éveil de la conscience révoltée.} Arg. 1 : la prise de conscience de Meka, qui comprend la supercherie (« il avait été le chien fidèle »). Arg. 2 : le geste final de rejet — la médaille jetée dans les latrines — acte de libération et de dignité retrouvée.
\end{itemize}
\emph{(Barème : 1 pt par grande partie conforme, 1 pt pour la précision des références textuelles.)}

\textbf{3. Rédaction (6 pts) :}
\begin{itemize}
\item[a)] \textbf{Introduction (3 pts) :} reprendre le modèle de l'exercice 6 : amorce sur la littérature africaine engagée (0,5 pt), présentation de l'œuvre et de l'auteur (0,5 pt), problématique (1 pt), annonce du plan en trois mouvements (1 pt).
\item[b)] \textbf{Premier paragraphe (3 pts) — modèle :} « \emph{Idée directrice :} Oyono peint d'abord un vieillard aveuglé par le mirage de la gloire coloniale. \emph{Argument :} Meka a intériorisé les valeurs du colonisateur au point de croire que la médaille couronnera une vie de fidélité ; il se prépare avec une joie naïve à la cérémonie. \emph{Exemple/analyse :} le roman montre le vieillard revêtant ses plus beaux habits, conviant tout le village, rêvant à voix haute de l'honneur qui lui sera fait — signe que l'aliénation est d'abord mentale : le colonisé admire celui qui l'exploite. \emph{Mini-transition :} Mais cette gloire espérée va bientôt se révéler un piège, car la récompense coloniale n'est que l'envers d'un profond mépris. » \emph{(1 pt idée directrice, 1 pt argument + exemple précis, 1 pt transition et qualité de la langue.)}
\end{itemize}
\end{corrige}

\begin{attention}[Points de vigilance]
Citer l'œuvre avec précision (scènes, répliques) ; interdiction de résumer tout le roman ; chaque partie doit répondre à la problématique ; soigner les connecteurs entre les parties.
\end{attention}

\begin{corrige}[Exercice 10 — Sujet type Baccalauréat Technique : Dissertation complète (12 pts)]
\textbf{1. Analyse du sujet (2 pts) :} « \cle{Médaille} » : distinction honorifique, récompense matérielle et symbolique décernée par une autorité ; par extension, tout honneur officiel. « \cle{Héros} » : personnage qui accomplit des actes exceptionnels de courage, de sacrifice ou de grandeur morale. \emph{Reformulation :} les honneurs officiels créent-ils la valeur héroïque d'un homme, ou celle-ci existe-t-elle indépendamment d'eux ? \emph{(1 pt définitions, 1 pt reformulation en question.)}

\textbf{2. Problématique et plan dialectique (4 pts) :} Problématique : « Si la médaille consacre publiquement le mérite, ne risque-t-elle pas d'être un leurre, la vraie héroïcité ne tenant qu'aux actes et à la dignité ? »
\begin{itemize}
\item \textbf{Thèse :} la médaille reconnaît et encourage le mérite. Ex. : fierté initiale de Meka ; décorations des sportifs camerounais ; légion d'honneur des anciens combattants.
\item \textbf{Antithèse :} la médaille peut être dérisoire, injuste ou instrument de domination. Ex. : « un morceau de ferraille » pour des fils et des terres donnés (\emph{Le Vieux nègre et la médaille}) ; distinctions décernées par favoritisme dans l'actualité.
\item \textbf{Synthèse :} l'héroïsme véritable est intérieur : courage, lucidité, dignité. Ex. : Meka rejetant la médaille ; Um Nyobè et les héros méconnus de l'indépendance ; lecture personnelle : \emph{Une vie de boy} d'Oyono.
\end{itemize}
\emph{(1 pt problématique dialectique, 3 pts plan équilibré avec exemples variés : œuvre au programme, lecture personnelle, actualité camerounaise.)}

\textbf{3. Conclusion rédigée (3 pts) — modèle :} « En définitive, si la médaille a l'apparence de consacrer le mérite, l'exemple de Meka montre qu'elle peut n'être qu'un leurre, voire un instrument d'aliénation. Le véritable héros n'a pas besoin de ferraille officielle : il se reconnaît à ses actes, à son courage et à sa dignité, fussent-ils ignorés des puissants. \emph{(bilan)} Ainsi, à l'heure où notre société multiplie les honneurs et les titres, ne conviendrait-il pas de s'interroger : que vaut une récompense qui ne s'accompagne pas de justice et de respect ? \emph{(ouverture)} » \emph{(1,5 pt bilan fidèle au plan, 1 pt ouverture pertinente, 0,5 pt langue.)}

\textbf{4. Qualité de la langue (3 pts) :} syntaxe correcte et variée (1 pt), vocabulaire précis et soutenu (1 pt), connecteurs logiques et orthographe (1 pt).
\end{corrige}

\begin{attention}[Points de vigilance]
Dans un plan dialectique, la synthèse ne doit pas être une simple répétition de la thèse : elle doit dépasser l'opposition. La conclusion ne doit introduire aucun argument nouveau ; l'ouverture doit rester liée au sujet.
\end{attention}

\begin{corrige}[Exercice 11 — Exercices combinés : Méthodologie \& Langue (8 pts)]
\textbf{Partie A — Méthodologie de la problématique (3 pts)}
\begin{itemize}
\item \emph{Problématique pour un plan dialectique :} « L'école est-elle réellement la voie unique de la réussite sociale, ou d'autres chemins mènent-ils aussi à l'accomplissement de l'individu ? »
\item \emph{Problématique pour un plan thématique :} « Quels sont les différents rôles de l'école dans la réussite sociale, et quelles en sont les limites ? »
\item \emph{Problématique pour un plan analytique :} « Pourquoi l'école est-elle considérée comme le principal instrument de la réussite sociale, et comment expliquer qu'elle ne suffise pas toujours ? » (causes puis conséquences)
\end{itemize}
\textbf{Justification :} la problématique dialectique est la plus pertinente car la consigne « discuter » suppose de confronter deux opinions opposées (l'école comme seule voie / l'école comme une voie parmi d'autres) avant de trancher de manière nuancée. \emph{(Barème : 0,75 pt par problématique adaptée, 0,75 pt justification.)}

\textbf{Partie B — Paragraphe transformé (5 pts) — modèle :}

« \textbf{Bien que} les jeunes soient diplômés, ils ne trouvent pas d'emploi, \textbf{si bien que} leurs diplômes, \textbf{qui} devraient être un passeport pour la vie active, semblent ne servir à rien. \textbf{C'est pourquoi} beaucoup partent à l'étranger, \textbf{espérant} y gagner leur vie, \textbf{même si} cet exil les arrache à leur famille. \textbf{En partant}, ils privent le pays de ses forces vives, \textbf{ce qui} appauvrit la nation entière. \textbf{Par conséquent}, il faut créer des entreprises ici \textbf{afin que} les jeunes restent et \textbf{que} le pays se développe. »

\textbf{Outils utilisés :} subordinations conjonctives de concession (\textbf{bien que}), de but (\textbf{afin que} + subjonctif) ; adverbes et locutions de liaison (\textbf{si bien que}, \textbf{c'est pourquoi}, \textbf{par conséquent}, \textbf{même si}) ; pronoms relatifs (\textbf{qui}, \textbf{ce qui}) ; gérondif et participe présent (\textbf{en partant}, \textbf{espérant}).

\emph{Barème : 1 pt par type d'outil exigé correctement employé (subordination, adverbe de liaison, pronom relatif, participe/gérondif) ; 1 pt pour la cohérence argumentative et la correction globale.}
\end{corrige}

\begin{attention}[Points de vigilance]
« Afin que » et « pour que » exigent le subjonctif (« restent », « se développe ») ; le participe présent est invariable (« en partant ») ; vérifier que chaque connecteur exprime bien le rapport logique voulu et que le paragraphe garde une progression argumentative (constat → conséquence → solution).
\end{attention}

\newpage

\coursec{Fiche bilan}

\coursec{Leçon 02 : Dissertation — Produire une dissertation}

\begin{popfigure}
\begin{center}\tcbox[colback=popdark,colframe=popdark,coltext=white,arc=9pt,boxsep=4pt,left=6pt,right=6pt]{\bfseries\small La dissertation (Bac technique)}\end{center}
\vspace{-2pt}
\begin{tcbraster}[raster columns=3,raster equal height=rows,raster column skip=6pt,raster row skip=6pt,enhanced,blankest]
\begin{tcolorbox}[enhanced,colback=white,colframe=popblue,boxrule=0.9pt,arc=3pt,title={L'introduction},fonttitle=\bfseries\scriptsize,coltitle=white,colbacktitle=popblue,left=4pt,right=4pt,top=2pt,bottom=2pt,boxsep=1pt,toptitle=1.5pt,bottomtitle=1.5pt,halign=left]
\begin{itemize}[nosep,leftmargin=*,itemsep=1pt,font=\scriptsize]
\item Accroche, sujet, problématique, plan
\end{itemize}
\end{tcolorbox}
\begin{tcolorbox}[enhanced,colback=white,colframe=poporange,boxrule=0.9pt,arc=3pt,title={Le paragraphe},fonttitle=\bfseries\scriptsize,coltitle=white,colbacktitle=poporange,left=4pt,right=4pt,top=2pt,bottom=2pt,boxsep=1pt,toptitle=1.5pt,bottomtitle=1.5pt,halign=left]
\begin{itemize}[nosep,leftmargin=*,itemsep=1pt,font=\scriptsize]
\item Idée directrice
\item Argument + exemple
\end{itemize}
\end{tcolorbox}
\begin{tcolorbox}[enhanced,colback=white,colframe=popgreen,boxrule=0.9pt,arc=3pt,title={La conclusion},fonttitle=\bfseries\scriptsize,coltitle=white,colbacktitle=popgreen,left=4pt,right=4pt,top=2pt,bottom=2pt,boxsep=1pt,toptitle=1.5pt,bottomtitle=1.5pt,halign=left]
\begin{itemize}[nosep,leftmargin=*,itemsep=1pt,font=\scriptsize]
\item Bilan + ouverture
\end{itemize}
\end{tcolorbox}
\begin{tcolorbox}[enhanced,colback=white,colframe=poppurple,boxrule=0.9pt,arc=3pt,title={Les liaisons},fonttitle=\bfseries\scriptsize,coltitle=white,colbacktitle=poppurple,left=4pt,right=4pt,top=2pt,bottom=2pt,boxsep=1pt,toptitle=1.5pt,bottomtitle=1.5pt,halign=left]
\begin{itemize}[nosep,leftmargin=*,itemsep=1pt,font=\scriptsize]
\item Conjonctions, adverbes, relatifs
\end{itemize}
\end{tcolorbox}
\begin{tcolorbox}[enhanced,colback=white,colframe=popgold,boxrule=0.9pt,arc=3pt,title={L'argumentation},fonttitle=\bfseries\scriptsize,coltitle=white,colbacktitle=popgold,left=4pt,right=4pt,top=2pt,bottom=2pt,boxsep=1pt,toptitle=1.5pt,bottomtitle=1.5pt,halign=left]
\begin{itemize}[nosep,leftmargin=*,itemsep=1pt,font=\scriptsize]
\item Thèse, arguments, connecteurs
\end{itemize}
\end{tcolorbox}
\begin{tcolorbox}[enhanced,colback=white,colframe=poppink,boxrule=0.9pt,arc=3pt,title={Le plan détaillé},fonttitle=\bfseries\scriptsize,coltitle=white,colbacktitle=poppink,left=4pt,right=4pt,top=2pt,bottom=2pt,boxsep=1pt,toptitle=1.5pt,bottomtitle=1.5pt,halign=left]
\begin{itemize}[nosep,leftmargin=*,itemsep=1pt,font=\scriptsize]
\item Problématique $\rightarrow$ 2 ou 3 parties
\end{itemize}
\end{tcolorbox}
\end{tcbraster}

\legende{Carte mentale de la leçon : les six compétences clés pour produire une dissertation réussie au Baccalauréat technique.}
\end{popfigure}

\begin{aretenir}[L'essentiel de la leçon]
\textbf{Définitions clés}
\begin{itemize}
\item \cle{Dissertation} : exercice écrit qui consiste à discuter un sujet de manière organisée, en défendant une réponse argumentée à une problématique.
\item \cle{Problématique} : question centrale (souvent formulée par un ou deux interrogatifs) qui naît de l'analyse du sujet et qui guide tout le devoir.
\item \cle{Idée directrice} : phrase qui annonce l'idée principale d'un paragraphe.
\item \cle{Argument} : raison logique qui justifie l'idée ; \cle{exemple} : fait précis (œuvre, événement, situation) qui illustre l'argument.
\item \cle{Connecteurs logiques} : mots de liaison (conjonctions, adverbes, pronoms relatifs) qui assurent la cohérence du texte.
\end{itemize}

\textbf{Architecture à connaître par c\oe ur}
\begin{center}
\begin{tabularx}{\linewidth}{|l|X|}
\hline
\rowcolor{popblueL} \textbf{Partie} & \textbf{Contenu obligatoire} \\
\hline
Introduction & Phrase d'accroche ; présentation du sujet ; problématique ; annonce du plan \\
\hline
Développement & 2 ou 3 grandes parties ; chaque paragraphe = idée directrice + argument(s) + exemple(s) + phrase de transition \\
\hline
Conclusion & Bilan (réponse à la problématique) ; élargissement ou ouverture \\
\hline
\end{tabularx}
\end{center}

\textbf{Méthodes express}
\begin{itemize}
\item \textbf{Analyser le sujet} : souligner les mots-clés, définir les termes, dégager la tension ou le paradoxe, formuler la problématique.
\item \textbf{Mobiliser la culture} : puiser des exemples dans \emph{Le Vieux nègre et la médaille} (Ferdinand Oyono) et dans les autres œuvres étudiées.
\item \textbf{Rédiger au brouillon} le plan détaillé avant de passer à la rédaction complète.
\item \textbf{Relire et améliorer} : vérifier les liaisons, l'orthographe, la cohérence entre l'introduction et la conclusion.
\end{itemize}
\end{aretenir}

\begin{aretenir}[Checklist avant le Bac]
Avant de rendre votre copie de dissertation, vérifiez :
\begin{itemize}
\item[$\square$] J'ai bien compris le sujet et défini tous ses mots-clés.
\item[$\square$] Ma problématique est une vraie question, claire et liée au sujet.
\item[$\square$] Mon introduction contient les quatre étapes : accroche, sujet, problématique, plan.
\item[$\square$] Mon développement comporte au moins deux grandes parties équilibrées.
\item[$\square$] Chaque paragraphe suit le schéma : idée directrice, argument, exemple.
\item[$\square$] J'ai cité au moins un exemple tiré du \emph{Vieux nègre et la médaille} ou d'une autre œuvre au programme.
\item[$\square$] J'ai utilisé des connecteurs logiques variés (d'abord, ensuite, cependant, en effet, par conséquent...).
\item[$\square$] Mes phrases sont reliées par des conjonctions de coordination et de subordination correctement employées.
\item[$\square$] Ma conclusion répond à la problématique et propose une ouverture.
\item[$\square$] J'ai relu ma copie pour corriger l'orthographe, la grammaire et la ponctuation.
\end{itemize}
\end{aretenir}

\findecours

\end{document}
